% concatenated from 20250826_Arxiv.tex
\documentclass[11pt]{amsart}
\usepackage[latin1]{inputenc}
\usepackage{amsfonts}
\usepackage[toc,page,title,titletoc,header]{appendix}
\usepackage{graphicx,psfrag,epsfig,multirow,caption,subcaption}
\usepackage{amssymb,amsmath,amscd,amsthm,amssymb,verbatim,setspace}
\numberwithin{equation}{section}
\usepackage{mathrsfs,wrapfig}
\usepackage{indentfirst}
\usepackage{extarrows}
\usepackage{float}
%\setlength{\parskip}{10pt}

\usepackage[colorlinks=true, backref=page]{hyperref}


\headheight=6.15pt \textheight=8in \textwidth=6.5in
\oddsidemargin=0in \evensidemargin=0in \topmargin=0in


\setcounter{section}{0}
%\theoremstyle{definition}

\newtheorem{theorem}{Theorem}[section]
\newtheorem{definition}[theorem]{Definition}
\newtheorem{conjecture}[theorem]{Conjecture}
\newtheorem{proposition}[theorem]{Proposition}
\newtheorem{lemma}[theorem]{Lemma}
\newtheorem{remark}[theorem]{Remark}
\newtheorem{corollary}[theorem]{Corollary}
\newtheorem{example}[theorem]{Example}
\newtheorem{question}[theorem]{Question}
\newtheorem{fact}[theorem]{Fact}
\newtheorem*{acknowledgements}{Acknowledgements}

\newtheorem*{remark*}{Remark}

\numberwithin{equation}{section}


\newcommand{\dv}{\text{div}}
\newcommand{\ee}{{\text{e}}}
\newcommand{\hps}{{\text{hypersurface}}}
\newcommand{\sk}{{\text{shrinker}}}
\newcommand{\sks}{{\text{shrinkers}}}
\newcommand{\mf}{\mathbf}
\newcommand{\mc}{\mathcal}
\newcommand{\mb}{\mathbb}
\newcommand{\mfk}{\mathfrak}
\newcommand{\oli}{\overline}
\newcommand{\eps}{\varepsilon}
\newcommand{\wti}{\widetilde}

\newcommand{\dt}{\delta}
\newcommand{\al}{\alpha}

\newcommand{\bn}{\mathbf{n}}
\newcommand{\br}{\mathbf{r}}
\newcommand{\bN}{\mathbf{N}}
\newcommand{\bM}{\mathbf{M}}
\newcommand{\bS}{\mathbb{S}}
\newcommand{\bT}{\mathbf{T}}


\newcommand{\cN}{\mathcal{N}}
\newcommand{\cA}{\mathcal{A}}
\newcommand{\cP}{\mathcal{P}}
\newcommand{\cM}{\mathcal{M}}
\newcommand{\cH}{\mathcal{H}}
\newcommand{\cE}{\mathcal{E}}
\newcommand{\cB}{\mathcal{B}}
\newcommand{\cS}{\mathcal{S}}
\newcommand{\cC}{\mathcal{C}}
\newcommand{\cL}{\mathcal{L}}
\newcommand{\cK}{\mathcal{K}}
\newcommand{\cT}{\mathcal{T}}
\newcommand{\cR}{\mathcal{R}}
\newcommand{\cQ}{\mathcal{Q}}
\newcommand{\cJ}{\mathcal{J}}
\newcommand{\cG}{\mathcal{G}}
\newcommand{\cU}{\mathcal{U}}
\newcommand{\cV}{\mathcal{V}}
\newcommand{\cZ}{\mathcal{Z}}
\newcommand{\cF}{\mathcal{F}}
\newcommand{\cI}{\mathcal{I}}

\newcommand{\bH}{\mathbf{H}}
\newcommand{\sR}{\mathscr{R}}
\newcommand{\sL}{\mathscr{L}}
\newcommand{\sT}{\mathscr{T}}
\newcommand{\mtr}{\mathfrak{r}}

\newcommand{\Si}{\Sigma}
\newcommand{\la}{\lambda}
\newcommand{\R}{\mathbb{R}}
\newcommand{\ind}{\mathrm{ind}}
\newcommand{\A}{\mathbb{A}}
\newcommand{\pr}{\partial}


\DeclareMathOperator{\vol}{Vol}
\DeclareMathOperator{\area}{Area}
\DeclareMathOperator{\Index}{index}
\DeclareMathOperator{\genus}{genus}
\DeclareMathOperator{\Ric}{Ric}
\DeclareMathOperator{\Hess}{Hess}
\DeclareMathOperator{\Clos}{Clos}
\DeclareMathOperator{\spt}{spt}
\DeclareMathOperator{\Graph}{Graph}
\DeclareMathOperator{\rad}{radius}
\DeclareMathOperator{\ang}{angle}
\DeclareMathOperator{\dist}{dist}
\DeclareMathOperator{\diam}{diam}
\DeclareMathOperator{\loc}{loc}
\DeclareMathOperator{\tr}{tr}
\DeclareMathOperator{\II}{II}
\DeclareMathOperator{\id}{\text{Id}}




\title[Cylindrical singularities I]{Generic mean curvature flows with cylindrical singularities I: the normal forms and nondegeneracy}
\author{Ao Sun}
\address{Lehigh University, Department of Mathematics, Chandler-Ullmann Hall, Bethlehem, PA 18015}
\email{aos223@lehigh.edu}
\author{Jinxin Xue}
\address{New Cornerstone Science Laboratory, Department of Mathematics, Rm A115, Tsinghua University, Haidian District, Beijing, 100084}
\email{jxue@tsinghua.edu.cn}
\date{\today}



\begin{document}
	\maketitle
	\begin{abstract}
		This paper studies the dynamics of mean curvature flow as it approaches a cylindrical singularity. We proved that the rescaled mean curvature flow converging to a smooth generalized cylinder can be written as a graph over the cylinder in a ball of radius $K\sqrt{t}$, and a normal form of the asymptotics. Using the normal form, we can define the nondegeneracy of cylindrical singularities, and we show that nondegenerate cylindrical singularities are isolated in space, have a mean convex neighborhood, and are type-I. 
	\end{abstract}
	
	
	
	\section{Introduction}
	This paper focuses on mean curvature flows with cylindrical singularities. Mean curvature flow is a fundamental geometric flow that has attracted considerable attention in diverse fields, including geometry, partial differential equations, and applied mathematics.
	
	Mean curvature flow is defined as a family of hypersurfaces $\{\bM_\tau\}_{\tau\in I}$ evolving in $\mathbb{R}^{n+1}$ according to the equation $\pr_\tau x=\vec{H}(x)$, where $\vec H$ is the mean curvature vector. Singularities must appear in mean curvature flows that originate from smooth embedded closed hypersurfaces. Therefore, understanding singularities is crucial to understanding mean curvature flows. The analysis of singularities involves a blow-up procedure. Huisken \cite{Hu1} introduced the concept of rescaled mean curvature flow to study these singularities. The rescaled mean curvature flow is a family of hypersurfaces $\{M_t\}_{t\in[0,\infty)}$ satisfying the equation the equation
	$
	\pr_t x=\vec{H}+\frac{x^\perp}{2}.
	$
	When the first singularity of the mean curvature flow $\{\bM_\tau\}_{\tau\in[-1,0)}$ occurs at the spacetime point $(0,0)\in \R^{n+1}\times \R$, the corresponding rescaled mean curvature flow $\{M_t\}_{t\in[0,\infty)}$ that captures this singularity is defined as follows:
	\begin{equation}\label{EqMCFRMCF}
		M_t=e^{t/2}\bM_{-e^{-t}}, \quad t\in [0,\infty).
	\end{equation}
	The singularity models, known as shrinkers, are obtained as a (subsequential) limit of the rescaled mean curvature flow, and they satisfy the equation $\vec H+\frac{x^\perp}{2}=0$.
	
	The rescaled mean curvature flow can be viewed as a dynamical system, a perspective introduced by Colding-Minicozzi in a series of papers \cite{CM1,CM2,ColdingMinicozzi19_Dynamics,CM7}. 
	To describe this system, let $\Sigma\subset\mathbb{R}^{n+1}$ be a hypersurface, and define its Gaussian area as follows:
	$$
	\cF(\Sigma):=(4\pi)^{-n/2}\int_{\Sigma}e^{-\frac{|x|^2}{4}}d\cH^n(x),
	$$ where $\cH^n$ represents the $n$-dimensional Hausdorff measure. In this context, a rescaled mean curvature flow $M_t$ is the negative gradient flow of the Gaussian area, and a shrinker is a fixed point of this flow. 
	
	In a pioneering work \cite{CM1}, Colding-Minicozzi proved that the only stable shrinkers from a variational point of view are the sphere and the generalized cylinders $\cC_{n,k}:=\mathbb{S}^{n-k}(\sqrt{2(n-k)})\times\R^k$. Based on their result, there has been much recent research on generic mean curvature flows that avoid other shrinkers as the singularity models; c.f. \cite{ColdingMinicozzi19_Dynamics, SX1, SX2, CCMS, CCS, BamlerKleiner23_multiplicity}. On the other hand, although the generalized cylinders are stable from a variational point of view, the singular sets modeled by them can still be very complicated. One famous example is the ``marriage ring'', namely, a very thin rotationally symmetric torus in $\R^3$, and it shrinks to the circle as the singular set, showing that the singular set can have high dimensions. Another example is the ``peanut surface'' constructed by Altschuler-Angenent-Giga, \cite{AAG}, showing that a compact mean curvature flow can shrink to a point while the tangent flow at the point is the noncompact cylinder.
	
	In the problem list \cite{I1}, Ilmanen proposed the following conjecture. This conjecture was also noted by Colding-Minicozzi-Pedersen in \cite[Conjecture 7.1]{CMP}.
	
	\begin{conjecture}[Generic isolatedness conjecture]\label{ConIsolate}
		Mean curvature flow with generic initial data only has isolated spacetime singularities. 
	\end{conjecture}
	
	
	In this paper, we study the cylindrical singularities to establish the necessary tool to study the dynamics of cylindrical singularities and we present a perspective from dynamics to hint at why this conjecture should be true; in another forthcoming paper (which is the second part of the first version of this paper), we will discuss further dynamical properties of the generic cylindrical singularities. Our focus is on the mean curvature flow of closed embedded hypersurfaces in $\R^{n+1}$. The basic assumption is that 
	\begin{equation}\label{eq:StarCondition}\tag{$\star$}
		\parbox{5.5in}{%
			\emph{$\{\bM_\tau\}_{\tau\in[-1,0)}$ is a mean curvature flow, with bounded entropy, with a cylindrical singularity modeled on $\cC_{n,k}$ at the spacetime point $(0,0)$, and $\{M_t\}_{t\in [0,\infty)}$ is its associated rescaled mean curvature flow that converges to $\cC_{n,k}$ in the $C^\infty_{loc}$-sense.}%
		}
	\end{equation}
	
	
	
	
	\subsection{The normal forms}
	To describe our main results, we start with some notations. We fix $1\leq k< n$. Let $\varrho:=\sqrt{2(n-k)}$, and we use coordinates denoted as $x=(\theta, y)\in \cC_{n,k}:=\mathbb{S}^{n-k}(\varrho)\times\R^k \subset \mathbb{R}^{n-k+1}\times\mathbb{R}^k$. Here, $\{y_i\}_{i=1}^k$ denotes the coordinates on $\mathbb{R}^k$, and $\{\theta_j\}_{j=1}^{n-k+1}$ denotes the restriction of the coordinate functions of $\mathbb{R}^{n-k+1}$ to $\mathbb{S}^{n-k}(\varrho)$. We use $B_R$ to denote the Euclidean open ball of radius $R$ centered at the origin.
	
	On $\cC_{n,k}$, it is natural to introduce the Gaussian weighted $L^2$-space, because the linearized operator $L_{\cC_{n,k}}$ of the generalized cylinder is self-adjoint with respect to the Gaussian weighted inner product. Throughout this paper, for any domain $\Omega\subset\cC_{n,k}$ and $f:\Omega\to\R$, we define
	\[
	\|f\|^2_{L^2(\Omega)}=\int_{\Omega}|f(x)|^2 e^{-\frac{|x|^2}{4}}d\cH^n(x),
	\quad 
	\|f\|^2_{H^1(\Omega)}=\int_{\Omega}(|f(x)|^2 +|\nabla f(x)|^2)e^{-\frac{|x|^2}{4}}d\cH^n(x).
	\]
	
	Our first main theorem is the asymptotic profile of the rescaled mean curvature flow converging to a generalized cylinder. Such an expression is called a {\bf normal form}, which is a terminology in dynamical systems after Birkhoff. Throughout this paper, we use the big $O$ notation as follows: given a positive function $f:\R\to\R$, $O(f(t))$ means a term satisfying $\limsup\limits_{t\to\infty}|O(f(t))|\cdot (f(t))^{-1}<+\infty$
	
	\begin{theorem}[$H^1$-normal form theorem]\label{ThmNF-sqrtt} Assume \eqref{eq:StarCondition}. Then  for any $K>0$ and $\vartheta\in(0,1)$, there exists $T>0$  such that for $t>T$, the rescaled mean curvature flow is a graph over $\cC_{n,k}\cap B_{K\sqrt{t}}$, and up to a rotation in $\R^k$, the graphical function $u(\cdot,t):\ \cC_{n,k}\cap B_{K\sqrt{t}}\to \R$ has the following asymptotic in (weighted) $H^1$-norm
		\begin{equation}
			\left\|u(\theta,y,t)-\sum_{i\in \mathcal I}\frac{\varrho }{4t}(y_i^2-2)\right\|_{H^1(\cC_{n,k}\cap B_{K\sqrt{t}})}=O(t^{-1-\vartheta})
		\end{equation}
		as $t\to\infty$, where $\mathcal I\subset \{1,2,\ldots,k\}$. Moreover, when $\mathcal I=\emptyset$, $\|u(\cdot,t)\|_{H^1(\cC_{n,k}\cap B_{K\sqrt{t}})}=O(e^{-K^2t})$ for some $K>0$. 
	\end{theorem}
	
	We also prove the following $C^1$-normal form theorem.
	
	\begin{theorem}[$C^1$-normal form theorem]\label{ThmNF-sqrtt-C1}
		Assume \eqref{eq:StarCondition}. Then for any $K>0$ and $\vartheta\in(0,1)$, there exists $T>0$ such that for $t>T$, the rescaled mean curvature flow is a graph over $\cC_{n,k}\cap B_{K\sqrt{t}}$, and the graphical function $u(\cdot,t):\ \cC_{n,k}\cap B_{K\sqrt{t}}\to \R$
		has the following asymptotic, up to a rotation in $\R^k$, in $C^1$-norm in $B_{K\sqrt{t}}$
		\begin{equation}
			\left\|u(\theta,y,t)-\varrho\left(\sqrt{1+\frac{ \sum_{i\in \mathcal I}(y_i^2-2)}{2t}}-1\right)\right\|_{C^1(\cC_{n,k}\cap B_{K\sqrt{t}})}=O(t^{-\vartheta}),
		\end{equation}
		as $t\to\infty$, where $\mathcal I\subset \{1,2,\ldots,k\}$. 
	\end{theorem}
	Combining the above two normal forms with the classical parabolic Schauder theory, we have the following estimate: 
	
	\begin{corollary}[$C^2$-normal form on a bounded domain]\label{ThmNF-sqrtt-C1bdd}
		For any $R>2n$, $\vartheta\in(0,1)$, there exists $T>0$ such that for $t>T$, the rescaled mean curvature flow is a graph over $\cC_{n,k}\cap B_{R}$, and the graphical function $u(\cdot,t):\ \cC_{n,k}\cap B_{R}\to \R$
		has the following asymptotic, up to a rotation in $\R^k$, in $C^2$ norm
		\begin{equation}
			u(\theta,y,t)=\sum_{i\in \mathcal I}\frac{\varrho }{4t}(y_i^2-2)+O(t^{-1-\vartheta}),
		\end{equation}
		as $t\to\infty$, where $\mathcal I\subset \{1,2,\ldots,k\}$. 
	\end{corollary}
	
	The normal form theorems generalize previous results of Gang Zhou in a series of works \cite{G1, G2} for the case of $n=4$ and $k=3$. However, the novelty in our work lies in the significantly simpler proof of the normal form. Our approach consists mainly of two ingredients: the analysis of the dynamical system for the rescaled mean curvature flow near the cylinder and the regularizing effect of the Ornstein-Uhlenbeck operator discovered by Vel\'azquez \cite{V1, V2}.
	
	
	We would like to remark that the $H^1$-normal form and $C^1$-normal form capture the behavior of the rescaled mean curvature flow in different regions: the $H^1$-normal form mainly captures the behavior inside the region of radius $O(t^{\kappa})$ with $\kappa$ close to $0$, while the $C^1$-normal form mainly captures the behavior near the boundary of the region of radius $O(t^{1/2})$.
	
	
	It is worth mentioning that in the investigation of ancient flows, Angenent-Daskalopoulos-Sesum \cite{ADS1} made the novel discovery that the asymptotic form can provide insights into the geometry of ancient flows. Furthermore, in recent research on the classification of ancient solutions of mean curvature flow \cite{ADS1, ADS2, BC1, BC2, CHH, CHHW, DuZhu22_spectral}, the asymptotic expansion of the rescaled mean curvature flow over the cylinders at $-\infty$ time plays a crucial role. However, it is important to note that the backward expansion differs from the forward expansion presented in this work, yielding distinct geometric information. For example, the linear asymptotic modes in the study of ancient flows form a finite-dimensional space, while in the study of forward flow, the linear asymptotic modes form an infinite-dimensional space.
	
	\subsection{Nondegeneracy and isolatedness}
	From the normal form that we derived, a natural definition arises.
	
	\begin{definition}\label{DefDeg}
		A cylindrical singularity as in Theorem \ref{ThmNF-sqrtt} $($equivalently as in Theorem \ref{ThmNF-sqrtt-C1}$)$ is \textbf{nondegenerate} if $\mathcal I=\{1,2,\ldots,k\}$, is {\bf partially nondegenerate} if $\mathcal I\subsetneqq\{1,2,\ldots,k\}$, and is \textbf{degenerate}  if $\cI=\emptyset$.   
	\end{definition}
	
	The concept of nondegeneracy for neckpinching singularities was initially introduced and studied by Angenent-Vel\'azquez \cite{AV} in the context of rotationally symmetric mean curvature flows. In the rotationally symmetric setting, the cylindrical singularities are modeled by $\cC_{n,1}$. Angenent-Vel\'azquez also constructed examples of mean curvature flows with degenerate singularities. It is worth noting that Schulze-Sesum \cite{SS} also defined a notion of nondegenerate neckpinch, namely a singularity modeled by $\cC_{n,1}$. Their nondegenerate singularities are characterized by the blow-up of limits. It seems that our notion of nondegeneracy for $\cC_{n,1}$ implies theirs, but not vice versa: in fact, some of those degenerate singularities constructed by Angenent-Vel\'azquez in \cite{AV} seem to satisfy Schulze-Sesum's nondegenerate notion in \cite{SS}.
	
	We discovered several geometric and dynamic properties of nondegenerate singularities. The first significant result from the normal form theorem is the following theorem concerning the isolatedness of nondegenerate cylindrical singularities.
	
	\begin{theorem}[Isolatedness theorem]\label{ThmIsolate}
		Assume \eqref{eq:StarCondition} and the singularity is nondegenerate, then the singularity is isolated. In other words, there exists $\dt_1>0$ such that $(0,0)$ is the only singularity of $\bM_\tau$ in the spacetime neighborhood $B_{\dt_1}\times(-\delta_1^{1/2},0]$. 
	\end{theorem}
	
	
	Another relevant property is the mean convex neighborhood. It was conjectured that any cylindrical singularity has a mean convex neighborhood. This conjecture was proved by Choi-Haslhofer-Heshkovitz \cite{CHH} for dimension $n=2$, by Choi-Haslhofer-Heshkovitz-White \cite{CHHW} for dimension $\geq 3$ with a $2$-convexity assumption, and by Gang Zhou \cite{G1} for the singularity model $\mathbb{S}^1\times\mathbb{R}^3$ with a nondegeneracy assumption. 
	
	We prove that all the nondegenerate cylindrical singularities have a mean convex neighborhood.
	
	\begin{theorem}[Mean convex neighbourhood of a nondegenerate singularity]\label{prop:Mconnbhd}
		Assume \eqref{eq:StarCondition} and that the singularity is nondegenerate. Then there exist constants $\dt_2>0$ and $\tau_0>0$, such that for any $\tau\in(-\tau_0,0)$, $\bM_\tau\cap B_{\dt_2}(0)$ has positive mean curvature and is diffeomorphic to $\cC_{n,k}$. 
	\end{theorem}
	
	The mean convex neighborhood property is related to the nonfattening property of level set flow. In \cite{HershkovitsWhite20_Nonfattening}, Hershkovits-White proved that if all the singularities are of mean convex/mean concave type, then the level set flow does not fatten. The condition in Theorem \ref{prop:Mconnbhd} is sufficient to show that a nondegenerate cylindrical singularity is of such a type. Thus, if a mean curvature flow only has nondegenerate singularities, its level set flow does not fatten. 
	
	
	Another crucial property of nondegenerate singularities is the type-I curvature condition. A singularity $(y,T)$ of a mean curvature flow is called \textbf{type-I} if the curvature blows up with a speed of at most $O((T-\tau)^{-1/2})$. This means that there exist constants $r>0$ and $C>0$ such that for any $t<T$, in $M_t\cap B_r(y)$, we have $|A|(x,t)=O((T-\tau)^{-1/2})$, where $A$ represents the second fundamental form. On the other hand, if the curvature blows up faster than this rate, the singularity is called \textbf{type-II}. There are examples of type-II singularities; for instance, in \cite{AAG}, the authors constructed a surface called a ``peanut," and the mean curvature flow starting from such a peanut develops type-II singularities known as ``degenerate neckpinching." In addition, Angenent-Vel\'azquez \cite{AV} constructed a large family of type-II singularities in their work.
	
	\begin{theorem}[Type-I curvature condition of a nondegenerate singularity]\label{ThmtypeI}
		Assume \eqref{eq:StarCondition} and that the singularity is nondegenerate. Then there exist constants $\dt_2>0$ and $\tau_0>0$, such that for any $\tau\in(-\tau_0,0)$, $\bM_\tau\cap B_{\dt_2}(0)$ has curvature bound $|A|\leq C(-\tau)^{-1/2}$ for some constant $C$.
	\end{theorem}
	
	\subsection{Subsequent work}
	In a subsequent work with Zhihan Wang \cite{SunWangXue1_Passing}, we further explored some finer properties of nondegenerate singularities, based on the normal form theorems of this paper. Particularly, with Zhihan Wang, we proved that Theorems \ref{ThmIsolate}, \ref{prop:Mconnbhd}, \ref{ThmtypeI} hold not only in the backward parabolic neighborhood, but also hold in the forward parabolic neighborhood. We also have a complete description of the change of geometry and topology of the mean curvature flow passing through a nondegenerate singularity. As a consequence, we obtain the topological information of the mean curvature flow with only nondegenerate singularities.
	
	In \cite{SunWangXue2_RegSing}, we further studied those singularities that are partially nondegenerate and degenerate, showing that their singular sets have some regularity.
	
	\subsection{Strategy of proof}
	Let us first comment on two important ingredients of the proof. The first is the \textbf{pseudolocality} theorem of mean curvature flow. It is a nonlinear effect that was first discovered by Ecker-Huisken in the setting of mean curvature flow. The pseudolocality shows that if a mean curvature flow is sufficiently close to a cylinder in a sufficiently large domain, then in a definite amount of time, it behaves similarly to the flow of a shrinking cylinder. In particular, when we work with the rescaled mean curvature flow, the pseudolocality shows that if a rescaled mean curvature flow is sufficiently close to a cylinder in a sufficiently large domain, then in a definite amount of time, it will be close to the cylinder inside a larger domain, and its higher order derivatives will have better control. This is the key to the nonlinear estimate to study the rescaled mean curvature flows, even if the nonlinear error is much more complicated than the semilinear equations.
	
	The second is Vel\'azquez's argument, which is also known as the \textbf{Ornstein-Uhlenbeck regularization}. For a nonnegative function $Z(\cdot,t)$ defined on $\cC_{n,k}$ satisfying the inequality
	\[
	\pr_t Z-LZ\leq \eps_0 Z+\text{terms decaying in time}+\text{terms from cut-off},
	\]
	this tool provide an estimate of $\|Z(\cdot,t')\|_{C^0}$ inside ball of radius $r(t)=O(t^{1/2})$ by $(r(t))^2\|Z(\cdot,t)\|_{L^2}$ for $t'>t$. 
	
	In applications, we will choose $Z=|\chi u|+|\nabla (\chi u)|$, where $u$ is the graph function of the rescaled mean curvature flow over $\cC_{n,k}$ and $\chi$ is a cut-off function; we will also choose $Z=|\chi w|+|\nabla (\chi w)|$, where $w$ is the difference of $u$ and the function given by the normal form. Then Vel\'azquez's argument provides, roughly speaking, a $C^1$-control of the graph by the $H^1$-norm. In particular, the $C^1$-control is the ingredient to apply the pseudolocality theorem. In order to get a very small $C^1$-bound, we need a decay of $\|Z(\cdot,t)\|_{L^2}$ faster than $(r(t))^{-2}$. Hence, the graphical radius relies on the $H^1$-norm of the graph function. That is the reason that we need to bootstrap between the $C^1$-norm and the $H^1$-norm.
	
	
	The proof is divided into the following steps. The initial input is the $L^2$-bound of the graphical function (decay like $t^{-\xi}$, $\xi>0$ that can be very small) and the graphical radius $\br(t)\geq C\sqrt{(\log t)}$ by \cite{CM2} obtained using the {\L}ojasiewicz inequality. 
	\begin{enumerate}
		\item We improve the graphical radius to $\br(t)=O(t^{\kappa})$ for $2\kappa<\xi$ as follows: starting at some large time $T$, we use the pseudolocality of mean curvature flow to improve the rough graphical radius to $T+\delta$. The graph function $u$ can have very large $C^0$-norm, but $|\nabla^k u|$ is small for $k=1,2,3$. Then we use Vel\'azquez's argument to show that the graph function indeed has a small $C^1$-norm. Then we iteratively use the pseudolocality and Vel\'azquez's argument to extend the graphical radius to all time $t>T$.
		\item Using the methods of dynamical systems, we obtain an $H^1$ normal form over $B_{\br(t)}$ and improve the $H^1$-estimate to $\|u-\frac{\varrho}{4t}\sum (y_i^2-2)\|_{L^2}=O(t^{-1-\vartheta})$. The dynamical argument is similar to the authors' previous work \cite{SX1, SX2, SX3}. To derive the explicit normal form, we follow the idea of Vel\'azquez \cite{V1} in the study of semilinear PDE.
		\item Let $w=u-\frac{\varrho}{4t}\sum (y_i^2-2)$. Then use the existing proof in Section \ref{SS_Extend_Graphical_Radius_to_K_0t^1/2} to extend the graphical radius of $w$ (hence $u$) to $K_0\sqrt{t}$, for some small $K_0$. $K_0$ can only be a small number for the following reasons: we need $\frac{\varrho}{4t}\sum (y_i^2-2)$ to be very small to ensure that $\|w\|_{C^2}$ small can imply $\|u\|_{C^2}$ small. 
		\item Then we use pseudolocality to prove $|\nabla^k u|$ bound on the ball of radius $(1+\alpha) K_0\sqrt{t}$, $\alpha>0$, for $k=1,2,3$. With this bound, we can write the equation for the graph of rescaled mean curvature flow on the ball of radius $(1+\alpha) K_0\sqrt{t}$. This allows us to apply Vel\'azquez's argument to prove the $C^1$-normal form on the ball of radius $(1+\alpha) K_0\sqrt{t}$.
		\item Once we got the $C^1$-normal form on $(1+\alpha) K_0\sqrt{t}$, we repeatedly using pseudolocality to prove $|\nabla^k u|$ bound on the ball of radius $(1+\alpha)^m K_0\sqrt{t}$ for $m\geq 1$ and $k=1,2,3$. Particularly, we have the $C^1$-normal form on the ball of radius $K\sqrt{t}$ for any $K>0$.
	\end{enumerate}
	
	
	
	\subsection{Organization of the paper}
	Section \ref{S_preliminaries} contains some preliminaries; in Section \ref{S_OU_Estimate}, we present the adaptation of Vel\'azquez's Ornstein-Uhlenbeck regularization; in Section \ref{S:Extending_t_kappa} we proved the first step, showing that the graphical radius is $O(t^\kappa)$ for some $\kappa\in(0,1/2)$; in Section \ref{SDynamics}, we use the dynamical argument to study the evolution of the $0$-eigenmodes, and prove the $H^1$-normal form; in Section \ref{SS:The global variational equation: scale up to}, we prove the $C^1$-normal form based on the $H^1$-normal form and the Ornstein-Uhlenbeck regularization; in Section \ref{S_Geometric} we studied some geometric properties of the nondegenerate singularities. We also discuss some nonlinear estimates in the Appendix.
	
	
	
	\subsection*{Acknowledgment}
	We would like to thank Professor Bill Minicozzi for the stimulating discussions, and Professor Natasa Sesum for her interest and comments. We also thank Zhihan Wang for valuable suggestions and comments -- some of them lead to new statements and arguments in the current version. J. X. is supported by NSFC grants (No. 12271285) in China, the New Cornerstone investigator program, and the Xiaomi endowed professorship of Tsinghua University. 
	
	\section{Preliminaries}\label{S_preliminaries}
	In this section, we provide some background and preliminary results and set up some notations for later sections. 
	
	\subsection{Eigenvalues and Eigenfunctions of the $L$-operator}\label{SPre}
	In this section, we summarize some previously known results on cylindrical singularities. 
	We will be working in the weighted Sobolev space. Given two functions $f,g$ defined on a hypersurface $\Sigma$, we define an inner product
	$
	\langle f,g\rangle_{L^2} =\int_{\Sigma}f(x)g(x)e^{-\frac{|x|^2}{4}}d\cH^n(x).
	$
	Then we define the weighted $L^2$-norm by
	$\|f\|_{L^2}=\langle f,f\rangle_{L^2}^{1/2},$
	and the weighted $L^2(\Sigma)$ space consists of function $f$ with $\|f\|_{L^2}<+\infty$. Similarly, we define the weighted higher Sobolev space $H^{k}$ with
	$
	H^{k}(\Sigma):=\left\{f: \sum_{i=0}^k \|\nabla^i f\|^2_{L^2}<+\infty\right\}.
	$
	Throughout the paper, for simplicity of notations, we use $\|\cdot\|$ to denote the $H^1$-norm if not otherwise mentioned. 
	
	Recall that the linearized operator on a shrinker is defined as
	$L:=\Delta-\frac{1}{2}\langle x,\nabla\cdot\rangle+(|A|^2+1/2).$
	In the special case that the shrinker is $\cC_{n,k}$, we have
	\begin{equation}\label{EqL}
		L_{\cC_{n,k}} u= \Delta_{\bS^{n-k}(\varrho)} u +L_{\R^k}u,\quad L_{\R^k}u=\Delta_{\R^{k}} u -\sum_{i=1}^{k}\frac{1}{2}y_i\pr_{y_i}u+u.	
	\end{equation}
	Throughout this paper, an eigenvalue $\lambda$ of an operator $-L$ is a number such that there exists a nonzero function $f$ satisfying $Lf+\lambda f=0$. The following fact was proved in \cite[Section 5.2]{SunWangZhou20_MinmaxShrinker}: \emph{Suppose the eigenvalues of $-\Delta_{\bS^{n-k}(\varrho)}$ are given by $\mu_0\leq \mu_1\leq \mu_2\leq \cdots$ with corresponding eigenfunctions $\phi_0,\phi_1,\phi_2,\cdots$, and suppose the eigenvalues of $L_{\R^k}$ on $\R^{k}$ is given by $\nu_0\leq \nu_1\leq \nu_2\leq \cdots$ with corresponding eigenfunctions $\psi_0,\psi_1,\psi_2, \cdots$, then
		the eigenvalues of $L_{\cC_{n,k}}$ are given by $\{\mu_i+\nu_j\}_{i,j=0}^\infty$ with corresponding eigenfunctions $\{\phi_i\psi_j\}_{i,j=0}^\infty$.}
	
	The spectrum of $-\Delta_{\bS^{n-k}(\varrho)}$ is $0,1/2,\frac{n-k+1}{n-k},\cdots$, and the eigenfunctions are known to be the restriction of homogeneous harmonic polynomials. The first several eigenfunctions are listed as follows: constant functions for eigenvalue $0$; $\theta_i$, the restriction of linear functions in $\R^{n-k+1}$ to $\bS^{n-k}(\varrho)$, for eigenvalue $1/2$; and $\theta_i^2-\theta^2_j,\cdots$ for eigenvalue $\frac{n-k+1}{n-k}$. The spectrum of $-L_{\R^k}$ on $\R^{k}$ is given by half integers $\frac{m}{2}-1,\ m=0,1,2,\ldots$, and the eigenfunctions for eigenvalue $\frac{m}{2}-1$ are given by $h_{m_1}(y_1)\cdots h_{m_k}(y_k)$ with $m_1+\ldots+m_k=m$, where $h_{m_i}(y_i)=c_{m_i}\tilde h_{m_i}(y_i/2)$, and $\tilde h_{m_i}$ are standard Hermite polynomials with $c_{m_i}=2^{-m_i/2}(4\pi)^{-1/4}(m_i!)^{-1/2}$ are normalizing factors such that $\|h_{m_i}\|_{L^2(\R)}=1$. In particular, we have $$\tilde h_0(x)=1,\ \tilde h_1(x)=x,\ \tilde h_2(x)=4x^2-2 \mathrm{\ and\ } c_0=(4\pi)^{-1/4}, c_1=4^{-1/2}\pi^{-1/4}, c_2=2^{-2}\pi^{-1/4}.$$ We shall use the following fact (c.f. Appendix B of \cite{HV2}).
	\begin{lemma}\label{LmTriple}
		Let $A_{m,n,\ell}=\int_\R h_m(x)h_n(x)h_\ell(x)e^{-\frac{|x|^2}{4}}dx.$ Then $A_{m,n,\ell}=0$ unless we have $m+n+\ell$ is even and $n\leq m+\ell,\ m\leq n+\ell,\ \ell\leq m+n$, in which case we have $$A_{m,n,\ell}=(4\pi)^{-1/4}(m!n!\ell!)^{1/2}\left(\left((\frac{m+n-\ell}{2}\right)!\left(\frac{n+\ell-m}{2}\right)!\left(\frac{m+\ell-n}{2}\right)!\right)^{-1}.$$ In particular, we have $A_{2,2,2}=2\pi^{-1/4}=8c_2.$
	\end{lemma}
	
	
	Combining the spectra together, we obtain the first three eigenvalues and their corresponding eigenfunctions of $L_{\cC_{n,k}}$, see Table \ref{TableEigen}. These eigenfunctions have geometric meanings: 
	\begin{itemize}
		\item Constant $1$ is the mean curvature on the generalized cylinder, representing infinitesimal (spacetime) dilation.
		\item $\theta_i$ and $y_j$ are both infinitesimal translations. Specifically, $\theta_i$'s represent the translations in the directions of the spherical components, and $y_i$'s represent the translations in the axis directions.
		\item $\theta_i y_j$ represents the infinitesimal rotation.
		\item $h_2(y_i)=c_2(y_i^2-2)$ is known to be the non-integrable Jacobi field. It represents \emph{non-degenerate neckpinching} in the related direction on the axis. 
		\item $y_iy_j$'s show up if we rotate $h_2(y_i)$'s in the $\R^k$ space. e.g. $\left(\frac{y_i+y_j}{\sqrt2}\right)^2-2=\frac{y_i^2-2}{2}+\frac{y_j^2-2}{2}+\sqrt{2}y_iy_j$.
	\end{itemize}
	
	\begin{table}[H]
		\begin{tabular}{|l|l|}
			\hline
			eigenvalues of $-L_{\cC_{n,k}}$ & corresponding eigenfunctions \\ \hline
			$-1$ & $1$ \\ \hline
			$-1/2$ & $\theta_i,y_j,\ i=1,2,\ldots,n-k+1,\ j=1,2,\ldots,k$ \\ \hline
			$0$ & $\theta_iy_j,\ h_2(y_j)=c_2(y_j^2-2),\ y_{j_1}y_{j_2}$ \\ \hline
			$\max\{1/(n-k),1/2\}$ & $\ldots$ 
			\\ \hline
		\end{tabular}
		\caption{Eigenvalues and eigenfunctions of $-L_{\cC_{n,k}}$.}
		\label{TableEigen}
	\end{table}
	
	Note that while $h_{m_i}(y)$ is a normalized eigenfunction in $\|\cdot\|_{L^2(\R)}$, $h_{m_i}(y_i)$ is not a normalized eigenfunction in $\|\cdot\|_{L^2(\cC_{n,k})}$. Thus, we define $H_{m_1,m_2,\cdots,m_k}$ be a multiple of $h_{m_1}h_{m_2}\cdots h_{m_k}$, such that $H_{m_1,m_2,\cdots,m_k}$ is a normalized eigenfunction in $\|\cdot\|_{L^2(\cC_{n,k})}$. In particular, $H_2(y_i)=c_0^{k-1}\cG_{n,k}^{-1/2}h_2(y_i)$, $H_{1,1}(y_i,y_j)=c_0^{k-2}\cG_{n,k}^{-1/2}h_1(y_i)h_1(y_j)$. Here \begin{equation}\label{EqcG}
		\cG_{n,k}=\int_{\mathbb S^{n-k}(\varrho)}e^{-\frac{|x|^2}{4}}dx=\varrho^{n-k}e^{-\frac{\varrho^2}{4}}\omega_{n-k},        
	\end{equation}
	where $\omega_{n-k}$ is the area of the $(n-k)$-dimensional unit sphere.
	
	
	
	\subsection{Coarse graphical scale and $L^2$ estimate}\label{SSCM}
	
	By Colding-Minicozzi's {\L}ojasiewicz inequality \cite{CM2}, if one tangent flow of a bounded entropy mean curvature flow $\bM_{\tau}$ at $(0,0)$ is $\cC_{n,k}$, then \eqref{eq:StarCondition} holds. In \cite{CM2}, Colding-Minicozzi introduced the notion of \emph{cylindrical scale}.	
	
	\begin{definition}[Graphical radius; called cylindrical scale in \cite{CM2}]\label{DefGraphRad} Let $\eps_0>0,\ \ell\geq 2$ and $C_{\ell}$ be some fixed constants.
		Given a hypersurface $\Sigma$, the cylindrical scale $($throughout the paper, we call it the \emph{graphical radius}$)$ $\br(\Sigma)$ is the largest radius such that $\Sigma \cap B_{\br(\Sigma)}$ is the graph a function $u:\ \cC_{n,k}\cap B_{\br(\Sigma)}\to \R$ with $\|u\|_{C^{2,\alpha}}\leq \eps_0$ and the part of hypersurface $\cC_{n,k} \cap B_{\br(\Sigma)}$ has curvature bound $|\nabla^\ell A|\leq C_{\ell}$. 
	\end{definition}
	
	Because we will work with various graphical regions, and the behavior of the rescaled mean curvature flow can be different in different regions, we may \emph{not} always choose $\br$ to be the largest radius.
	
	We need some results from \cite{CM2} to serve as initial input for our proofs. The result is summarized in the following proposition. 
	\begin{proposition}\label{PropR}
		Assume \eqref{eq:StarCondition} and let $\br(t)$ be the graphical radius of the rescaled mean curvature flow $M_t$. Then there exists $t_0$ such that we have the estimate $\br(t)\geq (\al\log t)^{1/2}$ for some small $\al>0$ and all $t>t_0$. Moreover, choosing $r(t)= (\al\log t)^{1/2}$, the graphical function $u:\ \cC_{n,k}\cap B_{r(t)}\to \R$ has the estimate  $\|u(\cdot,t)\|_{H^1(B_{r(t)})}\leq t^{-\frac{3}{8}}$ for $t>t_0$. 
	\end{proposition}
	\begin{proof} 
		In Theorem 6.1 of \cite{CM2} we choose $1/\tau=3-\epsilon$, then we get the estimate $F(\Sigma_{t-1})-F(\Sigma_{t+1})\leq t^{-\frac{1}{\tau}}=t^{-(3-\epsilon)}$ by Lemma 6.9 of \cite{CM2} (we remark that the statement of Theorem 6.1 of \cite{CM2} gives only some $\tau\in (1/3,1)$. However, the proof there is constructive. It is not hard to see that $\tau$ can be chosen close to $1/3$. See \cite[Footnote 10]{CM2} on page 244). 
		Defining $R(t)$ (shrinker scale in \cite{CM2}) through $e^{-\frac{R(t)^2}{2}}=F(\Sigma_{t-1})-F(\Sigma_{t+1})$,
		we get $R(t)\ge \sqrt 2\sqrt{\frac{1}{\tau}\log t} .$
		By Theorem 5.3 of \cite{CM2} on the graphical scale $\br(t)\geq (1+\mu)R(t)$ for some $\mu>0$, we get the estimate $$\br(t)\geq \frac{\sqrt2}{\sqrt\tau}(1+\mu)\sqrt{\log t}=\sqrt{2(3-\eps)} (1+\mu)\sqrt{\log t}.$$ 
		
		To estimate $\|u\|_{H^1(B_{r(t)})}$, we integrate the square of \cite[(2.51)]{CM2} with respect to the Gaussian area to yield (with $r(t)$ possibly getting smaller by a constant)
		$$\|u(t)\|^2_{H^1(B_{r(t)})}\leq C r(t)^\rho(\|\phi\|_{L^1(B_{r(t)})}^{1-\eps}+e^{-(1-\epsilon)\frac{r(t)^2}{4}})
		\leq C(\al\log t)^{\rho/2}
		(\|\phi\|^{1-\eps}_{L^1(B_{r(t)})}+
		t^{-\frac{3}{2}(1+\mu)^2(1-\eps)}
		)
		,$$
		for some uniform constants $C$ and $\rho$, and arbitrarily small $\epsilon$,
		where $\phi$ is a function with the estimate (see equation (6.6) of \cite{CM2})
		$$\|\phi\|^4_{L^1(B_{r(t)})}\leq C \|\phi\|^2_{L^2(B_{r(t)})}\leq C(F(\Sigma_{t-1})-F(\Sigma_{t+1}))\leq t^{-1/\tau}.$$
		Combining the two estimates, we get the estimate of $\|u\|_{H^1(B_{r(t)})}\leq t^{-3/8}$. 	\end{proof}
	We remark that the graphical radius growth rate and the decay rate of $u$ are uniform in entropy for all cylindrical singularities (see footnote 5 of \cite{CM2} on page 223). 
	
	
	\subsection{Pseudolocality}\label{SSPL}
	The following pseudolocality argument was first studied by Ecker-Huisken in \cite{EH}. Later, Ilmanen-Neves-Schulze \cite[Section 9]{INS} gave a simple proof using Brakke-White regularity theorem \cite{Wh4}. There, they studied the mean curvature flow of graphs over a hyperplane, and we study the mean curvature flow of graphs over a cylinder.
	
	\begin{theorem}\label{thm:pseudo-locality}
		For any $\eps\in(0,1)$, there exist $R_0>0$ and $\eta_0>0$ with the following significance: suppose $M_t$ is a rescaled mean curvature flow and $M_{t_0}$ is the graph of a function $u(\cdot,t_0)$ over $\cC_{n,k}\cap B_R(p)$ for some $R>R_0$, $p\in\{0\}\times\R^k$, and $\|u(\cdot,t_0)\|_{C^1(B_R(p))}\leq\eta\leq \eta_0$. Then there exists $\bar\delta >0$ such that when $t\in[t_0,t_0+\bar\delta]$, $M_t$ is the graph of a function $u(\cdot,t)$ over $\cC_{n,k}\cap B_{e^{(t-t_0)/2}(R-1)}(e^{(t-t_0)/2}p)$ with $$\|\nabla^ku(\cdot,t)\|_{L^\infty(\cC_{n,k}\cap B_{e^{(t-t_0)/2}R}(e^{(t-t_0)/2}p))}\leq e^{-(k-1)(t-t_0)/2}\eps, \quad k=0,1,2,3.$$ Moreover, $\bar\delta\to\infty$ as $\eta\to 0$.
		
	\end{theorem}
	
	In the theorem, the factor $e^{-(k-1)(t-t_0)/2}$ is from the exponential expansion property of the rescaled mean curvature flow. In particular, when $k=0$, the $L^\infty$ norm of $u$ increases very fast. We need to use the equation of the rescaled mean curvature flow $\partial_t u=Lu+\cQ$ to show that the growth rate is actually not that fast.
	
	Theorem \ref{thm:pseudo-locality} follows from the following theorem for mean curvature flow, considering the scaling between mean curvature flow and rescaled mean curvature flow. 
	\begin{theorem}
		For any $\tau_0\in(-1,0)$ and $\eps\in(0,1)$, there exists $\eta_0>0$ such that for any $\eta\in(0,\eta_0)$, if $\bM_{-1}\cap B_{\eta^{-1}}$ is the graph of a function $v$ over $\cC_{n,k}$ with $\|v\|_{C^1}\leq \eta$, then 
		\begin{enumerate}
			\item $\bM_{\tau}\cap B_{\eta^{-1}-1}$ is smooth mean curvature flow for $\tau\in(-1,\tau_0)$;
			\item $\bM_{\tau}\cap B_{\eta^{-1}-1}$ has a singularity before time $-\tau_0$;
			\item $\bM_{\tau}\cap B_{\eta^{-1}-1}$ can be written as the graph of a function $u(\cdot,\tau)$ over $\sqrt{-\tau}\cC_{n,k}$ for $\tau\in(-1,\tau_0)$, with $\|u(\cdot,\tau)\|_{C^3}\leq \eps$.
		\end{enumerate}
	\end{theorem}
	
	\begin{proof}
		We prove by contradiction. Suppose there is no $\eta_0$ such that all of the three items hold, then we can find a sequence of hypersurfaces $\bM_{-1}^k$, such that each one of them is the graph of a function $v_k$ inside the ball of radius $k$, with $\|v_k\|_{C^1}\leq k^{-1}$, but at least one of the three items fails. Then we can use the compactness of Brakke flow (see \cite[Section 7]{I2}) to pass to a limiting weak mean curvature flow $\bM_\tau^{\infty}$. By Brakke-White's regularity theorem \cite{Wh4}, and Ecker-Huisken's curvature estimate, $\bM_\tau^{\infty}$ is a smooth mean curvature flow when $\tau$ is close to $-1$. Also, the sequence $\bM_{-1}^k$ has varifold limit being the standard cylinder $\cC_{n,k}$. Thus, the limit flow $\bM_\tau^{\infty}$ is exactly the shrinking soliton, with only one singularity at time $0$. Then by the upper semi-continuity of Gaussian density, we know that the singularity of $\bM_{\tau}^k$ must have the singular time converging to $0$, and $\bM_{\tau}^k$ are graphs over $\sqrt{-\tau}\cC_{n,k}$ with $C^3$-norm converges to $0$. Then all three items hold, which is a contradiction.
	\end{proof}
	
	
	
	
	\subsection{RMCF graph equation and cut-off}	\label{SSCone}
	To study the asymptotics of the rescaled mean curvature flow, we need to study the evolution of the graphical function. 
	In Proposition \ref{PropQ} in Appendix \ref{appendix:Nonlinear estimate}, we write the graph of the rescaled mean curvature flow over a shrinker given by
	\[
	\pr_t u=Lu+\cQ(J^2u)
	\]
	where $u(\cdot, t):\ \cC_{n,k}\cap B_{\mathbf r(t)}\to \R$ whose graph is the rescaled mean curvature flow $M_t$ restricted to the ball $B_{\mathbf r(t)},$ and $\cQ$ is at least quadratic in $u$ given in \eqref{EqQ} and $J^2u:=(u,Du,D^2u)$ means the $2$-jet of $u$. We shall study the evolution of $u$ under the differential equation. 
	
	Since $M_t$ cannot be written as a global graph over $\cC_{n,k}$, we introduce a smooth cutoff function $\tilde \chi(t):\ \R_{\geq 0}\to \R$ that is $1$ over $[0,1]$ and vanishes outside $[0,2]$, so that $|\nabla^\ell \tilde \chi|\leq C_{\ell}$ for any $\ell\in\{0,1,2\}$. Then for any $r>1$ we define $\chi_r:\cC_{n,k}\to\R$ by 
	\[
	\chi_r(\theta,y)=
	\begin{cases}
		1,& y\leq r,
		\\
		\tilde \chi(|y|-r+1), & y\in[r,r+1],
		\\
		0,& y\geq r+1.
	\end{cases}
	\]
	In particular, given a differentiable function $f:\R\to\R_{\geq 1}$, we have $\pr_t \chi_{f(t)}(\theta,y)$ is only nonzero in $[f(t),f(t)+1]$, and $|\pr_t \chi_{f(t)}(\theta,y)|\leq C_1|f'(t)|$. In particular, when $f(t)=C(\log t)^{\alpha}$ or $t^\kappa$ for $\alpha>0$ and $\kappa<1$, we have $|\pr_t \chi_{f(t)}(\theta,y)|\to 0$ as $t\to\infty$.
	
	Given a function $\mathbf{r}:\R\to\R_{\geq 0}$, we define $\mathbb A_{\mathbf r(t)}$, or simply $\A_t$ if $\mathbf{r}$ is the graphical function and is fixed, to be the annulus region $\cC_{n,k}\cap(B_{\br(t)}\backslash B_{\br(t)-1})$. Then we derive the equation for $\chi u$:
	\begin{equation}\label{EqCutOff}
		\pr_t (\chi u)
		=
		L(\chi u)+\chi\cQ(J^2u)+(u(\Delta\chi+\pr_t\chi)+2\langle \nabla\chi,\nabla u\rangle +\frac{1}{2}\langle x,\nabla \chi\rangle u),
	\end{equation}
	where the last term on the RHS is supported on the annulus $\mathbb A_{\mathbf r(t)}$. For simplicity, we write
	\begin{equation}\label{Eq}
		\pr_t (\chi u)
		=
		L(\chi u)+\cB(J^2u).
	\end{equation}
	
	
	
	
	\section{Vel\'azquez's Ornstein-Uhlenbeck regularization estimate}\label{S_OU_Estimate}
	
	One of the key technical ingredients in this paper is the Ornstein-Uhlenbeck regularization for the solutions to the equation $\pr_t u=\Delta u-\frac{1}{2}\sum_{j=1}^k y_j \pr_{y_j}u+u$ studied by Vel\'azquez in \cite{V1}; c.f. \cite{HV1,HV2}. This estimate gives $C^0$-bounds over the region with expanding radius, provided we have a bound on the $L^2$-norm. The main result that we shall prove in this section is the following proposition. 
	\begin{proposition}\label{PropVelazquez}
		Given $\kappa \in (0,1/2]$, there exist constants $\bar\eps>0$, $K_0>0$, $T_0>0$ such that for all $\eps_0<\bar \eps$, $K<K_0$ the following holds: suppose  $Z(\cdot,t):\cC_{n,k}\times[T_0,\infty)\to\R_{\geq 0}$ satisfies the following differential inequality (in the sense of distribution) \begin{equation}\label{EqVelazInequality}
			\begin{split}
				\pr_t Z-L Z
				\leq
				\eps_0 Z+C\left(\frac{|y|^2+1}{t^2}+|y|\check\chi(y,t)\right),
			\end{split} 
		\end{equation}
		where $\check\chi(y,t)$ is a smooth cutoff function that is $1$ outside $B_{Kt^{\kappa}+1}$, and $0$ inside $B_{Kt^{\kappa}}$.
		For any $t_0>T_0$, suppose we have $ \|Z(\cdot,t)\|_{L^2(\cC_{n,k})}\leq C_0t^{-\xi}$ for $\xi>0$, when $t\in[t_0,t_0+\mathbf{T}]$ with $\mathbf{T}\in(2,2\log(K t_0^\kappa)]$. Then we have the estimate 
		$\|Z(\cdot,t)\|_{C^0(Kt^\kappa)}\leq C(C_0)t^{2\kappa-\xi}$ when $t\in[t_0+2,t_0+\mathbf{T}]$. 
	\end{proposition}
	
	
	We remark that the only place where we need $\eps_0<\bar\eps$ to be small is the proof of Lemma \ref{LmJ134}.
	
	While the linear effect of the drift Laplacian is the same, our problem is very different from the semilinear heat equation $\partial_t u=\Delta u+u^p$ considered by Vel\'azquez due to the nonlinear part. In particular, our nonlinearity contains second-order derivatives. In later proofs, we shall recast our problem into the form of the differential inequality \eqref{EqVelazInequality} that differs from the differential inequality considered in \cite{V1}. Moreover, our $L$-operator contains a spherical Laplacian summand. For the convenience of readers, we provide detailed proof in this section. 
	
	
	
	Suppose $S(t,s)$ is the heat kernel of $\partial_t u=Lu$ such that $u(t)=S(t,s)u(s)$. Then we can write \eqref{EqVelazInequality} as the following integral inequality
	$$Z(\cdot, t)\leq S(t,t_0)Z(\cdot,t_0)+\int_{t_0}^t S(t,s)\left(\eps_0 Z(\cdot,s)+\frac{|y|^2+1}{s^2}+|y|\check{\chi}(y,s)\right)ds.$$
	
	To estimate $\|Z\|_{C^0(Kt^\kappa)}$, we need a good understanding of the behavior of the heat kernel, which is done in the first subsection below. The integral of $S(t,s)(\frac{|y|^2+1}{s^2}+|y|\check{\chi}(y,s))$ is estimated by explicit integrations. We shall show that the estimate of $\|Z\|_{C^0(Kt^\kappa)}$ is mainly dominated by $S(t,t_0)Z(\cdot,t_0)$. The heat kernel estimate is summarized in the following proposition, which will be proved in Section \ref{SSOU}.  We will show in Section \ref{SSVelaz} that the integral of $S(t,s)\eps_0 Z(\cdot,s)$ is small compared to the linear part $S(t,t_0)Z(\cdot,t_0)$ thus completes the proof of Proposition \ref{PropVelazquez}. 
	
	A particular application of this section is the following $C^0$-estimate of the linearized rescaled mean curvature flow equation, Proposition \ref{PropBoundary-global}. Although we will prove stronger estimates for nonlinear equations, the linear analysis may have independent interest, and we include it here.
	
	\begin{proposition}[\cite{V1}, see also Proposition 2.13 of \cite{KZ}]\label{PropBoundary-global}
		Given $\kappa>0$ and $K_0>0$, there exist $T_0>0$, $C>0$ such that for all $t_0>T_0$, 
		if $v(\cdot,t):\ \cC_{n,k}\times[t_0,t_0+\mathbf{T}]\to \R$ be a nonnegative solution to the equation
		$\pr_tv= Lv$
		with initial condition $\|v(\cdot,t_0)\|_{L^2(\cC_{n,k})}<\infty$ and $\mathbf{T}\in(2,2\log (K_0 t_0^\kappa)]$. 
		Then we have
		\begin{equation*}
			\begin{aligned}
				\|v(\cdot,t)\|_{C^0(B_{K_0t^{\kappa}})}
				&\leq C t^{2\kappa}\|v(\cdot,t_0)\|_{L^2},
			\end{aligned}
		\end{equation*}
		whenever $t\in[t_0+2,t_0+\mathbf T]$.
	\end{proposition}
	\subsection{Vel\'azquez's inequality for the heat kernel of the drift Laplacian}\label{SKZ}
	
	
	The main goal is understanding the evolution of the nonnegative function $Z$ satisfies $\pr_t Z\leq LZ$. The starting point is the heat kernel of the operator $L_{\cC_{n,k}}$. Following \cite{V1}, we use $S(\tau)$ to denote the semigroup generated by $L_{\cC_{n,k}}$ over $\cC_{n,k}$, and we introduce the kernel of the semigroup generated to the operator $L_{\cC_{n,k}}$ on $\cC_{n,k}$
	\begin{equation}\label{EqS}
		S(\tau,(\theta,y),(\eta,z))
		=G(\tau,\theta,\eta)\frac{e^\tau}{(4\pi(1-e^{-\tau}))^{k/2}}
		\exp\left(-\frac{|ye^{-\tau/2}-z|^2}{4(1-e^{-\tau})}\right)
	\end{equation}
	where $G$ is the heat kernel for the semigroup associated to $\Delta_{\bS^{n-k}(\varrho)}$ on $\bS^{n-k}(\varrho)$. Define the following norm
	\begin{equation}
		N_r^q(\Psi)=\sup_{|\xi|\leq r}\left(\int_{\cC_{n,k}}
		|\Psi(\theta,y)|^qe^{-\frac{|y-\xi|^2}{4}}d\theta dy
		\right)^{1/q}.
	\end{equation}
	Particularly, we use $N_r$ to denote $N^2_r$, and it is clear that $N^2_r(v)\leq \|v\|_{L^2}$. The first important lemma is the generalization of Vel\'azquez inequality \cite{V1}.
	
	\begin{lemma}\label{lem:estNr}
		For $q_0,q\in (1,\infty)$, $q'=\frac{q_0}{q_0-1}$ and $r_0\geq 0$, $r\geq 0$, we have the following estimate
		\begin{equation}\label{eq:estNr}
			\begin{aligned}
				N^{q}_{r}(S(\tau)\psi)
				&\leq C(n,k)
				\frac{e^{\tau}}{(4\pi(1-e^{-\tau}))^{k/2+\frac{n-k}{2q_0}}}\left(\frac{4\pi q_0 (1-e^{-\tau})}{q'(q_0-(1-e^{-\tau}))} \right)^{k/(2q')}\\
				&\left(\frac{4\pi(q_0- (1-e^{-\tau}))}{(q_0-1)-(q-1)e^{-\tau} } \right)^{k/(2q)}
				\exp\left(
				\frac{e^{-\tau}(r-r_0e^{\tau/2})_+^2}{4((q_0-1)-(q-1)e^{-\tau})}
				\right)N^{q_0}_{r_0}(\psi).\\
			\end{aligned}
		\end{equation}
	\end{lemma}
	
	
	
	\begin{remark}\label{RkVelaz}
		In later applications, we will be focusing on some specific choices of $q$ and $q_0$:
		\begin{enumerate}
			\item  We shall mainly use the lemma for the two cases $q_0 > q$ and $q_0=q$. We have the following  bounds   provided we have a bound on $e^{-\tau/2}r$
			$$N^q_r(S(\tau)\psi)\leq 
			\frac{ C e^{\tau}}{(4\pi(1-e^{-\tau}))^{n/(2q_0)}} N^{q_0}_{r_0}(\psi), \quad (q_0>q),$$
			$$N^q_r(S(\tau)\psi)\leq 
			\frac{ C e^{\tau}}{(4\pi(1-e^{-\tau}))^{n/(2q_0)+k/(2q)}} N^{q_0}_{r_0}(\psi), \quad (q_0=q).$$
			\item We will need the integral $\int_0^t N^q_r(S(\tau)\psi)d\tau$ to be bounded. For that purpose, we shall take $q_0$ to be large, so that the factor $(1-e^{-\tau})^{-\frac{n}{2q_0}}$ is integrable near $\tau=0$ for $2q_0>n$. 
		\end{enumerate}
	\end{remark}
	
	
	The proof is almost the same as \cite{V1}, with the only difference being the spherical component. 
	
	\begin{proof}
		For any fixed $\xi\in\R^k$, let us consider
		\[
		\begin{split}
			I
			=&
			\int_{\cC_{n,k}}
			|S(\tau)\psi(\theta,y)|^q e^{-\frac{|y-\xi|^2}{4}}d\theta dy
			=
			\int_{\cC_{n,k}}
			e^{-\frac{|y-\xi|^2}{4}}
			\left(
			\int_{\cC_{n,k}} S(\tau,(\theta,y),(\eta,z))\psi(\eta,z)d\eta dz
			\right)^q d\theta dy
			\\
			=&
			\left(\frac{e^\tau}{(4\pi(1-e^{-\tau}))^{k/2}}\right)^q
			\int_{\cC_{n,k}}
			e^{-\frac{|y-\xi|^2}{4}}
			\left(
			\int_{\cC_{n,k}} G(\tau,\theta,\eta)
			\exp\left(-\frac{|ye^{-\tau/2}-z|^2}{4(1-e^{-\tau})}\right)
			\psi(\eta,z)d\eta dz
			\right)^q d\theta dy	.
		\end{split}
		\]
		Now let $w\in\R^k$ be any point satisfying $|w|\leq r_0$. Using H\"older's inequality, we obtain that 
		\[
		\begin{split}
			&
			\left(
			\int_{\cC_{n,k}} G(\tau,\theta,\eta)
			\exp\left(-\frac{|ye^{-\tau/2}-z|^2}{4(1-e^{-\tau})}\right)
			\psi(\eta,z)d\eta dz
			\right)^q
			\\
			\leq &
			\left(
			\int_{\cC_{n,k}} 
			e^{-\frac{|z-w|^2}{4}}
			|\psi(\eta,z)|^{q_0} d\eta dz
			\right)^{q/q_0}
			\left(
			\int_{\cC_{n,k}} G^{q'}(\tau,\theta,\eta)
			\exp\left(-\frac{q'|ye^{-\tau/2}-z|^2}{4(1-e^{-\tau})}+\frac{q'|z-w|^2}{4q_0}\right)d\eta dz
			\right)^{q/q_0'}
			\\
			=&
			\left(
			\int_{\cC_{n,k}} 
			e^{-\frac{|z-w|^2}{4}}
			|\psi(\eta,z)|^{q_0} d\eta dz
			\right)^{q/q_0}
			\left(
			\int_{\bS^{n-k}} G^{q'}(\tau,\theta,\eta)
			d\eta
			\right)^{q/q'}
			\\
			&\qquad \cdot
			\left(
			\int_{\R^k} 
			\exp\left(-\frac{q'|ye^{-\tau/2}-z|^2}{4(1-e^{-\tau})}+\frac{q'|z-w|^2}{4q_0}\right) dz
			\right)^{q/q'}.
		\end{split}
		\]
		We also have the identity
		\[
		\begin{split}
			-\frac{q'|ye^{-\tau/2}-z|^2}{4(1-e^{-\tau})}+\frac{q'|z-w|^2}{4q_0}
			=&
			-\frac{q'(q_0-(1-e^{-\tau}))}{4q_0(1-e^{-\tau})}
			\left|
			z-w-\frac{q_0 (ye^{-\tau}-w)}{q_0-(1-e^{-\tau})}
			\right|^2
			+\frac{q'e^{-\tau}|y-w e^{\tau/2}|^2}{4(q_0-(1-e^{-\tau}))}.
		\end{split}
		\]
		
		Therefore, we obtain that
		\[
		\begin{split}
			I^{1/q}
			\leq &
			N^{q_0}_{r_0}(\psi)
			\left(
			\int_{\bS^{n-k}}\left(
			\int_{\bS^{n-k}} G^{q'}(\tau,\theta,\eta)
			d\eta
			\right)^{q/q'}d\theta\right)^{1/q}
			\left(\frac{e^\tau}{(4\pi(1-e^{-\tau}))^{k/2}}\right)
			\\
			&
			\left(\frac{4\pi q_0 (1-e^{-\tau})}{q'(q_0-(1-e^{-\tau}))} \right)^{k/(2q')}
			\left(\frac{4\pi(q_0- (1-e^{-\tau}))}{(q_0-1)-(q-1)e^{-\tau})} \right)^{k/(2q)}
			\exp\left(
			\frac{e^{-\tau}(r-r_0e^{\tau/2})_+^2}{4((q_0-1)-(q-1)e^{-\tau})}
			\right).
		\end{split}
		\]
		The simplification can be found in \cite[page 1573]{V1}. The only part that was not calculated by Vel\'azquez is the integral over the spherical part. By the decay order of the heat kernel on the sphere (for example, see \cite{NSS19_SphereHeatKernel}) and the $L^{q'}$ upper bound of the Gaussian function, we have 
		\begin{align*}
			\left(
			\int_{\bS^{n-k}}\left(
			\int_{\bS^{n-k}} G^{q'}(\tau,\theta,\eta)
			d\eta
			\right)^{q/q'}d\theta\right)^{1/q}
			\leq &
			\left(
			\int_{\bS^{n-k}}\left(
			C\tau^{-(1-1/q')(n-k)/2}
			\right)^{q}d\theta\right)^{1/q}
			=
			C \tau^{-\frac{n-k}{2q_0}}.
		\end{align*}
		
		Thus we obtain the lemma.
	\end{proof}
	
	
	\bigskip
	
	\subsection{The Ornstein-Uhlenbeck regularization}\label{SSOU}
	
	In this section, we use the inequality of the previous section to prove Proposition \ref{PropBoundary-global}. 
	\begin{proof}[Proof of Proposition \ref{PropBoundary-global}]
		
		Let $Z=|v|$, then by Kato's inequality, $\pr_t Z\leq LZ$. For $t_0$ and $\mathbf T$ to be determined, we estimate $N_{K_0s^{\kappa}}(Z(s))$ for $s \in[t_0+\delta,\mathbf T]$ with $\delta=1/2$. Suppose $K_0$ is a fixed constant, and $s_0$ is chosen such that $e^{(s_0-t_0)/2}\leq K_0 {s_0}^{\kappa}$. Notice that whenever $K_0=t_0^{-\kappa}$, $s_0=t_0$, and for any fixed $K_0$, when $t_0$ is sufficiently large, we always have $s_0>K_0$. We also define $n'=n/4+k/4$. Denote by $\mtr(s,t)=e^{(s-t)/2}$, and when $t=t_0$, we simplify the notation $\mtr(s)=e^{(s-t_0)/2}$. Because $\frac{\pr Z}{\pr t}\leq LZ$, we have
		\[
		Z(s)\leq S(s-t_0)Z(t_0)=S(s-t_0-\delta)S(\delta)Z(t_0).
		\]
		
		Then in the $N_r$-norm we have
		\[
		N_{\mtr(s)}(Z(s))\leq
		N_{\mtr(s)} (S(s-t_0-\delta)S(\delta)Z(t_0)).
		\]
		We observe $\mtr(t_0+\delta)e^{(s-t_0-\delta)/2}=\mtr(s)$, which implies $(\mtr(s)-\mtr(t_0+\delta)e^{(s-t_0-\delta)/2})_+= 0$. Lemma \ref{lem:estNr} shows that
		\[
		N_{\mtr(s)}(Z(s))\leq 
		C\frac{e^{s-t_0-\delta}}{(1-e^{-(s-t_0-\delta)})^{n'}}
		N_{(\mtr(t_0+\delta))} (S(\delta)Z(t_0)).
		\]
		Finally, we use Lemma \ref{lem:estNr} again to obtain
		\[
		N_{\mtr(t_0+\delta)}(S(\delta)Z(t_0))
		\leq
		C\frac{e^{\delta}}{(1-e^{-\delta})^{n'}}\exp\left(\frac{e^{-\delta}\mtr(t_0+\delta)^2}{4(1-e^{-\delta})}\right)\|Z(t_0)\|_{L^2}.
		\]
		Combining the estimates above we get
		\begin{equation}\label{eq:NrboundbyL2}
			\begin{split}
				N_{\mtr(s)}(Z(s))\leq& C\frac{e^{s-t_0}}{(1-e^{-(s-t_0-\delta)})^{n'}(1-e^{-\delta})^{n'}}
				\exp\left(\frac{1}{4(1-e^{-\delta})}\right)\|Z(t_0)\|_{L^2}.
			\end{split}
		\end{equation}
		
		\medskip
		Next, we use $N_{K_0s^{\kappa}}(Z(s))$ to bound $Z(\cdot,s')$ in the domain $\{|y|\leq K_0(s')^\kappa\}$ with some $s<s'$. 
		
		
		\begin{lemma}\label{LmC0byL2}
			We have  
			$
			Z(\theta,y,s')\leq \tilde C(d)N_{K_0s^{\kappa}}(Z(s))
			$ for $s=s'-d$, $s'\geq \frac{d}{1-e^{-\frac{d}{2\kappa}}}$, and some constant $\tilde C(d)$.
		\end{lemma} 
		
		We prove this lemma later, and at this moment, we conclude the proof by using this lemma. Using Lemma \ref{LmC0byL2} and \eqref{eq:NrboundbyL2}, we get that when $|y|\leq K_0(s')^{\kappa}$,
		\begin{equation}\label{EqC0byL2}
			Z(y,\theta,s')
			\leq
			C(d,\delta)\frac{e^{s-t_0}}{(1-e^{-(s-t_0-\delta)})^{n'}(1-e^{-\delta})^{n'}}
			\exp\left(\frac{1}{4(1-e^{-\delta})}\right)\|Z(t_0)\|_{L^2}.
		\end{equation}
		
		
		In conclusion, we have
		\begin{equation}
			\|v(\cdot,s')\|_{C^0(B_{K_0(s')^{\kappa}})}
			\leq C(d,\delta)s^{2\kappa}\|v(\cdot,t_0)\|_{L^2}.
		\end{equation}
		
		Therefore, if we choose $d$ sufficiently close to $0$ (so that $\frac{d}{1-e^{-\frac{d}{2\kappa}}}$ is sufficiently close to $2\kappa\leq 1$) and $\delta=1/2$, then for $t_0$ sufficiently large such that $t_0>\frac{d}{1-e^{-\frac{d}{2\kappa}}}$ and $\mathbf T$ solved from $e^{\mathbf{T}/2}\leq K_0 t_0^\kappa$, we have \begin{equation*}
			\begin{aligned}
				\|v(\cdot,t)\|_{C^0(B_{K_0t^{\kappa}})}
				&\leq C t^{2\kappa}\|v(\cdot,t_0)\|_{L^2},
			\end{aligned}
		\end{equation*}
		whenever $t\in[t_0+2,t_0+\mathbf T]$.
		
		
	\end{proof}
	\begin{proof}[Proof of Lemma \ref{LmC0byL2}]   
		We use the heat kernel of $L_{\cC_{n,k}}$ and the upper bound of $G$ to conclude
		\[
		\begin{split}
			0\leq&
			Z(\theta,y,s')
			\leq S(d)Z(\theta,y,s'-d)
			\\
			\leq&
			C(d)\int_{\cC_{n,k}} G(d,\theta,\eta)\exp\left(-\frac{|y e^{-d/2}-z|^2}{4(1-e^{-d})}\right)
			Z(\eta,z,s'-d)d\eta dz.
		\end{split}
		\]
		Now for any $|w|\leq K_0(s'-d)^{\kappa}$, we have
		\[
		Z(\theta,y,s')
		\leq
		C(d)\int_{\cC_{n,k}} G(d,\theta,\eta)\exp\left(-\frac{|y e^{-d/2}-z|^2}{4(1-e^{-d})}
		+\frac{|z+w|^2}{8}
		\right)e^{-\frac{|z+w|^2}{8}}
		Z(\eta,z,s'-d)d\eta dz.
		\]
		Then Cauchy-Schwarz inequality yields that
		\[
		Z(\theta,y,s')
		\leq
		C(d)
		\left(\int_{\cC_{n,k}}
		|Z(\eta,z,s'-d)|^2e^{-\frac{|z+w|^2}{4}}
		d\eta dz
		\right)^{1/2}
		I(d,w,y)^{1/2},
		\]
		where
		\[
		\begin{split}
			I(d,w,y)
			=&
			C\int_{\cC_{n,k}}(G(d,\theta,\eta))^2
			\exp\left(-2\frac{|y e^{-d/2}-z|^2}{4(1-e^{-d})}
			+\frac{|z+w|^2}{4}
			\right) d\eta dz
			\\
			=&
			C(d)
			\exp\left(
			\frac{|w+ye^{-d/2}|^2}{2(1+e^{-d})}
			\right).
		\end{split}
		\]
		If we take the infimum in $w$ with $|w|\leq K_0(s'-d)^{\kappa}$, and then take the supremum in $y$ with $|y|\leq K_0(s')^{\kappa}$, whenever $s'\geq \frac{d}{1-e^{-\frac{d}{2\kappa}}}$, we get 
		\[\sup_{|y|\leq K_0(s')^{\kappa}}\left(\inf_{|w|\leq K_0(s'-d)^{\kappa}}|w+ye^{-d/2}|\right)=0.\]
		Thus, we have proved the lemma. 
	\end{proof}
	
	
	\subsection{Estimates of the nonlinear equation, Proof of Proposition \ref{PropVelazquez}}\label{SSVelaz}
	In this section, we consider the nonlinear estimate for \eqref{EqVelazInequality} and complete the proof of Proposition \ref{PropVelazquez}. As the passage from the $N^2$ estimate to the $C^0$ estimate was given Lemma \ref{LmC0byL2} and \eqref{EqC0byL2} of the last section, it is sufficient to prove the following. 
	
	\begin{proposition}\label{prop_OU_var_const}
		Under the assumption of Proposition \ref{PropVelazquez}, we have the following estimate when $t\in[t_0+2,t_0+\mathbf{T}]$: 
		$$N_r^2(Z(\cdot,t))\leq C\epsilon e^{t-t_0}/t^\xi.$$
	\end{proposition}
	\begin{proof}
		
		
		
		Recall that $\mathbf T$ is chosen such that $e^{{\mathbf T}/{2}}\leq  Kt_0^\kappa$ and we consider $t\in [t_0+2,\mathbf T]$. We should understand this choice of $t$ as the space expansion $e^{\frac{t-t_0}{2}}$, due to the passage from mean curvature flow to rescaled mean curvature flow, which arrives at the scale $t^\kappa$ after time $\mathbf T. $
		Note that for large $t_0$, we have $t/t_0\approx 1$. 
		Using the method of variation of constants, we get
		\begin{equation}\label{EqDuHamel}
			Z(\cdot,t)\leq \left(S(t-t_0)Z(\cdot,t_0)+\int_{t_0}^t S(t-s)\left(\eps_0 Z(\cdot,s)+C\frac{|y|^2+1}{s^2}+C|y|\check\chi(\cdot,s)\right)ds  \right).    
		\end{equation}
		
		Taking the $N^2_r$-norm with $r=\mtr(t,t_0)=e^{(t-t_0)/2}$, we get
		\begin{equation}
			\begin{aligned}
				&N_r^2(Z(\cdot,t))\leq N^2_r(S(t-t_0)Z(\cdot,t_0))
				+
				\int_{t_0}^t N^2_r(S(t-s)\eps_0Z(\cdot,s))ds\\
				&+
				C\int_{t_0}^tN^2_r\left(S(t-s)\left(\frac{|y|^2+1}{s^2}+|y|\check\chi(\cdot,s)\right)\right)ds=:J_1+J_2+J_3.
			\end{aligned}
		\end{equation}
		$J_1+J_3$ can be estimated as follows, and we will postpone the proof to the end of this section. 
		\begin{lemma}\label{LmJ134}
			For $t\in[t_0+2,t_0+\mathbf {T}]$, we have $J_1+J_3\leq\bar C\epsilon e^{t-t_0}t^{-\xi}$ for some $\bar C>0$. We have a similar bound if we replace the $N^2_r$ norm by $N^q_r$-norm in $J_1+J_3$ for any $q>2.$
		\end{lemma}
		
		
		
		
		For $J_2$, we will bounding $N_{\mtr(t)}^2$ by $N_{\mtr(s)}^{q}$ for a large $q>n/2$. The reason for doing so is given in item (2) of Remark \ref{RkVelaz}, since we want an integrable prefactor $\frac{e^{\tau}}{(1-e^{-\tau})^{n/q }}$ near $\tau=0$. 
		Thus, we get 
		\begin{equation}\label{EqN2}
			\begin{aligned}
				&N_{\mtr(t)}^2(Z(\cdot,t))\leq \bar C\epsilon e^{t-t_0}/t^\xi+ \eps_0\int_{t_0}^{t}  \frac{e^{(t-s)}}{(1-e^{-(t-s)})^{n/q }} N^{q}_{\mtr(s)}(Z(\cdot,s))ds.   
			\end{aligned}
		\end{equation}
		We next estimate $N^{q}_{\mtr(s)}(Z)$ on the RHS. From equation \eqref{EqDuHamel}, we have 
		\begin{equation*}
			\begin{aligned}
				N_{r(t)}^{q}(Z(\cdot,t))&\leq  N_{r(t)}^{q}(S(t-t_0)Z(\cdot,t_0))+\int_{t_0}^t N_{r(t)}^{q}(S(t-s)(\eps_0Z(\cdot,s))\\
				&+C\int_{t_0}^tN_{r(t)}^{q}\left(S(t-s)\left(\frac{|y|^2+1}{s^2}+|y|\check\chi(\cdot,s)\right)\right)ds=:J_1'+J_2'+J_3'.
				& \end{aligned}   
		\end{equation*}
		We next bound $J_1'+J_3'$ by $\hat C\epsilon e^{t-t_0}/t^\xi$ for the same reason as Lemma \ref{LmJ134}. Next, we get
		$$N_{\mtr(t)}^{q}(Z(t,y))\leq \hat C\epsilon e^{t-t_0}/t^\xi+\eps_0 \int_{t_0}^t\frac{e^{(t-s)}}{(1-e^{-(t-s)})^{\frac{n}{q}}} N^{q}_{\mtr(s)}(Z(s,\cdot))ds.$$
		Note that the last equation is in a closed form for $N^{q}$. We shall choose $q>n/2.$ We then conclude that $N_{\mtr(t)}^{q}\leq 2\hat C\epsilon e^{t-t_0}/t^\xi$. Indeed, assuming that $N_{\mtr(t)}^{q}(t)\leq 2\hat C\epsilon e^{t-t_0}/t_0^\xi$ for $t\in [t_0,t_1]$, substituting to the above integral inequality, we find $N_{\mtr(t)}^{q}(t)\leq (1+C\eps_0)\hat C\epsilon e^{t-t_0}/t_0^\xi$, where $C$ only depends on the integral upper bound $t$. If $\eps_0$ is small so that $C\eps_0<1$, we can extend the estimate $N_{\mtr(t)}^{q}(t)\leq 2\hat C\epsilon e^{t-t_0}/t^\xi$ beyond time $t_1.$
		Substituting this back to equation \eqref{EqN2}, we get $N_{\mtr(t)}^2\leq 2\bar C\epsilon e^{t-t_0}/t^\xi.$ 
		
		
		
	\end{proof}
	
	
	\begin{proof}[Proof of Lemma \ref{LmJ134}]
		We apply Lemma \ref{lem:estNr} with $q=q_0=2$, $r=\mtr(t,t_0)$ and $r_0=\mtr(t_0+2,t_0)$ to get
		$$|J_1|=|N_r^2(S(t-t_0-1)S(1)Z( t_0,\cdot))|\leq Ce^{t-t_0-1} N^2_{\mtr(t_0+2,t_0)}(S(1)Z(t_0,\cdot)).$$
		We apply lemma \ref{lem:estNr} again with $q=q_0=2$, $r=\mtr(t_0+2,t_0),\ r_0=0$ to get 
		$$RHS\leq Ce^{t-t_0}\|Z(t_0,\cdot)\|_{L^2}\leq Ce^{t-t_0}/t^\xi\leq C\epsilon  t^{2\kappa-\xi} $$
		where the bound $\|Z(t_0,\cdot)\|_{L^2}\leq C \epsilon/t^\xi$ follows from the assumption and $t/t_0\to1$ when $t_0$ is sufficiently large. 
		
		To prove the part of the statement with $N^q_r$-norm in place of the $N^2_r$-norm. We can either apply Lemma \ref{lem:estNr} with $q>q_0=2$, or apply the $C^0$ estimate \eqref{EqC0byL2} of $Z(t)$ and calculate the $N^q_r$-norm from there. 
		
		For $J_3$, We  choose $q=2,\ q_0=20n,\ r=\mtr(t,r_0),\ r_0=\mtr(s,r_0)$ in Lemma \ref{lem:estNr}, then we have the following.
		$$N_r^2\left(S(t-s)\left(\frac{y^2+1}{s^2}\right)\right)\leq \frac{c e^{t-s}}{(1-e^{-(t-s)})^{1/40}}\left(N^{20n}_{r_0}\left(\frac{y^2+1}{s^2}\right)\right)\leq \frac{c e^{t-s}}{(1-e^{-(t-s)})^{1/40}} \frac{1}{s^2}(1+e^{s-t_0}).$$
		Next, we choose $t_0$ large enough with $s\leq 2t_0$ to get 
		$$\int_{t_0}^tN^2_r\left(S(t-s)(\frac{|y|^2+1}{s^2})\right)ds\leq \frac{C}{t_0^2}\int_{t_0}^t (e^{t-t_0}+e^{t-s})(1-e^{-(t-s)})^{1/40}ds\leq 
		C K^2 t^{2\kappa-2}. $$
		Finally, we get
		\begin{equation*}
			\begin{aligned}
				\int_{t_0}^tN^2_r\left(S(t-s)(|y|\check\chi(\cdot,s))\right)ds&\leq C\int_{t_0}^t \frac{c e^{t-s}}{(1-e^{-(r-s)})^{1/40}}(N^{20n}_{r_0}(|y|\check\chi(\cdot, s))ds\\
				&\leq C\int_{t_0}^t \frac{c e^{t-s}}{(1-e^{-(t-s)})^{1/40}}\left(\int_{|\lambda|\geq K s^{\kappa}}\lambda^{20n} e^{-\lambda^2/4}d\lambda\right)^{1/(20n)}ds.         
			\end{aligned}
		\end{equation*}
		The integral is mainly bounded by the term in the parentheses, i.e. $\leq e^{-t^{2\kappa}/100n}$. Thus we get the stated estimate.
	\end{proof}
	
	As an application, we also show that the square of the quadratic integral can bound the quartic integral of $Z$.
	
	\begin{lemma}\label{lem_OU_control_L4byL2}
		Suppose $Z$ satisfies the assumptions in Proposition \ref{PropVelazquez}. Then
		\begin{equation}\label{eq_in_lem_OU_control_L4byL2}
			\|Z(\cdot,t_0+2)\|_{L^4}\leq C(\|Z(\cdot,t_0)\|_{L^2}+e^{-t^{2\kappa}/100n}).
		\end{equation}
	\end{lemma}
	\begin{proof}
		We simply repeat the proof of Proposition \ref{prop_OU_var_const}, but replacing $N^2_r$-norm by the $N^4_r$-norm. By the variation of constants, we also have \eqref{EqDuHamel}. Then we take $N^4_r$-norm rather than $N^2_r$-norm to get 
		\begin{equation}\label{eq_DH_q=4_q_0=2bound}
			\begin{aligned}
				&N_r^4(Z(\cdot,t))\leq N^4_r(S(t-t_0)Z(\cdot,t_0))
				+
				\int_{t_0}^t N^4_r(S(t-s)\eps_0Z(\cdot,s))ds\\
				&+
				C\int_{t_0}^tN^4_r\left(S(t-s)\left(\frac{|y|^2+1}{s^2}+|y|\check\chi(\cdot,s)\right)\right)ds=:J_1+J_2+J_3.
			\end{aligned}
		\end{equation}
		We can apply Lemma \ref{lem:estNr} by choosing $q=4$, $q_0=2$, then $J_1$ on the right-hand side of \eqref{eq_DH_q=4_q_0=2bound} is bounded by $C\|Z(\cdot,t_0)\|_{L^2}$. The estimate of $J_3$ is the same as in the proof of Lemma \ref{LmJ134}, where we choose $q=4$ and $q_0=20n$ in Lemma \ref{lem:estNr}, same analysis shows that $J_3$ is bounded by $Ce^{-t^{2\kappa}/100n}$. Finally, the analysis of $J_2$ is the same as the proof of Proposition \ref{prop_OU_var_const}, but instead of $C\eps e^{t-t_0}/t^\xi$, we use the same estimate but bound $J_2$ by $C\eps e^{t-t_0}(\|Z(\cdot,t_0)\|_{L^2}+e^{-t^{2\kappa}/100n})$. Then we get the bound as \eqref{eq_in_lem_OU_control_L4byL2}.
	\end{proof}
	
	
	
	
	\begin{remark}
		While in this paper we only consider the graphical radius with $O(t^{\kappa})$ growth, one can also consider the exponential graphical radius growth when the $H^1$ norm of the graph function decays exponentially fast, which occurs when $\mathcal I=\emptyset$ in the normal form theorems. In \cite{AV}, Angenent-Vel\'azquez proved that in the rotationally symmetric case, the graphical radius of a degenerate singularity can be exponential. In \cite{SunWangXue2_RegSing}, joint with Zhihan Wang, we proved an exponential growth of the graphical radius for fully degenerate cylindrical singularities.
	\end{remark}
	
	
	
	\section{Extending the graphical radius to $O(t^{\kappa})$}\label{S:Extending_t_kappa}
	In this section, we first extend the graphical radius to a polynomial rate $O(t^{\kappa})$ for some small $\kappa>0$ using Proposition \ref{PropVelazquez}. We shall use the $L^2$-decay rate in Proposition \ref{PropR} as the input, which gives $\|u\|_{L^2}=O(t^{-\xi})=O(t^{-\frac{3}{8}})$ and we expect the graphical radius to be $O(t^{\kappa})$ with $\kappa=3/16-\eps$ as the output of Proposition \ref{PropR}. Thus, in this section, we devote most of the efforts to showing that $Z=|u|+|Du|$ satisfies the differential inequality \eqref{EqVelazInequality}. Note that the nonlinearity in our rescaled mean curvature flow equation depends on second-order derivatives of the graphical function $u$, which is the main difficulty for us to apply the machinery of Vel\'azquez. In this section, we show how to apply the pseudolocality to give a bound on higher-order derivatives of the graphical function over expanding domains.  
	
	\subsection{Extending the graphical radius}
	In this section, we extend the graphical radius from $O(\sqrt{\log t})$ to $O(t^\kappa)$, where $\kappa\in(0,1/2)$ is a fixed number, determined in Proposition \ref{PropBoundary-global}. 
	
	\begin{proposition}\label{thm:mainthmfor sqrtt}
		Suppose $M_t$ is a rescaled mean curvature flow converging to $\cC_{n,k}$ in the $C^\infty_{loc}$-sense. Then there exists $\eps_2>0$ such that for any $\eps\in(0,\eps_2)$, there exist $K>0$, and $T=T(\eps,\lambda(M_0))>0$, such that when $t>T$, $M_t$ can be written as the graph of a function $u(\cdot,t)$ over $\cC_{n,k}\cap B_{K t^\kappa}$, with $\|u(\cdot,t)\|_{C^2(B_{ K t^\kappa})}\leq \eps$.
	\end{proposition}
	\begin{lemma}\label{LmVelaz}
		Let $Z=|\chi_{Kt^\kappa}u|+|\nabla (\chi_{Kt^\kappa}u)|$. There exists $\eps_1\in(0,\eps_0)$, where $\eps_0$ is the one in Proposition \ref{SKZ}, such that if $|D^ku|\leq \eps_1,\ k=0, 1,2,3,$ over the ball $B_{Kt^\kappa+1}$, then for sufficiently large $t$, $Z$ satisfies the differential inequality \eqref{EqVelazInequality}.
	\end{lemma}   
	\begin{remark}
		The term $\frac{|y|^2+1}{t^2}$ is redundant in this section. This term will appear later when we prove the $C^1$-norm form. We shall treat both cases together in the proof. 
	\end{remark}
	
	\begin{proof}[Proof of Lemma \ref{LmVelaz}]
		The proof is based on an observation on the nonlinearity $\mathcal Q$ in Appendix \ref{appendix:Nonlinear estimate}. 
		Note that in $\mathcal Q$, the dependence on $Du$ is always of the form $|Du|^2$ with possibly further dependence on $D^2u.$ We use pseudolocality to bound $|D^ku|\leq \eps_1,\ k=0, 1,2,3,$ on the ball $B_{Kt^\kappa}$, and when $\eps_1$ is sufficiently small, $\max\{|\mathcal Q|,|D\mathcal Q|\}\leq \eps_0 (|u|+|Du|).$
		
		We next derive the differential equation for $\nabla u$. By $\pr_\theta(Lu)=L(\pr_\theta u)$, we can see that the equation for $\pr_\theta u$ is $$\partial_t (\pr_\theta u)=L (\pr_\theta u)+\pr_{\theta}\mathcal{Q}(J^2u).$$ When we take the derivative in the spine direction, note that there is a shift of frequency, i.e. $\pr_{y_i}(Lu)=L(\pr_{y_i}u)-\frac{1}{2}\pr_{y_i}u,$ thus, the equation for $\pr_{y_i}u$ is of the form $$\partial_t (\pr_{y_i}u)=L (\pr_{y_i} u)-\frac{1}{2}\pr_{y_i} u+\pr_{y_i}\mathcal{Q}(J^2u).$$
		The extra term $-\frac{1}{2}\pr_{y_i} u$ will be discarded when deriving the differential inequality since it has the correct sign. 
		
		Next, we introduce a cutoff function $\chi=\chi_{Kt^\kappa}$ and consider the equations of motion for $\chi u$ and $D(\chi u)$. As in \eqref{EqCutOff}, there is a term supported in the annulus region $B_{Kt^\kappa+1}\setminus B_{Kt^\kappa}$, for which we use $|y|\check\chi$ to give an upper bound. 
		
		Finally, recall Kato's inequality $\Delta g\cdot \text{sgn}(w)\leq \Delta|w|$ (in the sense of distribution), if we multiply $\pr_tw-(\Delta w-\frac{1}{2}y\cdot\nabla w+\ell w)$ with $\text{sgn}(w)$, for $\ell\in(0,1]$, we obtain 
		\[
		\left(\pr_tw-(\Delta w-\frac{1}{2}y\cdot\nabla w+\ell w)\right)\text{sgn}(w)
		\geq
		\pr_t|w|-(\Delta |w|-\frac{1}{2}y\cdot\nabla |w|+|w|)
		\]
		This leads to \eqref{EqVelazInequality} if we plug in $w=\chi_{Kt^\kappa}u$, $w=\pr_\theta (\chi_{Kt^\kappa}u)$ (with $\ell=1$) and $w=\pr_y(\chi_{Kt^\kappa}u)$ (with $\ell=1/2$).
	\end{proof}
	
	Now we come back to the proof of Proposition \ref{thm:mainthmfor sqrtt}. The proof uses the ``extension-improvement'' that Colding-Minicozzi \cite{CM2} used to prove the uniqueness of cylinder tangent flows. 
	
	\begin{proof}[Proof of Proposition \ref{thm:mainthmfor sqrtt}]
		We denote by $\xi=3/8$ and $\kappa\in(0,1/2)$ such that $2\kappa<\xi$. The proof is divided into several steps:
		
		\noindent{\bf Step 1. Initiating the setup.} Let us first fix $\eps_2$ to be determined, and let $\eps\in(0,\eps_2)$. Then we choose $T_0>0$ sufficiently large with the following significance: 
		\begin{itemize}
			\item Proposition \ref{PropVelazquez} is applicable.
			\item for $t>T_0$, $M_t$ is a graph of a function $u(\cdot,t)$ inside $\cC_{n,k}\cap B_{\bar R}$, such that $\|u(\cdot,t)\|_{C^1(B_{\bar R})}\leq \eta_0$, where $\bar R>R_0$ and $\eta_0$ as in the Pseudolocality Theorem \ref{thm:pseudo-locality} can be applicable for $\eps_2$. We also assume we $\eta_0$ is chosen very small such that $\bar\delta$ in Theorem \ref{thm:pseudo-locality} is much larger than $4$. We choose $K$ such that $K(T_0+2)^{\kappa}=\bar R$. In particular, this shows that $\|u(\cdot,t)\|_{C^1(B_{Kt^\kappa})}<\eta_0$ for $t\in[T_0,T_0+2]$.
		\end{itemize}
		
		\noindent{\bf Step 2. Extend the graphical region.}
		Apply the Pseudolocality theorem Theorem \ref{thm:pseudo-locality}, we have for $t\in[T_0+2,T_0+4]$, $u(\cdot,t)$ is a graph of function over $\cC_{n,k}\cap B_{(1+\alpha)Kt^{\kappa}}$, for some $\alpha>0$ depending on $\mathrm{t}$, such that $$\|\nabla^ku(\cdot,t)\|_{L^\infty(\cC_{n,k}\cap B_{(1+\alpha)Kt^{\kappa}})}\leq e^{-(k-1)(t-t_0)/2}\eps<\eps_1, \quad k=0,1,2,3.$$ Here we can choose $\eps$ sufficiently small (hence $\eta_0$ very small and possibly $T_0$ very large) so that the above inequality holds for $\eps_1$ as in Lemma \ref{LmVelaz}. 
		
		
		\noindent{\bf Step 3. Improvement of $C^1$-norm.} Now we extend the $H^1$-norm of $u$ also to a larger scale. In fact,
		\begin{align*}
			\|u\|_{H^1(B_{(1+\alpha)Kt^{\kappa}})}^2
			=&
			\|u\|_{H^1(B_{\br(t)})}^2
			+
			\int_{B_{(1+\alpha)Kt^{\kappa}}\backslash B_{\br(t)}}
			(|u|^2+|\nabla u|^2) e^{-\frac{|x|^2}{4}}dx
			\\
			\leq &
			\|u\|_{H^1(B_{\br(t)})}^2
			+
			\eps_1^2 \cF(\cC_{n,k}\backslash B_{\br(t)})
			\\
			\leq &2t^{-2\xi}.
		\end{align*}
		Here we used the fact that 
		\begin{align*}
			\cF(\cC_{n,k}\backslash B_{\br(t)})
			= \sum_{j=0}^\infty \cF(\cC_{n,k}\cap(B_{\br(t)+j+1}\backslash B_{\br(t)+j}))
			\leq 
			C\sum_{j=0}^\infty (\br(t)+j+1)^n e^{-\frac{(\br(t)+j)^2}{4}},
		\end{align*}
		and when $t$ is sufficiently large, this infinite sum is bounded by $Ce^{-(1-\epsilon)\frac{\br(t)^{2}}{4}}$. Then if we fix sufficiently small $\eps_2<\bar\eps$, where $\bar\eps$ is given in Proposition \ref{PropVelazquez}, we have $\|u\|_{H^1(B_{(1+\alpha)Kt^{\kappa}})}\leq 2t^{-\xi}$. Then we can apply Proposition \ref{PropVelazquez} to show that for $t\in[T_0+2, T_0+4]$,
		\begin{equation}
			\|\chi_{Kt^\kappa} u(\cdot,t)\|_{C^1(B_{K t^{\kappa}})}\leq C\eps_2 t^{2\kappa-\xi}\leq C\eps_2 {T_0}^{2\kappa-\xi}.
		\end{equation}
		In particular, if initially $T_0$ is chosen sufficiently large, we have $\|u(\cdot,t)\|_{C^1(B_{K t^{\kappa}})}\leq \min\{\eps,\eta_0/2\}.$
		
		\noindent{\bf Step 4. Iteration.} Now we repeat Step 2 with $T_0$ replacing by $T_0+2$. Note that after step 3, we have improved the estimate $\|u(\cdot,t)\|_{C^1(B_{Kt^\kappa})}<\eta_0$ to $t\in[T_0+2,T_0+4]$, thus Step 2 is applicable. Once we have Step 2, we can apply Step 3.
		
		Then keep repeating Step 2 and Step 3, we can extend the estimate $\|u(\cdot,t)\|_{C^1(B_{K t^{\kappa}})}< \eps$ to all $t\geq T_0$. Finally, the $C^2$-estimate of $u$ is a direct consequence of the Pseudolocality estimate again.
	\end{proof}
	
	From the proof, we can get an even better $C^2$ bound for $u$.
	
	\begin{corollary}\label{cor_C2_bound_for_kappa<1/2}
		Suppose the assumptions in Proposition \ref{thm:mainthmfor sqrtt}. Further assume $\xi=3/8$ and $\kappa\in(0,1/2)$ such that $2\kappa<\xi$. Then for any $\eps'\in(0,1)$, there exist $C>0$ and a possibly smaller $K>0$ such that $\|u(\cdot,t)\|_{C^2(B_{ K t^\kappa})}\leq Ct^{2\kappa-\xi}$.
	\end{corollary}
	
	\begin{proof}
		From the proof of Proposition \ref{thm:mainthmfor sqrtt}, we have $\|u(\cdot,t)\|_{C^1(B_{ K t^\kappa})}\leq Ct^{2\kappa-\xi}$. Therefore, it suffices to improve the estimate to $C^2$. First, we have already proved that $\|u\|_{C^2(B_{ K t^\kappa})}<\eps_0$, and hence, by the mean value theorem, $u$ satisfies a uniformly parabolic PDE in $B_{4n}$, and then we can apply the interior parabolic Schauder estimate to show that $\|u(\cdot,t)\|_{C^{2,\alpha}(B_{3n})}\leq \sup_{s\in[t-1,t]}\|u(\cdot,s)\|_{C^0(B_{4n})}\leq C(t-1)^{2\kappa-\xi}\leq \tilde C t^{2\kappa-\xi}$, if $t$ is sufficiently large. Next, we observe that for any $\mathbf{y}\in\{0\}\times\R^k$, we have $M_t-e^{(t-t_0)/2}\mathbf{y}$ is also a RMCF. In particular, this shows that $\|u(\cdot,t)\|_{C^{2,\alpha}(B_{3n}(e^2\mathbf{y}))}\leq \tilde C t^{2\kappa-\xi}$, whenever $B_{4n}(\mathbf{y})\subset B_{K(t-1)^{\kappa}}$. This implies that $\|u(\cdot,t)\|_{C^1(B_{ K' t^\kappa})}\leq C't^{2\kappa-\xi}$, for some constant $C'$ larger than $C$ in the proof of Proposition \ref{thm:mainthmfor sqrtt} and $K'$ slightly smaller than $K$ in the proof of Proposition \ref{thm:mainthmfor sqrtt}.
	\end{proof}
	
	
	\section{The invariant cones and $H^1$-normal form}\label{SDynamics}
	In Section \ref{S:Extending_t_kappa}, we have extended the graphical scale to $Kt^\kappa$ for some small $K>0$ and $\kappa=3/8$. However, this graphical radius is still not sufficient for us to derive the geometric consequence, such as the isolatedness of nondegenerate singularities, which requires a graphical radius of order $t^{1/2}$.  Note that by Proposition \ref{PropVelazquez}, the estimate of the graphical radius is mainly constrained by the $L^2$-estimate of the graphical $u$ provided in Proposition \ref{PropR}. To further extend the graphical scale, we need to improve the $L^2$-estimate of $u$. In this section, we prove that the linear equation $\partial_t u=L_\Sigma u$ controls the dynamics of the rescaled mean curvature flow, and prove the $H^1$-normal form theorem. 
	
	Recall that $\|\cdot\|$ without a subscript denotes the Gaussian weighted $H^1$-norm of a function over $\cC_{n,k}$. We consider the region of radius $O(t^\kappa)$ for some $\kappa\in(0,1/2)$, which is the improved graphical radius from Section \ref{S:Extending_t_kappa}. Suppose $u(\cdot,t)$ is the graph function of the rescaled mean curvature flow over $\cC_{n,k}\cap B_{Ct^\kappa}$. We let $\bar \eps(t):=t^{2\kappa-\xi}$, where we fix a constant $\xi\in(2\kappa,1)$, and by Corollary \ref{cor_C2_bound_for_kappa<1/2}, $\|u\|_{C^2(\cC_{n,k}\cap B_{\br(t)})}\leq C\bar \eps(t)$.
	
	In the following, we define $2\vartheta:=2\kappa-\xi$, and we choose $\br(t)=t^{\vartheta/8}$, for the technical purpose in Section \ref{SS:BC-rotation}.
	
	
	\subsection{Brendle-Choi's argument on rotations}
	\label{SS:BC-rotation}
	The natural function space to study these equations is $H^1(\cC_{n,k})$, which admits a natural direct sum decomposition into eigenspaces of the $L$-operator. Lying in the kernel of $L$, those $\theta_\alpha y_j\in\mathfrak{so}(n+1)$ represent infinitesimal rotations. These terms do not influence the major decay behavior, see \cite{CM2}; on the other hand, they do show up in the nonlinear analysis. Following an idea of Brendle-Choi \cite{BC1, BC2}\footnote{While \cite{BC1, BC2} studied ancient RMCF, as the RMCF equations are the same, their arguments can be simply adapted to our setting by reversing the time range.}, we modulo these modes by considering the rotated rescaled mean curvature flow $\tilde M_t=S_t M_t,\ S_t\in \mathrm{SO}(n+1)$. Let us write $\tilde M_t$ over $\mathcal{C}_{n,k}\cap B_{\br(t)}$ as the graph of a function $\tilde u$, where the rotation $S_t$ is chosen to make sure that $\chi_{\br(t)} \tilde u$ does not have Fourier modes corresponding to rotations. The existence of $S_t$ can be proved as in \cite[Proposition 2.4]{BC2}.
	
	Explicitly, suppose $\br(t)=Ct^{\vartheta/8}$, serving as the role of $\rho$ in \cite[Proposition 2.4]{BC2}, we require the identity $\int \tilde u\chi_{\br(t)} \langle Ax,\mathbf{n}(x)\rangle e^{-\frac{|x|^2}{4}}dx=0$ to hold for all $A\in \mathfrak{so}(n+1)$ at every moment $t,$ where $\mathbf n(x)$ is the unit outer normal of the cylinder $\mathcal C_{n,k}$ at the point $x.$ We have that $\tilde u$ solves the equation $\partial_t \tilde u=L\tilde u+\mathcal Q(\tilde u)+\mathcal A$, where $\mathcal A=\langle S'_t S_t^{-1}x,\mathbf n(x)\rangle$ (see \cite{BC2}, Proposition 2.4) and the evolution of $\overline{v}=\tilde u\chi_{\br(t)}$ with cutoff satisfies the equation
	\begin{equation}\label{eq:bar equation}
		\pr_t \overline{v}=L\overline{v}+\overline{\cQ}(J^2\overline{v})+\chi_{\br(t)}\mathcal A,
	\end{equation}
	where $\|\bar v\|_{C^2(\br(t))}\leq t^{-\vartheta}$ (see \cite{BC2}, Proposition 2.4), $\overline{\cQ}(J^2\overline{v})=O(t^{-\vartheta/8})(|\bar v|+|\nabla \bar v|+|\cA(t)|)$ (see \cite[Lemma 2.4, 2.5]{BC1}), and the RHS is zero when projected to the eigenmodes corresponding to rotations, since $\bar v$ does not have those Fourier modes by the choice of $S_t$. We have the following lemma.
	\begin{lemma}\label{LmA}
		For sufficiently large $t$, $	\|\chi_{\br(t)}\cA\|_{L^2}\leq 2\|\overline{\cQ}\|_{L^2}$.
	\end{lemma}
	
	\begin{proof}[Proof of Lemma \ref{LmA}] 
		From the definition of $\cA$, we have 
		$\int (\bar{\cQ}(J^2u)-\chi\cA(t))\cA(t)e^{-\frac{|x|^2}{4}}\,dx=0$
		and by Plancherel and Cauchy-Schwarz inequality,
		\begin{equation}\label{EqRotation}
			\|\chi_{\br(t)}\cA\|_{L^2}^2\leq \left|\int \chi_{\br(t)}\cA\cA e^{-\frac{|x|^2}{4}} dx\right|\leq \|\bar{\cQ}(J^2u)\|_{L^2}\|\cA\|_{L^2}. 	
		\end{equation}
		As $\cA$ is a concrete linear function in $x$, when $t$ is sufficiently large, $\|\cA\|_{L^2}\leq 2\|\chi_{\br(t)}\cA\|_{L^2}$. 
	\end{proof}
	
	
	As a consequence, we may also write the equation of $\bar v$ as 
	\begin{equation}\label{eq:bar equation after reduction}
		\pr_t \overline{v}=L\overline{v}+\overline{\cB}(J^2\overline{v}),
	\end{equation}
	with $\|\overline{\cB}\|_{L^2}\leq O(t^{-\vartheta/8})\|\bar v\|_{H^1}$. In the rest of this section, we remove all the overline notations for notational simplicity.
	
	
	
	
	
	
	\subsection{The cone theorem}
	The invariant cone theorem is an important tool in dynamical systems to study the stable/unstable manifolds. Let $E=H^1(\cC_{n,k})/\text{span}\{\theta_\alpha y_j\}$ (recall we have modulo these rotations in Section \ref{SS:BC-rotation}) and we introduce the direct sum decomposition $E=E^+\oplus E^0\oplus E^-$ where $E^+$ (respectively $E^0$ and $E^-$) is spanned by eigenfunctions of $L$ with positive (respectively $0$ and negative) eigenvalues. We denote by $\Pi_+,\ \Pi_-,\ \Pi_0$ the $L^2$-projections to the spaces $E^+, E^-, E^0 $ respectively. 
	
	We introduce a double cone construction that is used to suppress the $E^+$-components and manifest the $E^0$-component. Let $\alpha>0$ be a positive number. We introduce two cones $\cK_{\geq 0}$ and $\cK_{0}$ as follows
	$$\cK_{\geq 0}(\alpha):=\left\{u=(u_+,u_0,u_-)\in E^+\oplus E^0\oplus E^-\ |\ \|u_++u_0\|\geq \alpha\|u_-\|\right\}$$
	is a $\alpha$-cone around $E^+\oplus E^0$
	and 
	$$\cK_{0}(\alpha):=\left\{u=(u_+,u_0,u_-)\in E^+\oplus E^0\oplus E^-\ |\ \|u_0\|\geq 
	\alpha\|u_++u_-\|\right\}$$
	is a $\alpha$-cone around $ E^0$. Both are narrower when $\alpha$ is large. 
	
	We shall also need to compare the difference between the perturbed rescaled mean curvature flow and the unperturbed one under evolution. For this purpose, we introduce the following setting. Let $u_1,u_2$ be two graphical rescaled mean curvature flows over $\cC_{n,k}$, and $\|u_i(\cdot,t)\|_{C^{2,\alpha}}\leq \eps_0$ inside $\cC_{n,k}\cap B_{\br(t)}$. Then we write $v=\chi(u_1-u_2)$, where $\chi$ is a smooth cutoff function that is 0 outside the ball $B_{\br(t)}$ and is $1$ inside the ball $B_{\br(t)-1}$, and we can calculate
	\begin{equation}\label{EqDifference}
		\pr_t v=L v+\delta\cB. 
	\end{equation}
	where $\delta\cB=\delta\cB(J^2u_1,J^2u_2)=\cB(J^2u_1)-\cB(J^2u_2)$. In $\delta\cB$, we replace $u_i$ by $\chi u_i$ creating an error supported on the annulus $\mathbb A_{\br(t)}$ denoted by $\mathcal E(J^2u_1,J^2u_2). $ Thus we can write 
	\begin{equation}\label{EqPE}
		\delta\cB =P v+\mathcal E(J^2u_1,J^2u_2), \quad P = \int_0^1 D \cB(sJ^2(\chi u_1)+(1-s)J^2(\chi u_2))ds.		
	\end{equation}
	
	In this paper, we only consider the case that $u_1=u$ is a given rescaled mean curvature flow and $u_2=0$ is the rescaled mean curvature flow of the shrinking cylinder. Nevertheless, the study of the difference of the graphs of two rescaled mean curvature flows is nowhere more complicated, and it is used in our forthcoming work. Therefore, we work with this general setting in this paper.
	
	Throughout the rest of this section, we consider discrete-time $n\in\mathbb Z_{+}$ for presentation simplicity, although the proof for continuous time is almost verbatim. We also remind the readers that $\mathbb A_{\mathbf r(t)}$ is the annulus region $\cC_{n,k}\cap(B_{\br(t)}\backslash B_{\br(t)-1})$.
	
	Let us state the main approximation Proposition, where the proof is in Section \ref{SS:proofofPropCone}. This Proposition shows that the linearized solution of the rescaled mean curvature flow equation can approximate the rescaled mean curvature flow equation nicely, whenever the cutoff does not bring in too much error.
	\begin{proposition}\label{PropApproximation} 
		Suppose $\eps\in(0,\eps_0)$, where $\eps_0$ is from Appendix \ref{appendix:Nonlinear estimate}. Let $v( t)=\chi(t)(u_1( t)-u_2(t)):\ \cC_{n,k}\cap B_{\br(t)}\to \R$ be as above a solution to \eqref{EqDifference} satisfying $\|v(t)\|_{C^2(B_{\br(t)})}\leq \eps$ and  $\eps\| v(m)\|\geq \|v(t)\|_{C^2(\mathbb A_{\mathbf r(t)})}\omega(t)^{1/2}$,  for all $t\in [m,{m+1}]$ for some  $m\in\mathbb Z_+$. 
		Then we have 
		$$	\|v({m+1})-e^{L}v(m)\|\leq C\eps\|v(m)\|. $$
		Here $e^Lv(m)$ means the heat semigroup generated by $L$ acting on $v(m)$ for time $1$.
	\end{proposition}
	The assumption $\eps\| v(m)\|\geq \|v(t)\|_{C^2(\mathbb A_{\mathbf r(t)})}\omega(t)^{1/2}$ in Proposition \ref{PropApproximation} requires that the cutoff does not bring in too much error. This assumption either holds for all sufficiently large $m$, or fails for all sufficiently large $m$. This fact is verified in the following lemma, which is proved in Section \ref{SS:ConeBoundary}.
	\begin{lemma}\label{LmContinuation}
		Suppose $\br(t)=Kt^{\kappa}$ for some $K>0$ such that there is no eigenvalue of $L$ in the interval $(-K^2/4,0)$ and $\kappa\in (0,1/2]$, $v(t)=\chi(t)(u_1( t)-u_2(t)):\ \cC_{n,k}\cap B_{\br(t)}\to \R$ be as above a solution to \eqref{EqDifference}. Then we have the following dichotomy:
		\begin{itemize}
			\item either $\|v(t)\|\geq \omega(t)^{1/2}$ for all sufficiently large $t$,
			\item or $\|v(t)\|\leq \omega(t)^{1/2}$ for all sufficiently large $t$.
		\end{itemize} 
		Moreover, for any $\alpha>0$, there exists $\eps_1$ such that if $\|v(t)\|_{C^2(B_{\br(t)})}<\eps_1$, and $v(t)\not \in \cK_{\geq 0}(\alpha)$ for all sufficiently large $t$, then the second case of the dichotomy holds.
	\end{lemma}
	
	The following cone theorem follows immediately from Proposition \ref{PropApproximation} and Lemma \ref{LmContinuation}.  Briefly speaking, the cone theorem says that once the rescaled mean curvature flow of the graph enters a cone, it will stay inside the improved cone. 
	
	\begin{theorem}[Cone theorem]\label{ThmCone}
		For $i=1,2$, let $M_t^i$ be the rescaled mean curvature flow and $u_i:\ \Sigma^k\cap B_{\br(t)}\to \R$ be the graphical function of $M^i_t\cap B_{\br(t)}$ with $\|u_i(t)\|\to 0$ as $t\to\infty$, where the graphical function $\br(t)$ is as in the last lemma. Let $v=\chi(t)(u_1(t)-u_2(t))$. Suppose the first dichotomy in Lemma \ref{LmContinuation} holds. Then there exist $\alpha_0>0$, $m\in\mathbb Z_+$ and $\eta>0$ such that for $\alpha>\alpha_0$, the following holds:
		\begin{enumerate}
			\item if $v(m)\in \cK_{\geq 0}(\alpha)$, then $v(m')\in \cK_{\geq 0}((1+\eta)\alpha)$ for all $m'\in\mathbb Z_+$, $m'\ge m$;
			\item if $v(m)\in \cK_{0}(\alpha)$, then $v(m')\in \cK_{0}((1+\eta)\alpha)$ for all $m'\in\mathbb Z_+$, $m'\geq m. $
		\end{enumerate}
	\end{theorem}
	
	
	\begin{proof}
		We shall apply Proposition \ref{PropApproximation} with $u_2=0$ and $u_1=u$, the graphical function of the rescaled mean curvature flow $M_t$. 
		
		If the first dichotomy in Lemma \ref{LmContinuation} holds, the assumption of Proposition \ref{PropApproximation} holds. Thus, we can apply Proposition \ref{PropApproximation}, which means that the solution to the nonlinear equation \eqref{EqDifference} is well approximated by that of the linear equation $\partial_tu=Lu$. 
		
		For item (1), we denote $v=v_1+v_2$ with $v_1\in E^+\oplus E^0$ and $v_2\in E^-$. Given $\alpha>0$, suppose we have $v(m)\in \cK_{\geq 0}(\alpha),$ i.e. $\frac{\|v_1(m)\|}{\|v_2(m)\|}\geq \alpha$. Then we have by Proposition \ref{PropApproximation}
		$$\|v_1(m+1)\|\geq \|v_1(m)\|-C\eps(1+\alpha^{-2})^{1/2}\|v_1(m)\|,$$
		$$\|v_2(m+1)\|\leq e^{-\lambda}\| v_2(m)\|+C\eps (1+\alpha^{-2})^{1/2}\|v_1(m)\|,$$
		where $\lambda$ is the smallest positive eigenvalue of $L$. Taking the quotient, we get 
		\begin{equation}\label{Eqalpha}
			\alpha':=\frac{\|v_1(m+1)\|}{\|v_2(m+1)\|}\geq \frac{1-C\eps(1+\alpha^{-2})^{1/2}}{e^{-\lambda}+C\eps(1+\alpha^{-2})^{1/2}\alpha}\alpha,    
		\end{equation}
		which is greater than $(1+\eta)\alpha$ if $\alpha\geq \alpha_0$ is larger than a big multiple of $\eps_0>\eps$, and $\eta>0$ is chosen accordingly. Thus, we get $v(m+1)\in \cK_{\geq 0}((1+\eta)\alpha)\subset \cK_{\geq 0}(\alpha)$. 
		
		For item (2), we introduce $\cK_{>0}(\alpha)$ in a similar manner to $\cK_{\geq 0}(\alpha)$, then $u(m)\in \cK_{> 0}(\alpha)$ implies $u(m+1)\in \cK_{> 0}(\alpha)$, which implies $\|u_+(m)\|$ grows exponentially due to the presence of positive eigenvalue of $L$. By assumption, we have $\|u_+(m)\|\to 0$, then we get $u(m)\notin \cK_{> 0}(\alpha)$ for all $\alpha>0$ and all $m$ large. From item (1) and the fact that $\cK_0(\alpha)\subset \cK_{\geq 0}(\alpha)$, we get that $u(m)\in \cK_0(\alpha)$ implies $u(m+1)\in \cK_0((1+\eta)\alpha)$, which proves item (2). 
		
	\end{proof}
	\subsection{Proof of Proposition \ref{PropApproximation}}\label{SS:proofofPropCone}
	
	In this subsection, we prove Proposition \ref{PropApproximation}. 
	
	\begin{proof}[Proof of Proposition \ref{PropApproximation}]
		
		Let $w(\cdot,t)$ for $t\in[m,m+1]$ be a solution to the linearized rescaled mean curvature flow equation, i.e. $\pr_t w=Lw$ with the initial condition $w(m)=\chi(m)v(m)$ for some large $m$. We next estimate the evolution of $v-w$. We first compute the time derivative of $\|v-w\|^2$, which is
		\begin{equation*}
			\begin{split}
				\pr_t\int (v-w)^2 e^{-\frac{|x|^2}{4}}
				=&
				2\int (v-w)(L(v-w)+\dt\cB) e^{-\frac{|x|^2}{4}}
				\\
				=&
				-2\int |\nabla(v-w)|^2e^{-\frac{|x|^2}{4}}
				+2\int (v-w)^2e^{-\frac{|x|^2}{4}}
				+2 \int (v-w)\dt\cB e^{-\frac{|x|^2}{4}}.
			\end{split}
		\end{equation*}
		Similarly, we compute the time derivative of $\|\nabla(v-w)\|^2$:
		\begin{equation}\label{EqEnergy}
			\begin{split}
				&\pr_t\int |\nabla (v-w)|^2 e^{-\frac{|x|^2}{4}}
				=
				2\int \nabla(v-w)\cdot\nabla(L(v-w)+\dt\cB )e^{-\frac{|x|^2}{4}}
				\\
				=&
				-2\int |\cL(v-w)|^2e^{-\frac{|x|^2}{4}}
				+2\int |\nabla(v-w)|^2e^{-\frac{|x|^2}{4}}
				+2 \int (\cL(v-w))\dt\cB e^{-\frac{|x|^2}{4}}.
			\end{split}
		\end{equation}
		We bound $2\int (\cL(v-w))\dt\cB e^{-\frac{|x|^2}{4}}$ by $\frac{1}{10}\|\cL(v-w)\|_{L^2}^2+10\|\dt \cB\|^2_{L^2}$. 
		Since we have $\dt\cB=Pv+\mathcal E$, we get $\|\dt \cB\|^2_{L^2}\leq 2\|Pv\|^2_{L^2}+2\|\mathcal E\|^2_{L^2}$. 	
		Notice that $\dt\cB$ is defined in equation \eqref{EqDifference} independent of $w$, and $\dt \mathcal B=Pv+\mathcal E$, where $\mathcal E$ is only supported on $\A_{\br(t)}$. As $\|v\|_{C^2(B_{\br(t)})}\leq \eps$ by the definition of graphical radius, it is straightforward that 
		\begin{equation}\label{eq:cB L_2}
			\|\dt\cB\|_{L^2}^2\leq \eps^2\|J^2 v\|_{L^2}^2+\omega(t) \|v\|^2_{C^2(\mathbb A_{\br(t))}},
		\end{equation}
		where we used Proposition \ref{PropR}.
		
		The term $\|J^2v\|_{L^2}^2$ is bounded by a multiple of $\|\mathcal L v\|_{L^2}^2+\|\nabla v\|_{L^2}^2$, and $\|\mathcal L v\|_{L^2}^2$ is further bounded by $\|\nabla^2 v\|^2+\|\nabla v\|^2$ by integrating the Bochner formula (details can be found in \cite[Section 5]{CM3} and \cite[Proposition 4.6]{SX3})
		\[
		\frac{1}{2}\cL|\nabla v|^2=|\nabla^2 v|^2+\langle \nabla\cL v,\nabla v\rangle +(\Ric+\Hess_{|x|^2/4})(\nabla v,\nabla v).
		\]
		In summary, we have
		\begin{equation}\label{eq:LinearApproxWithoutAnnBound}
			\pr_t\|v-w\|^2\leq C\|v-w\|^2+\eps^2\|v\|_{H^2}^2+\|v\|^2_{C^2(\mathbb A_{\br(t)})} \omega(t).
		\end{equation}
		
		
		Moreover, by assumption, we have $\|v\|^2_{C^2(\mathbb A_{\mathbf r(t)})}\omega(t)\leq \eps^2\|v(m)\|^2$. 
		In conclusion, we have
		\begin{equation}\label{EqDiffH1}
			\pr_t \|v-w\|^2
			\leq C\|v-w\|^2+\eps^2\|v(t)\|_{H^2}^2+\eps_0^2\|v(n)\|^2
		\end{equation}
		with initial condition $\|v(n)-w(n)\|^2=0$. 
		Note that $v(m+1)-w(m+1)=v(m+1)-e^Lv(m)$, by Gronwall inequality,
		$$\|v(m+1)-e^Lv(m)\|^2\leq e^{C}\eps_0^2(\int_0^1\|v(m+s)\|_{H^2}^2ds+\|v(m)\|^2).$$
		It remains to estimate $\int_0^1\|v(m+s)\|_{H^2}^2ds$. Repeating the above calculations with $w=0$, we get \eqref{EqEnergy} with $w=0$ and the same estimate \eqref{eq:cB L_2}. Substituting \eqref{eq:cB L_2} into \eqref{EqEnergy} with $w=0$ and integrating over time $1$, we get 
		\begin{equation}\label{eq:LinearApproxH2}
			\int_0^1\|v(m+s)\|_{H^2}^2ds\leq C(\|v(m)\|^2+\sup_{t\in[m,m+1]}\omega(t)\|v(\cdot,t)\|^2_{C^2(\mathbb A_{\br(t)})}).
		\end{equation}
		Substituting the last estimate to the above estimate of $\|v(m+1)-e^Lv(m)\|^2$, and using $\|v\|^2_{C^2(\mathbb A_{\mathbf r(t)})}\omega(t)\leq \eps^2\|v(m)\|^2$ derives the desired estimate in the statement. 
		
	\end{proof}
	\subsection{Control the boundary term}\label{SS:ConeBoundary}
	In this section, we prove Lemma \ref{LmContinuation}. 
	\begin{proof}[Proof of Lemma \ref{LmContinuation}]
		
		
		For any $t>0$, we follow the proof of Proposition \ref{PropApproximation} all the way until equation \eqref{eq:LinearApproxWithoutAnnBound}, and use \eqref{eq:LinearApproxH2} to estimate $\|v\|_{H^1}^2$, where neither needs the assumption of $\omega(s)^{1/2}\|v(s)\|_{C^2(\mathbb A_{\mathbf r(s)})}\leq \eps\|v(t)\|$ for $s\in[t,t+1]$. Let $x=\|v_-(t)\|^2,\ y=\|v_0(t)\|^2, z=\|v_+(t)\|^2$, then we get 
		\begin{equation}\label{Eq3Eqs}
			\begin{cases}
				\dot x&\leq -cx+\eps(x+y+z)+O(\eps_1\omega(t)),\\
				|\dot y|&\leq \eps(x+y+z)+O(\eps_1\omega(t)),\\
				\dot z&\geq z-\eps(x+y+z)+O(\eps_1\omega(t)),\\
			\end{cases}
		\end{equation}
		where $-c=2\max\{-1/(n-k),-1/2\}<0$.  
		
		\medskip
		
		\noindent{\bf Claim 1:} {\it If we have 
			\begin{enumerate}
				\item for all sufficiently large time, $|x(t)|\geq \alpha^{-1} (|y(t)|+|z(t)|)$ for some constant $\alpha>0$ with $\eps\al\ll c$ (i.e. $v(t)\notin \cK_{\geq 0}(\alpha)$ for all $t$), and 
				\item $|x(t)|+|y(t)|+|z(t)|\to0$,
			\end{enumerate}
			then there is a constant $C$ such that  we have $|x(t)|+|y(t)|+|z(t)|\leq C\eps_1\omega(t)$.}
		\medskip
		
		Note that with the claim, we conclude the ``moreover'' part of the lemma.
		
		\begin{proof}[Proof of Claim 1] The condition (1) allows us to consider only the $\dot x$-equation, and without loss of generality, we remove the term $\eps (x+y+z)$ by redefining a smaller $c$. Let $C_1$ be a  bound of the $O(\eps_1\omega(t))$ in the above inequalities and $C$ be a constant in the statement satisfying $C\gg C_1/c$, and suppose the conclusion fails at some time $t_1$, then the $\dot x$-equation gives $\dot x\leq -(c-\frac{C_1}{C}\eps_1)x$, whose solution decays faster than $\omega(t)$ for $c>\frac{C_1}{C}\eps_1$. Then we can find some $t_2$ such that $x(t_2)\leq C\eps_1\omega(t_2)$ and $x$ continues to decay exponentially for all future time or until some $t_3$ with $cx(t_3)\approx C_1\eps_1\omega(t_3)$ when $x$ no longer decays exponentially. Then either $x(t)$ continues to satisfy $\leq \frac{C_1}{c}\eps_1\omega(t)$, or it lies in the interval $(\frac{C_1}{c}\eps_1\omega(t),C\eps_1\omega(t))$. In both cases, we have proved the Claim. 
		\end{proof}
		
		Next, suppose the latter case in the dichotomy does not hold. Then we can find a sequence of times $t_j\to\infty$ such that $\omega(t_j)^{1/2}\leq \|v(t_j)\|$. By the Claim 1, for any $\alpha_0>0$, $|x(t)|\geq \alpha_0^{-1}\cdot (|y(t)|+|z(t)|)$ cannot hold for all time. Therefore, there exists a sequence $s_j\to\infty$ such that $v(s_j)\in\cK_{\geq 0}(\alpha_0)$. Furthermore, we have the following:
		
		\medskip
		
		\noindent{\bf Claim 2:} {\it We can choose some large $j$ with $s_j=t_j$ such that we have 
			both $\omega(t_j)^{1/2}\leq \|v(t_j)\|$ and $v(t_j)\in\cK_{\geq 0}(\alpha_0)$, where $\alpha_0\gg1$ is a large constant, to be determined later.} 
		
		\begin{proof}[Proof of Claim 2]
			Suppose this is impossible, then let $[t_j,t_j']$ be the sequence of maximal intervals on which we have  $\omega(t)^{1/2}\leq \|v(t)\|$ and suppose that on a subinterval $[a,b]$ of $[t_j',t_{j+1}]$, we have $v(t)\notin\cK_{\geq 0}(\alpha_0)$, in addition to $\omega(t)^{1/2}\geq \|v(t)\|$. We claim that we can repeat the ODE argument of Claim 1 to get  $\|v(t)\|<\omega(t)^{1/2}$ on the interval $t\in [a,b]$ and obtain a contradiction to the definition of $t_{j+1}$. Indeed, 
			\begin{itemize}
				\item (case 1) if $x(t)>\eps_1 \omega(t)$, then the $\dot x$ equation is dominated by the $-c x$ part, thus $x$ decays faster than  $\omega(t)$;
				\item (case 2) otherwise, at some time $t_0$, $x(t_0)<\eps_1 \omega(t_0)\ll  \omega(t_0)$, the $\dot x$ equation is dominated by $O(\eps_1 \omega(t))$ with initial condition $x(t_0)<\eps_1 \omega(t_0)$, thus we get $x(t)\ll  \omega(t)$ for a long time until case 1 occurs. 
			\end{itemize}
			In either case, we always have $\|v(t)\|<\omega(t)^{1/2}$ for $t\in [a,b]$.   Thus, the only way to get growth of $\|v(t)\|$ such that at time $t_{j+1}$ we have $\omega(t_{j+1})^{1/2}\leq \|v(t_{j+1})\|$, is to have $v(t_{j+1})\in\cK_{\geq 0}(\alpha_0)$, which is exactly the claim.
		\end{proof} 
		
		With Claim 2, we repeat the proof of Proposition \ref{PropApproximation} for a very short time interval $[t_j,t_j+\Delta]$, where $\Delta$ is chosen such that $\|v(t)\|_{C^2(\mathbb A_{\mathbf r(t)})}\omega(t)^{1/2}\leq \eps_1\|v(t_j)\|$ for $t\in [t_j,t_j+\Delta]$ by continuity.	Equation \eqref{EqDiffH1} remains true on the $\Delta$-interval. 
		
		Then Item (1) of the Cone Theorem \ref{ThmCone} holds, and the cone condition $v(t)\in\cK_{\geq 0}(\alpha_0)$ holds for all $t\in[t_j,t_j+\Delta]$. 
		Then \begin{enumerate}
			\item  either $\|v(t)\|$ grows if $v(t)\in \cK_{>0}(\alpha')$ for some $\alpha'>\alpha_0$, in which case, we continue to have $\omega(t)^{1/2}\leq \|v(t)\|$ since $\omega(t)^{1/2}$ decays; 
			\item or we have $v(t)\in \cK_0(\alpha_0)$, in which case $v_0:=\Pi_0v$ term dominates. We shall analyze this case in detail in the next subsection, but here we only need a crude bound as follows: by Lemma \ref{LmMatrix}, $\nu:=\|v_0\|$ satisfies the equation $\dot\nu=-\gamma\nu^2+(O(\alpha_0^{-1}+\eps_1))\nu^2+O(\eps_1 \omega(t))$, where $\gamma>0$ is a fixed constant. The assumption on the $\Delta$ interval gives $\nu^2>\omega(t)$. Thus we get $\dot\nu=(-\gamma+O(\alpha_0^{-1}+\eps_1))\nu^2$ with solution $\nu(t)\geq \frac{1}{\nu(t_j)^{-1}+\int_{t_j}^t(\gamma+O(\alpha^{-1}+\eps_1))ds}$. When $\alpha_0$ is chosen sufficiently large and $\eps_1$ is chosen sufficiently small, say $|O(\alpha^{-1}+\eps_1)|$ in the above expression is bounded by $\gamma/100$, it can be verified that $\frac{\nu(t)}{\omega(t)^{1/2}}>1$ for $t>t_j$ if we have $\nu(t_j)=\omega(t_j)^{1/2}$, i.e. $\nu(t)$ decays slower than $\omega(t)$. 
		\end{enumerate}
		In either case, we have proved that $\|v(t)\|\geq \omega(t)^{1/2}$ holds on $[t_j,t_j+\Delta]$, and we can continue the procedure to extend this interval further. This completes the proof. 
	\end{proof}
	
	
	
	\subsection{Behaviors of the neutral modes}
	To proceed, we have the following observation. Fourier modes in $E^-$ decay exponentially, and those in $E^+$ grow exponentially. The convergence of the rescaled mean curvature flow to the cylinder implies that exponential growth is impossible, which implies that the $E^+$ component should be small. Thus $E^0$ component dominates the graphical function. The normal form then follows by projecting the rescaled mean curvature flow equation to the $E^0$ component and analyzing the resulting ODE. 
	
	
	
	Throughout this subsection, we let $v$ be the graphical function over the part of the cylinder $\mathcal C_{n,k}\cap B_{\br(t)}$ with the radius $\br(t)=Kt^{\vartheta/8}$, after cut-off, as in Section \ref{SS:BC-rotation}. Recall that $v$ is perpendicular to $\text{span}\{\theta_\alpha y_j\}$ corresponding to rotations. 
	
	\begin{proposition}\label{LmCoefficient}
		Suppose the graphical radius satisfies $\br(t)=Kt^{\kappa}$ for some $K>0$ and $\kappa\in (0,1/2]$ as in Lemma \ref{LmContinuation}. We have the dichotomy:
		\begin{enumerate}
			\item either $\|v(t)\|^2\leq \omega(t)=e^{-\frac{\br(t)^2}{4}}$ for all sufficiently large $t$, in which case $\mathcal I$ in Theorem \ref{ThmNF-sqrtt} is empty;
			\item or $\|v(t)\|^2> \omega(t)$ for all sufficiently large $t$. Moreover, if $\kappa\in(0,1/2)$, there is a nonempty subset $\emptyset\neq \mathcal I\subset\{1,2\ldots, k\}$ such that, up to a rotation in $\R^k$-factor, the coefficients of $y_i^2-2$ in Theorem \ref{ThmNF-sqrtt} are explicitly given by $\frac{\varrho}{4t}+O(t^{-1-\vartheta})$ for $i\in \mathcal I$ and $O(t^{-1-\vartheta})$ for $i\notin \mathcal I.$
		\end{enumerate}
		
	\end{proposition}
	
	
	The idea is that after modulo $\text{span}\{\theta_\alpha y_j\}$, the remaining neutral modes are $h_2(y_i)$'s and $h_1(y_i)h_1(y_j)$'s. Recall that we have defined the normalized (in $\|\cdot\|_{L^2(\cC_{n,k})}$) eigenfunctions $$H_2(y_i)=c_0^{k-1}\cG_{n,k}^{-1/2}h_2(y_i),\quad H_{1,1}(y_i,y_j)=c_0^{k-2}\cG_{n,k}^{-1/2}h_1(y_i)h_1(y_j),$$ where $\cG_{n,k}$ is given in \eqref{EqcG}.
	Define 
	$
	\begin{cases}
		m_{ii}(t)&:=\langle v,H_2(y_i)\rangle
		=
		c_0^{k-1}\cG_{n,k}^{-1/2}\langle v,h_2(y_i)\rangle,\\
		m_{ij}(t)&:=\langle v,H_{1,1}(y_i,y_j)\rangle
		= c_0^{k-2}\cG_{n,k}^{-1/2}\langle v, h_1(y_i)h_2(y_j)\rangle.        
	\end{cases}$
	Then we get $$
	\begin{aligned}
		v=&\Pi_{0}v+\Pi_{\neq 0}v=\sum_i m_{ii}H_2(y_i)+\sum_{i< j}m_{ij}H_{1,1}(y_i,y_j)+\Pi_{\neq 0}v.
	\end{aligned}
	$$
	We also define $\bar m_{ii}(t)=a\langle v(t),h_2(y_i)\rangle$ with $a=\sqrt{2}c_0$ and $\bar m_{ij}(t)=\langle v(t), h_1(y_i)h_1(y_j)\rangle$. The reason that we put $a=\sqrt{2}c_0$ in $m_{ii}$ is to make sure that later, when we compute the derivatives of $\bar{m}_{ii}$ and $\bar m_{ij}$, we will get uniform expressions. The choice of $a$ will become evident in later proofs, in particular \eqref{Eqgamma}. Let $\mathrm{M}(t)$ be a symmetric $k\times k$ matrix, whose entries are given by $\bar{m}_{ii}$ and $\bar{m}_{ij}$. 
	
	
	
	\begin{lemma}\label{LmMatrix}
		Suppose on some time interval $I=[a,b]$, we have $v(t)\in \cK_0(\alpha)$ for some sufficiently large $\alpha>\alpha_0$, and $\|v\|_{C^2(B_{\br(t)})}\leq \bar\eps$. Then for $t\in [a+2,b]$, $\mathrm{M}(t)$ satisfies the following ODE:
		\begin{equation}\label{eq:ODE-1}
			\mathrm{M}'(t)=-\gamma \mathrm{M}^2(t) + O((\alpha^{-1}+\bar\eps(t)) |\mathrm{M}(t)|^2)+O(\bar\eps\omega(t)),
		\end{equation}
		where $\gamma=\frac{1}{\varrho}c_0^{2(k-1)}\cG_{n,k}^{-1}$ is a fixed constant. Moreover, the eigenvalues (may not be ordered from large to small) of $\mathrm{M}$ satisfy 
		$$\lambda_i'=-\gamma\lambda_i^2+O((\alpha^{-1}+\bar\eps(t))\sum_{i=1}^k \lambda^2_i)+O(\bar\eps(t)\omega(t)),\ i=1,2,\ldots,k.$$ 
	\end{lemma}
	
	\begin{proof} By Section \ref{SS:BC-rotation}, we may write the equation of $v$ as $\pr_t v=Lv+\cQ(J^2 v)+\chi_{\br}\cA(t)$. From Proposition \ref{PropQ}, we write the reminder $\cQ(J^2u)$ of the rescaled mean curvature flow equation as
		\[
		\cQ(J^2u)=-\frac{1}{2\varrho}(u^2+4u\Delta_\theta u+2|\nabla_\theta u|^2)+\mathcal C(J^2u),
		\]
		where $\mathcal C(J^2u)=O(\|u\|_{C^{2}}(|u|^2+|\nabla u|^2+|\nabla^2u|))$. 
		Then we take the time derivative of $\bar m_{ii}$ and $\bar m_{ij}$ to get
		\begin{equation}\label{eq:Aderivative}
			\begin{split}
				\bar m_{ii}'(t)
				&=
				a\int (Lv+\chi\cQ(v))h_2(y_i)e^{-\frac{|x|^2}{4}}+O\left(\bar\eps\omega(t)\right)
				\\
				&=
				a\int \left(-\frac{1}{2\varrho}(v^2+4v\Delta_\theta v+2|\nabla_\theta v|^2)+\mathcal C(J^2v)+\chi_{\br(t)}\cA(t)\right)h_2(y_i)e^{-\frac{|x|^2}{4}}
				+O\left(\bar\eps(t)\omega(t)\right),
				\\
				\bar m_{ij}'(t)
				&=
				\int \left(-\frac{1}{2\varrho}(v^2+4v\Delta_\theta v+2|\nabla_\theta v|^2)+\mathcal C(J^2v)
				+\chi_{\br(t)}\cA(t)\right)h_1(y_i)h_1(y_j)e^{-\frac{|x|^2}{4}}
				+O\left(\bar\eps(t)\omega(t)\right),
			\end{split}
		\end{equation}
		where the $O\left(\bar\eps(t)\omega(t)\right)$ error is created by the cutoff $\chi(t)$ since the graphical scale is $O(Kt^\kappa).$ 
		
		We next show that the RHS of \eqref{eq:Aderivative} is dominated by terms involving $v^2$. 
		On the RHS of \eqref{eq:Aderivative}, we substitute 
		$$v=\sum_i  a^{-1} \bar m_{ii}c_0^{2(k-1)}\cG_{n,k}^{-1}h_2(y_i)
		+\sum_{i< j}\bar m_{ij} c_0^{2(k-2)}\cG_{n,k}^{-1} h_1(y_i)h_1(y_j)+\Pi_{\neq 0}v.$$
		Then, 
		\begin{itemize}
			\item terms involving $4v\Delta_\theta v+2|\nabla_\theta v|^2$ depend only on $\Pi_{\neq 0}v$ since $\Pi_0v$ has no $\theta$-dependence. Hence it is estimated as $\|\Pi_{\neq 0}v\|^2\leq C\alpha^{-2}\|v\|^2$, because by assumption, we have $\|\Pi_{\neq 0}v\|\leq \alpha^{-1} \|v\|$;
			\item terms involving $\mathcal C(J^2v)$, after a Cauchy-Schwarz inequality, become $C\bar\eps\||v|+|\nabla v|\|_{L^4}^2$ (see the expression of $\mathcal C(J^2v)$ in the item (2) of Proposition \ref{PropQ}). By Lemma \ref{lem_OU_control_L4byL2}, it is bounded by $C\bar\eps(\|v(\cdot,t-2)\|^2_{H^1}+e^{-t^{2\vartheta}/100n})$. The term $C\bar\eps e^{-t^{2\vartheta}/100n}$ is bounded by $O(\bar\eps\omega(t))$. Then in the proof of Proposition \ref{PropApproximation}, without using the smallness of $\|v\|_{C^2}$ but using \eqref{eq:LinearApproxH2} directly, $C\bar\eps\|v(\cdot,t-2)\|^2_{H^1}\leq C\bar\eps\|v(\cdot,t)\|^2_{H^1}+\bar\eps \omega(t-2)\leq C\bar\eps\|v(\cdot,t)\|^2_{H^1}+C\bar\eps\omega(t)$. Then using the fact that $v\in \cK_{\alpha}$, this term is bounded by $C\bar\eps|\mathrm{M}(t)|^2+C\bar\eps \omega(t)$;
			\item terms involving $\chi_{\br(t)}\cA(t)$ is bounded by $O(\bar\eps(t)\omega(t))$. In fact, $\cA(t)$ is the linear combination of infinitesimal rotations, so it is orthogonal to $h_2(y_i)$ and $h_1(y_i)h_1(y_j)$. Thus, the integral involving $\chi_{\br(t)}\cA(t)$ can be estimated by $\|\cA(t)\|_{L^2}\omega(t)$, and by Lemma \ref{LmA}, it is bounded by $O(\|v\|_{C^2(B_{\br(t)})}\omega(t))=O(\bar\eps(t)\omega(t))$.
		\end{itemize}
		
		Thus, we have proved that the RHS of \eqref{eq:Aderivative} is dominated by terms involving $v^2$. Moreover, by assumption again, on the RHS of \eqref{eq:Aderivative}, we replace $v^2$ by $(\Pi_0 v)^2$ creating another error bounded by $\alpha^{-1}\|v\|^2$. For $\Pi_0v$, its Fourier expansion has only finitely many terms, then we can write it into the Fourier expansion and calculate $m'_{ii}$ and $m_{ij}'$ explicitly using Lemma \ref{LmTriple}. We have 
		\begin{equation}\label{EqNeutral}
			\begin{cases}
				\bar m'_{ii}&=-\gamma \sum_{\ell=1}^k \bar m^2_{i\ell}+O((\alpha^{-1}+\bar\eps)\|\mathrm{M}(t)\|^2)+O(\bar\eps\omega(t)),
				\\
				\bar m'_{ij}&=-\gamma \sum_{\ell=1}^k \bar m_{i\ell} \bar m_{\ell j}+O((\alpha^{-1}+\bar\eps)\|\mathrm{M}(t)\|^2)+O(\bar\eps\omega(t)).
			\end{cases}
		\end{equation}
		Using Lemma \ref{LmTriple}, we can calculate that 
		\begin{equation}\label{Eqgamma}
			\begin{aligned}
				\gamma
				=&\frac{A_{2,2,2}}{2\varrho \cG_{n,k}}(4\pi)^{\frac{k-1}{2}}
				c_0^{4(k-1)} a^{-1}
				=\frac{A_{2,1,1}}{2\varrho \cG_{n,k}} (4\pi)^{\frac{k-2}{2}}c_0^{4(k-2)}a=
				\frac{c_0^{2(k-1)}}{\varrho \cG_{n,k}}
				.
			\end{aligned}    
		\end{equation}
		We can also see why we need to choose $a=\sqrt{2}c_0$ from the calculation.
		
		Finally, the part of the statement on the eigenvalues follows from the classical perturbation theory for linear operator, saying that for symmetric matrices, the eigenvalues, may not be ordered from large to small, are differentiable with respect to given parameters if the matrices are differentiable; see \cite[Chapter Two, Section 6, Theorem 6.8]{Kato95_PertubLinearOp} and the proof of \cite[Lemma 3.1]{V2} in the semilinear PDE setting. 
	\end{proof}
	
	\begin{proof}[Proof of Proposition \ref{LmCoefficient}]
		By Lemma \ref{LmContinuation}, either we have the $H^1$-norm $\|v(t)\|$ decays faster than $\omega(t)^{1/2}$, or the assumption of the Lemma \ref{LmMatrix} is satisfied for all large time. Consider the equation for eigenvalues in the last lemma. We have
		\begin{enumerate}
			\item either $\lambda:=\max_i\{\lambda_i(t)\}<\omega(t)^{1/2}$ for all $t$ large enough; 
			\item or there is a time sequence $t_n\to\infty$ such that we have $\lambda(t_n)\geq \omega(t_n)^{1/2}$.
		\end{enumerate}
		
		
		In the former case, $\lambda(t)=o(t^{-1})$. In the latter case, if $\kappa\in(0,1/2)$, on the time interval $[t_n,t_n+\Delta]$ for some small $\Delta>0$, we can write the equation $\lambda'=-(\gamma+O(\alpha(t)^{-1}+\bar \eps(t)))\lambda^2$, where $\bar\eps(t)\leq Ct^{-\vartheta}$ and $\alpha(t)=Ce^{\ln(1+\eta)t}$, by Theorem \ref{ThmCone}. The eigenvalue attaining the maximum may switch from one to the other, but the equation for $\lambda$ remains of the same form under the assumption in item (2). In other words, the derivative $\lambda'$ may be discontinuous at the switching time, but the size of the discontinuity is absorbed in $O(\alpha^{-1}+\bar\eps).$
		We first claim that $\lambda>0$, since, otherwise, the solution to the last equation quickly blows up to $-\infty$ if the initial condition is negative. 
		Then we are in the same situation as the end of the proof of Lemma \ref{LmContinuation}. The equation can be integrated as $$\lambda(t)^{-1}-\lambda(t_n)^{-1}=\int_{t_n}^t(\gamma+O(\alpha^{-1}+\bar\eps)) ds,$$ which implies that $\lambda(t)\geq \omega(t)^{1/2}$ holds on the time interval $[t_n,t_n+\Delta]$. We can then repeat the argument to extend $\Delta$ to infinity and get $\lambda(t)/ \omega(t)^{1/2}\to\infty$ as $t\to\infty$. Restricted to the graphical radius $\br(t)$ of size $K t^{\vartheta/8}$, by Corollary \ref{cor_C2_bound_for_kappa<1/2}, we have that $\|v(t,\cdot)\|_{C^2(B_{\br(t)})}\leq \bar\eps(t)\to 0$. Thus, the $O(\bar\eps\omega(t))$ is negligible and we have $\lambda(t)\gamma t\to 1$ in the limit $t\to\infty$ since we have $\alpha(t)^{-1}\to0$ and $\bar\eps(t)\to0$. Therefore, we have proved that the set $\mathcal I\neq\emptyset$ in the latter case. 
		
		With the maximum eigenvalue known, we next study the subleading eigenvalue denoted by $\lambda_2$.  Then 
		we have (denoting $\Omega(t):=\max\{\omega(t),((\alpha(t)^{-1}+\bar\eps(t))\lambda^2(t))\}$)
		\begin{enumerate}
			\item either $\exists\ c>0$ such that for all $t$ sufficiently large we have $\lambda_2(t)^2<c\Omega(t)$, 
			\item or $\forall\ c>0$, there exists $t_n\to\infty $ such that $\lambda_2(t_n)^2\geq c\Omega(t_n)$.
		\end{enumerate}
		In the former case, we get $\lambda_2(t)=o(1/t)$. In the latter case, we choose $c$ small and get $$\lambda_2(t_n)\geq c\Omega(t_n),\ \mathrm{and}\ \lambda_2(t)\geq 0.5c\Omega(t)$$ for $t\in [t_n,t_n+\Delta]$ for some small $\Delta>0$. Then we get
		$\lambda_2'=-(\gamma+o(1))\lambda_2^2$ on the same time interval. Then the same reasoning as the above paragraph gives that we can extend $\Delta$ to $\infty$ and  $\lambda_2\gamma t\to1.$ The analysis for all other eigenvalues is the same, and we get $\lambda_i\gamma t\to 1$ or $0.$ Moreover, we get that in the expansion of $v$, the coefficient of $(y_i^2-2)=c_2^{-1}h_2(y_i)$ is given by $\frac{1}{\gamma t}a^{-1}c_0^{2(k-1)}\cG_{n,k}^{-1}c_2^{-1}=\frac{\varrho}{4t}$.
		
		Once we obtain the asymptotic expansion for the eigenvalues of the matrix $M(t)$, let $W(t)$ denote the eigenspace associated to the eigenvalues $\lambda_i$'s with $\lambda_i(t)\approx \frac{1}{\gamma t}$, and let $\Pi_{W(t)}$ denote the orthogonal projection to $W(t)$. Notice that then $\lambda_i$ satisfies the equation $\lambda_i'=-\gamma \lambda_i^2 + O(t^{-2-\vartheta})$, which further implies that $\lambda_i(t)=\frac{1}{\gamma t}+O(t^{-1-\vartheta})$ or $\lambda_i(t)=O(t^{-1-\vartheta})$, depending on whether $\lambda_i(t)\approx \frac{1}{\gamma t}$ or $0$. Then the classical perturbation theory of linear operators (see \cite[Lemma 3.6]{V2} and the reference therein) shows that $E:=\lim_{t\to\infty}\Pi_{W(t)}$ exists. In fact, if we define the resolvent $R(z,t)=(\mathrm{M}(t)-z)^{-1}$ in the complex plane, $\Pi_{W(t)}$ can be expressed as the contour integral of $R(z,t)$ around those $\lambda_i$'s with decay rate $\frac{1}{\gamma t}$. Then, if we further choose circles of radius $c/t$ around those eigenvalues and do the contour integral, following the proof of \cite[Lemma 3.6]{V2}, we see that $\left|\frac{d}{dt}\Pi_{W(t)}\right|\leq O(t^{-1-\vartheta})$, hence $\Pi_{W(t)}$ converges. Note that although our higher-order error $O(t^{-2-\vartheta})$ is not $O(t^{-3})$ as in \cite[Lemma 3.6]{V2}, but after multiplying by $t$, $O(t^{-1-\vartheta})$ is still integrable, which is suffices for the proof of \cite[Lemma 3.6]{V2}.
		
		Note that $E$ is the orthogonal projection on an $\ell$-dimensional space. Then up to a rotation, $E$ is a diagonal matrix with first $\ell$ entries $1$. Then, following the proof in \cite[Lemma 3.6]{V2} (where we replace the $o(1/t)$-errors in the proof of \cite{V2} by the explicit errors here), we get $\mathrm{M}(t)=-\frac{P_{W(t)}}{\gamma t}+O(t^{-2-2\vartheta})=-\frac{E}{\gamma t}+O(t^{-1-\vartheta})$. This gives the desired asymptotics. 
		
		Finally, if $\kappa=1/2$, we repeat the above discussion, but $\bar\eps$ in this case may not converge to $0$. However, if at the beginning we choose the bound $\eps$ for the $C^2$-norm very small (say, much smaller than $\gamma$) to get the graphical radius, the ODE for $\lambda$ becomes $-\gamma/10 \lambda^2\geq \lambda'\geq -\gamma/2 \lambda^2$, and we obtain that $\lambda(t)>c/t$. This is also enough to conclude that $\|v(t)\|^2>\omega(t)$ for all sufficiently large $t$.
	\end{proof}
	\subsection{The $H^1$-normal form, special case}\label{SNF}
	Let us first prove the $H^1$-normal form for a particularly chosen graphical graphical radius $O(K t^{\vartheta/8})$.
	\begin{proof}[Proof of Theorem \ref{ThmNF-sqrtt} on $H^1$-normal form, with graphical radius $O(Kt^{\kappa})$ and a particular $\vartheta$]
		We first consider the function $v$ after the Choi-Brendle rotation reduction. The cone theorem gives an $H^1$ normal form as in Theorem \ref{ThmNF-sqrtt} but with error $o(1/t)$ in $H^1$. Indeed, by assumption, we have $v(t)\to0$ in the $C^\infty_{loc}$-sense as $t\to\infty$, thus $\|v_+(t)\|\to0$ as $t\to\infty$. By item (2) of the Cone Theorem \ref{ThmCone}, we see that 
		\begin{enumerate}
			\item either $v(t)$ remains in the cone $\cK_0(\alpha)$ for all time sufficiently large, 
			\item or $v(t)$ does not enter the cone $\cK_0(\alpha)$ for all time sufficiently large. 
		\end{enumerate}
		In the former case, $\alpha$ grows exponentially fast, and we see that the $u_0$ part dominates. Thus, we have verified the assumption of Lemma \ref{LmMatrix} and Proposition \ref{LmCoefficient}, which gives the $H^1$ normal form as in Theorem \ref{ThmNF-sqrtt} with $O(t^{-1-\vartheta})$ error in $H^1$. In the latter case, Lemma \ref{LmContinuation} implies that $v$ decays exponentially fast.
		
		
		Now we consider the original function $u$. By Lemma \ref{LmA}, we have $\|\cA(t)\|_{L^2}\leq 2\|\chi_{\br(t)}\cA(t)\|_{L^2}\leq 4\|\cQ\|_{L^2}$, which is the quadratic error bounded by $\|v\|^2$. When $\|v\|=O(e^{-\br^2/4})$, we have $\|\cA(t)\|_{L^2}\leq O(e^{-\br^2/4})$, which implies that $\|S(t)-\id\|_{L^2}\leq O(e^{-\br^2/4})$.
		
		When $\|v\|=O(t^{-1})$, the proof is more involved, as the naive bound $\|S(t)-\id\|_{L^2}\leq O(t^{-1})$ is not enough for the accuracy of the $H^1$-normal form. What we do is to study the evolution of the whole $0$-eigenmodes, and obtain an ODE like \eqref{EqNeutral}, but including the coefficients of $\theta_i y_j$. In fact, suppose $R_{\alpha j}$ is the coefficients for $\theta_\alpha y_j$ in $u$, we can calculate, as in the proof of Lemma \ref{LmMatrix},
		\begin{equation}\label{eq_ODE_of_Rot}
			R_{\alpha j}'=O((\alpha^{-1}+\bar\eps)\|\mathrm{M}(t)\|^2)+O(\bar\eps\omega(t)).
		\end{equation}
		In fact, because the rotated generalized cylinders are also (static) RMCFs, we have $\cQ(cJ^2(R_{\alpha j}\theta_\alpha y_j))\equiv 0$, which implies that the contribution of the nonlinear term is only the cut-off $O(\bar\eps\omega(t))$. \eqref{eq_ODE_of_Rot} implies that, along the same line as the proof of Proposition \ref{LmCoefficient}, $R_{\alpha j}=O(t^{-1-\vartheta})$. This implies the $H^1$-normal form for $u$, with the error $O(t^{-1-\vartheta}).$ Combining all the ingredients above proves the statement.
		
	\end{proof}
	
	Once we obtain the $H^1$-normal form of $u$ over the graphical radius $O(Kt^{\vartheta/8})$, we can improve the error estimate to $O(t^{-1-\vartheta})$ for any $\vartheta\in(0,1)$ and some graphical radius $O(t^{\bar\kappa})$ for $\bar\kappa\in(0,1/2)$.
	
	\begin{proof}[Proof of Theorem \ref{ThmNF-sqrtt} on $H^1$-normal form, with any $\vartheta\in(0,1)$ and graphical radius $O(Kt^{\bar\kappa})$]
		Given $\vartheta\in(0,1)$, we choose $\bar\kappa\in(0,1/2)$ and $\bar\xi\in(0,1)$ so that $\bar\xi-2\bar\kappa=\vartheta$. Because the special case of the $H^1$-normal form we have obtained so far showing that $\|u\|_{H^1(Kt^{\kappa})}=O(t^{-1})$, we repeat the proof of Proposition \ref{thm:mainthmfor sqrtt} in Section \ref{S:Extending_t_kappa} and Corollary \ref{cor_C2_bound_for_kappa<1/2} to get that 
		$\|u\|_{C^2(Kt^{\bar\kappa})}=O(t^{2\bar\kappa-\bar\xi})=O(t^{-\vartheta})$. Then by repeating all the proof of this section so far, we obtain the $H^1$-normal form, with any given $\vartheta$ and graphical radius $O(Kt^{\bar\kappa})$.
	\end{proof}
	
	
	\begin{remark}\label{rmk_small_radius_better_H1}
		If one studies the region with radius $\sqrt{c\log t}$ with $c>4$, it seems plausible to obtain an $H^1$-normal form with error estimate $O(t^{-2}\log t)$.
	\end{remark}
	
	To obtain the $H^1$-normal form with the largest possible graphical radius $O(Kt^{1/2})$, we first need to extend the graphical radius to that size, and we will discuss that in the following sections.
	
	
	
	
	\subsection{Extend the graphical radius to $K_0\sqrt{t}$}\label{SS_Extend_Graphical_Radius_to_K_0t^1/2}
	With the $H^1$-normal form, in Proposition \ref{PropVelazquez} we can choose $\xi=1$ and $\kappa=1/2$. Thus, Proposition \ref{PropVelazquez} allows us to extend the graphical scale to  $K_0t^{1/2}$ for some small $K_0$. 
	
	\begin{proposition}\label{PropSmallSqrt}
		For for $\vartheta\in(0,1)$, there exist $K_0>0$ and $T>0$ such that we have the graphical scale $\mathbf{r}(t)\geq K_0t^{1/2}$ for $t\geq T$, and
		\[
		\left\|
		u(\theta,y,t)-\varrho\left(\sqrt{1+\frac{\sum_{i\in\cI}(y_j^2-2)}{2t}}-1\right)
		\right\|_{C^1(\cC_{n,k}\cap B_{K_0 t^{1/2}})}=O(t^{-\vartheta}).
		\]
	\end{proposition}
	
	The proof, again, uses the ``extension-improvement'' argument, and is similar to the proof of Proposition \ref{thm:mainthmfor sqrtt}. However, there is one major difference that we no longer have a desired decay estimate for the graph function $\|u(\cdot,t)\|_{L^2}$ directly. Instead, we have the desired decay estimate for the difference between the graph function and the normal form. Thus, there is an intermediate step to handle this difference.
	
	\begin{proof}[Proof of Proposition \ref{PropSmallSqrt}]
		The proof is divided into several steps:
		
		\noindent{\bf Step 1. Initiating the setup.} Let us first fix a very small $\eps$ to be determined. Then we choose $T_0>0$ sufficiently large with the following significance: 
		\begin{itemize}
			\item Proposition \ref{PropVelazquez} is applicable.
			\item for $t>T_0$, $M_t$ is a graph of a function $u(\cdot,t)$ inside $\cC_{n,k}\cap B_{\bar R}$, such that $\|u(\cdot,t)\|_{C^1(B_{\bar R})}\leq \eta_0$, where $\bar R>R_0$ and $\eta_0>0$ as in the Pseudolocality Theorem \ref{thm:pseudo-locality} such that it can be applied for $\eps$. We also assume $\eta_0$ is chosen very small such that $\bar\delta$ in Theorem \ref{thm:pseudo-locality} is much larger than $2$. 
			
			We choose $K_0$ such that $K_0(T_0+2)^{1/2}=\bar R$. In particular, this shows that $\|u(\cdot,t)\|_{C^1(B_{K_0t^{1/2}})}<\eta_0$ for $t\in[T_0,T_0+2]$. On the other hand, by choosing $T$ large, we will also choose $K_0$ such that $K_0^2\varrho/4<\min\{\eps,\eta_0/3\}$.
		\end{itemize}
		
		\noindent{\bf Step 2. Extend the graphical region.}
		Apply the Pseudolocality theorem, we have for $t\in[T_0+2,T_0+4]$, $u(\cdot,t)$ is a graph of function over $\cC_{n,k}\cap B_{(1+\alpha)K_0t^{1/2}}$, for some $\alpha>0$, such that $$\|\nabla^ku(\cdot,t)\|_{L^\infty(\cC_{n,k}\cap B_{(1+\alpha)K_0t^{1/2}})}\leq  e^{-(k-1)(t-t_0)/2}\eps<\eps_1, \quad k=0,1,2,3.$$ 
		Here we choose $\eps$ sufficiently small (hence $\eta_0$ very small and $T_0$ very large) so that the above inequality holds for $\eps_1$ in Lemma \ref{LmVelaz}. 
		
		
		\noindent{\bf Step 3. Improvement of $C^1$-norm.}
		Let $w(\cdot,t)=u(\cdot,t)-\sum_{i\in \mathcal I}\frac{\varrho }{4t}(y_i^2-2)$. By the $H^1$-normal form Theorem \ref{ThmNF-sqrtt} with graphical radius $O(K t^{\kappa})$ for some $\kappa\in(0,1/2)$, $\|w(\cdot,t)\|_{H^1(\cC_{n,k}\cap B_{K t^\kappa})}\leq Ct^{-2}$. We extend this $H^1$ bound to a larger scale. In fact,
		\begin{align*}
			\|w(\cdot,t)\|^2_{H^1(\cC_{n,k}\cap B_{(1+\alpha)K_0t^{1/2}})}
			= &
			\|w(\cdot,t)\|^2_{H^1(\cC_{n,k}\cap B_{K t^\kappa})}
			+
			\int_{\cC_{n,k}\cap (B_{(1+\alpha)K_0t^{1/2}}\backslash B_{K t^\kappa})}(|w|^2+|\nabla w|^2)e^{-\frac{|x|^2}{4}}dx
			\\
			\leq & 
			\|w(\cdot,t)\|^2_{H^1(\cC_{n,k}\cap B_{K t^\kappa})}
			+
			C\int_{\cC_{n,k}\backslash B_{K t^\kappa}} (\eps_1^2+|y|^4/t^2)e^{-\frac{|x|^2}{4}}dx
			.
		\end{align*}
		The last integral can be estimated by
		\[
		\sum_{j=0}^\infty (1+(K t^\kappa+j+1)^4)(K t^\kappa+j+1)e^{-\frac{(K t^\kappa+j)^2}{4}}
		\leq 
		Ce^{-\frac{(Kt^{\kappa})^2}{8}}
		\leq Ct^{-4},
		\]
		whenever $t$ is sufficiently large. Thus we have $\|w(\cdot,t)\|_{H^1(\cC_{n,k}\cap B_{(1+\alpha)K_0t^{1/2}})}\leq Ct^{-1-\vartheta}$, with possibly slightly larger $C$, but when $T_0$ is sufficiently large, it can be a fixed constant.
		
		Then we can apply Proposition \ref{PropVelazquez} to show that for $t\in[T_0+2, T_0+4]$,
		\begin{equation}
			\|\chi_{K_0t^{1/2}}w(\cdot,t)\|_{C^1(B_{K_0 t^{1/2}})}\leq C\eps_1 t^{-\vartheta}\leq C\eps_1 {T_0}^{-\vartheta}.
		\end{equation}
		In particular, if initially $T_0$ is chosen sufficiently large, we have $\|w(\cdot,t)\|_{C^1(B_{K_0 t^{1/2}})}< \eta_0/2.$ This also shows that $\|u(\cdot,t)\|_{C^1(B_{K_0 t^{1/2}})}< \eta_0$.
		
		
		\noindent{\bf Step 4. Iteration.} Now we repeat Step 2 with $T_0$ replacing by $T_0+2$. Note that after step 3, we have improved the estimate $\|u(\cdot,t)\|_{C^1(B_{K_0t^{1/2}})}<\eta_0$ to $t\in[T_0,T_0+4]$, thus Step 2 is applicable. Once we have Step 2, we can apply Step 3.
		
		Then by repeating Step 2 and Step 3, we can extend the estimate $\|u(\cdot,t)\|_{C^1(B_{K_0 t^{1/2}})}<\eps$ to all $t\geq T_0$. Finally, the $C^2$-estimate of $u$ is a direct consequence of the Pseudolocality estimate again.
	\end{proof}
	
	
	
	
	\section{The $C^1$ normal form and isolatedness of nondegenerate singularities}\label{SS:The global variational equation: scale up to}
	
	For the purpose of getting geometric properties such as the isolatedness of nondegenerate singularities, we need to upgrade the $H^1$-normal form to the $C^1$-normal form. Moreover, we can also prove that the graphical radius can be $K\sqrt{t}$ for any $K>0$. This larger graphical scale is not necessary in this paper, but we keep it for later purposes. 
	
	The main difficulty is as follows. In the equation of motion for $u$, we have $\pr_t u=Lu+\cQ(J^2u)$, where $\cQ$ is a quadratic error, (see Lemma \ref{PropQ}). However, in the radius $O(\sqrt{t})$, $u(\theta,y,t)\approx \frac{\varrho}{4t}\sum_{i\in\cI}h_2(y_i)$, is not small when $\cI\not=\emptyset$: when $|y|\approx K\sqrt t$, $u(\theta,y,t)\approx \frac{|\cI|\varrho}{4}K^2$. Therefore, in the $C^0$ norm, the nonlinear term $\cQ(J^2u)$ is at least $O(K^4)$, which leads to cumulative errors to the estimate of $u$ as $t\to\infty$. For large $K$, the error is nonperturbative. 
	The way to address the difficulty is to subtract from $u$ a term responsible for the unboundedness of $u$ near $\partial B_{K\sqrt{t}}$, which is mainly the zero eigenmode of $L,$ and apply Proposition \ref{PropVelazquez} to the remaining part of the graphical function. We also observe that in terms of the pseudolocality theorem, while $|u|$ is not necessarily small, $|\nabla^k u|$ is small for $k=1,2,3$, and hence the error introduced by higher-order derivatives can be nicely controlled.
	
	\subsection{The $C^1$ normal form}
	In this section, we give the proof of Theorem \ref{ThmNF-sqrtt-C1}.	Using Proposition \ref{thm:mainthmfor sqrtt}, we have the equation of $u$ as in \eqref{EqRMCFFull}, inside the ball of radius $K\sqrt{t}$. Now we rewrite \eqref{EqRMCFFull} as
	$\pr_t u
	=
	\mathscr L( u)
	+\tilde m(u),$
	where we introduced the nonlinear operators $$\sL (u)= \frac{\varrho^2}{(\varrho+u)^2}\Delta_\theta u
	+\Delta_{\R^k} u
	-\frac12y_i\pr_i u
	+\frac{\varrho+u}{2}
	-\frac{\varrho^2}{2(\varrho+u)}.$$
	
	For the equation $\partial_t f=\mathscr L (f)$, we have an approximate solution. 
	\begin{lemma}\label{Lmf} Suppose $1\leq \ell\leq k$. Let	$
		f(\theta,y,t)
		=
		\varrho\sqrt{1+\frac{\sum_{i=1}^\ell (y_i^2-2)}{2t}}-\varrho.$ Then we have 	\begin{equation}\label{EqLf}
			\pr_t f - \sL (f)=
			-\frac{(n-k)^2(\sum_{i=1}^\ell y_i^2)}{t^2(\varrho^2+\frac{\varrho^2}{2t}\sum_{i=1}^\ell(y_i^2-2))^{3/2}}=:P(y,t).
		\end{equation}
		Moreover,
		\begin{equation}\label{eq:f close to H2}
			\left\|f(\theta,y,t)-\frac{\varrho}{4t}\sum_{i=1}^\ell (y_i^2-2)\right\|_{H^1(\cC_{n,k})}
			=
			O(t^{-2}).
		\end{equation}
	\end{lemma}
	
	\begin{proof}[Proof of Lemma \ref{Lmf}]
		It can be verified directly that   $g(\theta,\xi,t)=\varrho\sqrt{1+(\sum_{i=1}^\ell\xi_i^2)}$ solves the equation $-\frac{1}{2}\xi\cdot\nabla_\xi g+\frac{g^2-\varrho^2}{2g}=0$. Thus, $f$ solves  $\frac{\varrho^2}{(\varrho+f)^2}\Delta_\theta f-\frac12y_i\pr_i f
		+\frac{\varrho+f}{2}
		-\frac{\varrho^2}{2(\varrho+f)}=0$. 
		Then \eqref{EqLf} follows from the calculation
		$	\pr_{y_j} f=\frac{\varrho^2y_j\delta_{jl}}{2t\sqrt{\varrho^2+\frac{1}{2t}\varrho^2\sum_{i=1}^\ell (y_i^2-2)}}$ and 
		\[
		\pr^2_{y_jy_m} f
		=
		\frac{\varrho^4\delta_{jm}\delta_{jl}\delta_{ml}}{4t\sqrt{\varrho^2+\frac{\varrho^2}{2t}\sum_{i=1}^\ell (y_i^2-2)}}
		-
		\frac{\varrho^4y_jy_m\delta_{jl}\delta_{ml}}{(2t)^2(\varrho^2+\frac{\varrho^2}{2t}(\sum_{i=1}^\ell (y_i^2-2))^{3/2}}.
		\]
		This gives \eqref{EqLf}. 
		
		Finally, \eqref{eq:f close to H2} follows from the following direct computation
		\begin{equation}\label{eq:f-H2}
			f-\frac{\varrho}{4t}\sum_{i=1}^\ell (y_i^2-2)=
			\frac{\varrho}{8t^2}(\sum_{i=1}^\ell (y_i^2-2))^2
			\left(\sqrt{1+\frac{1}{2t}\sum_{i=1}^\ell (y_i^2-2)}+1\right)^{-2}.
		\end{equation}
	\end{proof}
	
	\begin{proof}[Proof of Theorem \ref{ThmNF-sqrtt-C1} for $C^1$ normal form]
		
		We can check that $w:=u-f$ satisfies the following equation
		\begin{equation}\label{eq:C1_diff_equation}
			\pr_t w
			=
			\sL u-\sL f
			+\tilde{m}(u)
			-P(y,t),
		\end{equation}
		which can be rewritten as	\begin{equation}\label{EqNoZeroModes}
			\pr_t w=L w-P+\mathscr Q(w),
		\end{equation}
		where we have \[
		\begin{split}
			\mathscr Q(w)=&\tilde m(f+w)+\left(\frac{\varrho^2\Delta_\theta w}{(\varrho+f+w)^2}-\Delta_\theta w\right)
			-\left(
			\frac{\varrho(2f+w)+(f+w)f}{2(\varrho+f+w)(\varrho+f)}
			\right)w
			\\&
			-\varrho^2\left(\frac{1}{(\varrho+f)^2}-\frac{1}{(\varrho+w+f)^2}
			\right)\Delta_\theta f.
		\end{split}
		\]
		
		We also notice that \begin{align*}
			\left(\frac{\varrho^2\Delta_\theta w}{(\varrho+f+w)^2}-\Delta_\theta w\right)&=\frac{\Delta_\theta w}{(\varrho+w+f)^2}f^2+\frac{\Delta_\theta w}{(\varrho+w+f)^2}(2fw+w^2),
			\\
			\left(\frac{\varrho(2f+w)+(f+w)f}{2(\varrho+f+w)(\varrho+f)}w\right)&=
			\frac{2w\varrho+w(w+f)}{2(\varrho+f+w)(\varrho+f)}f+\frac{\varrho}{2(\varrho+f+w)(\varrho+f)}w^2.
		\end{align*} Later we will handle terms multiplying with $f$ and terms multiplying with $w^2$ separately.
		
		Recall the cutoff function $\check\chi$ that is $1$ outside the ball $B_{t^\kappa}$ and vanishes inside $B_{t^\kappa-1}$, where $\kappa<1/2$ is in Proposition \ref{PropVelazquez}. Similar to Lemma \ref{LmVelaz}, we have the following differential inequality for $Z=\chi|u|+D(\chi |u|)$, where $\chi$ is a cutoff function that is 1 inside the ball $B_{K \sqrt{t}}$ and vanishes inside $B_{K\sqrt{t}+1}$
		
		\begin{equation}\label{EqZ}
			\pr_tZ-LZ
			\leq
			\eps_0Z+C(\frac{|y|^2+1}{t^2}+|y|\check\chi(y,t))    
		\end{equation}
		where the constant $C$ depends on $K$.
		The proof is based on the following observations: First of all, throughout, we are in the setting of the proof of Proposition \ref{thm:mainthmfor sqrtt}, especially Step 2 and Step 3. Given $\bar\eps>0$, we apply the pseudolocality theorem as in Step 2, see Theorem \ref{thm:K0sqrt t to K sqrt t}, we bound $\|D^ku\|_{C^0}\leq\bar\eps,\ k=1,2,3,$ and for $\|u\|_{C^0}$, we bound it by $CK^2$ over the ball $B_{K \sqrt{t}}$. We get the same bounds for $w$, up to a constant, by the expression of $f$.
		\begin{enumerate}
			\item In \eqref{EqNoZeroModes}, we bound $P$ and $\varrho^2\left(\frac{1}{(\varrho+f)^2}-\frac{1}{(\varrho+w+f)^2}
			\right)\Delta_\theta f$, $\frac{\Delta_\theta w}{(\varrho+w+f)^2}f^2$ and $\frac{2w\varrho+w(w+f)}{2(\varrho+f+w)(\varrho+f)}f$ in $\mathscr Q(w)$ by $\frac{|y|^2+1}{t^2}.$ 
			
			\item Similar to Lemma \ref{LmVelaz}, we bound $\mathscr{Q}$ and $D\mathscr Q$ by $C\bar\eps Z$. Indeed, in $\tilde m(u)=\tilde m(f+w)$, using the mean value theorem, the coefficients of $|w|$ and $|Dw|$ are bounded by the the combinations of $\|u\|_{C^3}$ and $\|f\|_{C^3}$. The remaining estimate follows from the same argument as Lemma \ref{LmVelaz}. Next, the term $\frac{\Delta_\theta w}{(\varrho+w+f)^2}(2fw+w^2)$ which is bounded by $C\bar\eps Z$ inside $B_{K t^\kappa}$, where we use the pseudolocality estimate to bound $|\Delta_\theta u|\leq\bar\eps.$ Finally, $-\frac{\varrho}{2(\varrho+f+w)(\varrho+f)}w^2\leq 0$, so it has the right sign and can be discarded from the inequality.
			
			
			\item The error created by the cutoff $\chi$ supported on $B_{K\sqrt{k}+1}\setminus B_{K\sqrt{t}}$ is bounded by $C|y|\check{\chi}$. 
		\end{enumerate}
		
		Finally, if we differentiate \eqref{EqNoZeroModes} by $\theta$ or $y_j$, the linearized part is already discussed in the proof of Lemma \ref{LmVelaz}. For the other parts, differentiating will introduce higher-order derivatives, and they are nicely controlled by the pseudolocality theorem. Then if we choose $\bar\eps$ sufficiently small, we get $Z=|w|+|\nabla w|$ satisfies the equation \eqref{EqZ}.
		
		
		We then apply  Proposition \ref{PropVelazquez} to  \eqref{EqZ}.  By the $H^1$-normal form and \eqref{eq:f close to H2}, we have $\|w\|_{H^1(B_{t^\kappa})}=O(1/t^2)$, which implies $\|\chi w\|_{H^1}=O(1/t^2)$ using the bound $|w|\leq CK^2$ and $|Dw|\leq\eps_0$ outside the ball $B_{t^\kappa}$ and the Gaussian weight.  Then Proposition \ref{PropVelazquez} implies that $\|w\|_{C^1(B_{K\sqrt{t}})}\leq O(1/t)$. 
		This completes the proof of Theorem \ref{ThmNF-sqrtt-C1} for the $C^1$-normal form. 
	\end{proof}
	
	
	\subsection{Extensions of the graphical scale to $K\sqrt{t}$ for all $K>0$}
	In this section, we show that once we obtain the $C^1$-normal form of the rescaled mean curvature flow inside the ball of radius $K_0\sqrt{t}$ for some $K_0>0$, we can extend the $C^1$-normal form to a larger scale, namely, there exists $\alpha>0$ that only depends on $n,k$, such that the $C^1$-normal form holds in the ball of radius $(1+\alpha)K_0\sqrt{t}$. Then an iteration argument shows that $C^1$-normal form holds inside the ball of radius $K\sqrt{t}$ for any $K>0$.
	
	
	
	The starting point is an extension of the graphical radius. Notice that in a nondegenerate direction, near the boundary of the ball of radius $K\sqrt{t}$, the $C^1$-normal form yields the graph function roughly $(\sqrt{1+K^2/2}-1)$, which is larger than any given $\eps$ if $K$ is sufficiently large. Therefore, we need to allow the graph function to have a possibly large $C^0$-norm. On the other hand, we still hope the higher-order derivatives are small, to ensure the nonlinear terms are small.
	
	We start with an application of the pseudolocality theorem.
	
	\begin{theorem}\label{thm:K0sqrt t to K sqrt t}
		For any $\bar \eps>0$ and $\vartheta\in(0,1)$, there exists $\alpha=\alpha(n,k,\bar \eps,\vartheta)>0$ with the following significance. Suppose for $t\geq T$, $M_t$ is a rescaled mean curvature flow that is the graph of a function $u(\cdot,t)$ over $ \cC_{n,k}\cap B_{K_0\sqrt{t}}$, and $\mathcal I\subset\{1,2,\cdots,k\}$, such that 
		\[
		\left\| u(\theta,y,t)-\left(\varrho\sqrt{1+\frac{\sum_{i\in\mathcal I}(y_i^2-2)}{2t}}-\varrho\right)
		\right\|_{C^1(\cC_{n,k}\cap B_{K_0\sqrt{t}})}=O(t^{-\vartheta}),\quad
		t\to\infty.
		\]
		Then there exists $\overline{T}>T$ such that $M_t$ is a graph of $u(\cdot,t)$ over $ \cC_{n,k}\cap B_{(1+\alpha)K_0\sqrt{t}}$ when $t>\bar T$, and 
		\begin{equation}\label{eq:Du and D2u bound}
			\sup_{ \cC_{n,k}\cap B_{(1+\alpha)K_0\sqrt{t}}}
			\sum_{j=1}^3|\nabla^j u(\cdot,t)|\leq \bar\eps.
		\end{equation}
	\end{theorem}
	
	\begin{proof}
		We consider three different cases.
		
		\noindent{\bf Case 1. The fully degenerate case} $\cI=\emptyset$. Let us fix $\bar\eps$ and choose $\eta_0$ in Theorem \ref{thm:pseudo-locality} so that Theorem \ref{thm:pseudo-locality} holds for $\eps=\bar\eps$. When $\bar t$ is sufficiently large such that $K_0\sqrt{\bar t}>R_0$ and $\|u(\cdot,\bar t)\|_{C^1(\cC_{n,k}\cap B_{K_0\sqrt{\bar t}})}<\eta$, the pseudolocality theorem shows that $M_t$ is a graph of $u(\cdot,t)$ over $\cC_{n,k}\cap B_{e^{t-\bar t}K_0\sqrt{\bar t}}$ for $t\in[\bar t,\bar t+\bar\delta]$ for some $\bar\delta>0$, and
		\[
		\sup_{\cC_{n,k}\cap B_{e^{\bar \delta}K_0\sqrt{\bar t}}}
		\sum_{j=1}^3|\nabla^j u(\cdot,\bar t+\bar\delta)|
		\leq \bar\eps.
		\]
		Notice that $e^{\bar\delta}\sqrt{\bar t}/\sqrt{\bar t+\bar \delta}\to e^{\bar \delta}>1$ as $\bar t\to\infty$, so when $\bar t$ is sufficiently large, we can choose $\alpha=e^{\bar \delta/2}-1$ and \eqref{eq:Du and D2u bound} holds.
		
		\noindent{\bf Case 2. The nondegenerate case} $\cI=\{1,2,\cdots,k\}$. Then by the $C^1$-normal form, for any $r>0$, and any $p\in\cC_{n,k}\cap (B_{K_0\sqrt{t}-r}\backslash B_{K_0\sqrt{t}-2r})$, the rescaled mean curvature flow is a graph of function $v(\cdot,t)$ over $(\sqrt{1+K_0^2/2}-1)\cC_{n,k}\cap B_r(p)$, and the $C^1$-norm of this function turns to $0$ as $t\to\infty$. In particular, when $\bar t$ is sufficiently large, we can apply the pseudolocality theorem to $M(\bar t)\cap B_r(p)$, to show that $v(\cdot,t)$ is a graph over $e^{t-\bar t}[(\sqrt{1+K_0^2/2}-1)\cC_{n,k}\cap B_r(p)]$ for $t\in[\bar t,\bar t+\bar \delta]$, and also notice that the gradient is invariant under dilation and the hessian is scaled down under dilation up, we have 
		\[
		\sup_{\cC_{n,k}\cap B_{(r-1)}}
		\sum_{j=1}^3|\nabla^j u(\cdot,\bar t+\bar \delta)|
		\leq \bar\eps.
		\]
		In particular, this shows that \eqref{eq:Du and D2u bound} holds for $\alpha=e^{\bar \delta/2}$.
		
		\noindent{\bf Case 3. Partially nondegenerate case} $\cI\subsetneq \{1,2,\cdots,k\}$. Without loss of generality, we assume $\cI=\{1,2,\cdots,m\}$. Now we divide the $\R^{n-k+1}\times\R^k$ into two parts $\R^k=\cA\cup\cB$, where
		\[
		\cA=\{(\theta,y):y_1^2+\cdots+y_m^2<\eta|y|^2\},
		\quad 
		\cB=\{(\theta,y):y_1^2+\cdots+y_m^2\geq \eta|y|^2\}.
		\]
		Then by the $C^1$-normal form, on $\cC_{n,k}\cap \cA\cap B_{\sqrt{t}}$, the $C^1$-norm of $u(\cdot,t)$ is bounded by $\sqrt{1+K^2\eta/2}-1+o(1)\leq K^2\eta+o(1)$; for any $r>0$, and any $p\in\cC_{n,k}\cap (B_{K_0\sqrt{t}-r}\backslash B_{K_0\sqrt{t}-2r})\cap \cB$, the rescaled mean curvature flow is a graph of function $v(\cdot,t)$ over $(\sqrt{1+K_0^2\varrho_p/2}-1)\cC_{n,k}\cap B_r(p)$, where $\varrho_p=\frac{p_1^2+\cdots+p_m^2}{|p|^2}\geq \eta$, and the $C^1$-norm of this function turns to $0$ as $t\to\infty$. Then by applying the argument in Case 1 to $\cC_{n,k}\cap \cA\cap B_{\sqrt{t}}$ and the argument in Case 2 to $\cC_{n,k}\cap \cB\cap B_{\sqrt{t}}$, we have \eqref{eq:Du and D2u bound} holds for some $\alpha>1$. This concludes the proof.
	\end{proof}
	
	
	One key ingredient in Vel\'azquez's regularization theorem is the weighted $L^2$ bound. The following lemma shows that the $L^2$ bound on small graphical radius can induce the $L^2$ bound on large graphical radius.
	
	\begin{lemma}\label{lem:K0sqrt t to K sqrt t L2 norm}
		Suppose $M_t$ is a rescaled mean curvature flow as in Theorem \ref{thm:K0sqrt t to K sqrt t}. Then 
		\[
		\left\| u(\theta,y,t)-
		\frac{\varrho}{4t}\sum_{i\in\cI }(y_i^2-2)
		\right\|_{H^1(\cC_{n,k}\cap B_{(1+\alpha)K_0\sqrt{t}})}=O(t^{-1-\vartheta}),\quad
		t\to\infty,
		\]
		where $\alpha$ is given in Theorem \ref{thm:K0sqrt t to K sqrt t}. 
	\end{lemma}
	\begin{proof}
		By the $C^1$-normal form, we already have the decay on $\cC_{n,k}\cap B_{K_0\sqrt{t}}$. So it suffices to show that the $H^1$-norm of the difference turns to $0$ on $\cA_t:=\cC_{n,k}\cap [B_{(1+\alpha)K_0\sqrt{t}}\backslash B_{K_0\sqrt{t}})]$.
		
		Note that on $\cA_t$, when $t>1$,
		\[
		\sum_{j=1}^3|\nabla^j u(\cdot,t)|
		\leq 
		(1+\alpha)^2 K_0^2/2,
		\]
		and by the proof of Theorem \ref{thm:K0sqrt t to K sqrt t}, 
		\[
		\sum_{j=1}^3|\nabla^j u(\cdot,t)|
		\leq
		C(K_0).
		\]
		Thus,
		\begin{align*}
			\left\| u(\theta,y,t)-
			\frac{\varrho}{4t}\sum_{i\in\cI }(y_i^2-2)
			\right\|^2_{H^1(\cA_t)}
			\leq& \int_{\cA_t} C(K_0)^2 e^{-\frac{|x|^2}{4}}d\cH^n(x)
			\\
			\leq &
			C(K_0)^2( (1+\alpha)K_0t^{1/2})^k e^{-K_0^2t/4}.
		\end{align*}
		Then by the exponential decay of this term, we have
		\[
		\left\| u(\theta,y,t)-
		\frac{\varrho}{4t}\sum_{i\in\cI }(y_i^2-2)
		\right\|_{H^1(\cC_{n,k}\cap B_{(1+\alpha)K_0\sqrt{t}})}=O(t^{-1-\vartheta}),\quad
		t\to\infty.
		\]
	\end{proof}
	
	
	\begin{proof}[Proof of Theorem \ref{ThmNF-sqrtt-C1} ($C^1$-normal form in $B_{K\sqrt{t}}$)]
		Suppose $\vartheta$ is given as in Section \ref{SNF}. The proof is the same as the proof of Proposition \ref{PropSmallSqrt}. The only difference is that in Step 2, we use Theorem \ref{thm:K0sqrt t to K sqrt t} to do the extension; and in Step 3, we use Lemma \ref{lem:K0sqrt t to K sqrt t L2 norm}, to get the desired $H^1$-norm decay as $O(t^{-1-\vartheta})$. All the rest of the proof is verbatim, which implies that we can extend the $C^1$-normal form from the radius $K_0\sqrt{t}$ to $(1+\alpha)K_0\sqrt{t}$. Then by iterating this process, we can get the $C^1$-normal form inside the ball of radius $K\sqrt{t}$ for any $K>0$.
	\end{proof}
	
	\begin{proof}[Proof of Theorem \ref{ThmNF-sqrtt}]
		Once we extended the graphical radius to $K\sqrt{t}$, the $H^1$-normal form inside $\cC_{n,k}\cap B_{K\sqrt{t}}$ is straightforward from the $H^1$-normal form proved in Section \ref{SNF} and repeatedly using Lemma \ref{lem:K0sqrt t to K sqrt t L2 norm}. The moreover part is proved in Proposition \ref{prop:decay order zero normal form}.
	\end{proof}
	
	
	
	\section{Geometric consequences of the $C^1$-normal form}\label{S_Geometric}
	In this section, we shall apply our $C^1$-normal form to obtain geometric consequences on the mean curvature flow. 
	
	\subsection{Isolatedness of nondegenerate singularities}\label{SIsolate}
	We first state a version of the pseudolocality. 
	
	\begin{lemma}[Pseudolocality of cylindrical MCF over $\cC_{n,k}$ with a slightly larger radius]\label{lem:plcyl}
		
		For any $\eps_0>0$, and a vector $V\in \R^k$ with $0<|V|<\eps_0$, there exist $T_0>0$, $R_1>0$ and $\eps_1>0$ with the following significance. Suppose $T>T_0$ and $\bM_\tau$ is a mean curvature flow, at time $\tau=0$ is the graph of a function $u$ over the cylinder $\cC_{n,k}$ inside a ball of radius $\eps_0\sqrt{T}$, with $\|u(\theta,y)-C|T^{-1/2}y-V|^2\|_{C^1}<\eps_1$, where $C>0$ is some constant. Then for $\tau\in[0,1]$, $\bM_\tau\cap B_{R_1}$ is a smooth mean curvature flow.
	\end{lemma}
	
	
	
	The proof of Lemma \ref{lem:plcyl} is a straightforward application of the pseudolocality. 
	\begin{proof}[Proof of Lemma \ref{lem:plcyl}]
		We use the pseudolocality property of mean curvature flow. We observe that the graph of the function $C|T^{-1/2}y-V|^2$ is a paraboloid over the cylinder when $|y|<\eps_0\sqrt{T}$. Moreover, for any fixed $R>0$, inside the ball of radius $R$, we have $C|T^{-1/2}y-V|^2\to C|V|^2$ as $T\to\infty$. So the graph of the function $C|T^{-1/2}y-V|^2$ converges to the cylinder $\mathbb S^{n-k}(\varrho+C|V|^2)\times\R^k$ smoothly compactly as $T\to\infty$. By the pseudolocality of mean curvature flow (See Theorem \ref{thm:pseudo-locality}), as $T\to\infty$, the mean curvature flow $\bM_\tau^T$ starting from the graph of function $C|T^{-1/2}y-V|^2$ converges to the shrinking cylinder mean curvature flow with initial radius $\varrho+C|V|^2$, which is smooth for time $\tau\in[0,1]$. This shows that when $T$ is sufficiently large, $\bM_\tau^T$ is sufficiently close to the shrinking cylinder mean curvature flow with initial radius $\varrho+C|V|^2$ with $\tau\in[0,1]$, inside the ball of radius $2R_1$. Finally, when $\eps_1$ is sufficiently small, a similar argument shows the lemma.	
	\end{proof}
	
	
	
	With this pseudolocality lemma, we can prove Theorem \ref{ThmIsolate}.
	
	\begin{proof}[Proof of Theorem \ref{ThmIsolate}] 
		Theorem \ref{ThmNF-sqrtt-C1} shows that the rescaled mean curvature flow is a graph over $\cC_{n,k}$ in a ball of radius $K_0\sqrt{t}$ when $t$ is sufficiently large. Let us cover the rescaled mean curvature flow $M_t\cap (B_{K_0\sqrt{t}}\backslash B_{2^{-1}K_0\sqrt{t}})$ with several balls $B_{R_1}(p_i)$ of radius $R_1:=2^{-1}K_0\sqrt{t}$ centered at $p_i\in B_{K_0\sqrt{t}}\backslash B_{2^{-1}K_0\sqrt{t}}$, $i=1,2,\ldots,m$, and Theorem \ref{ThmNF-sqrtt-C1} implies that we can apply Lemma \ref{lem:plcyl} to get that the mean curvature flow starting from $M_t$ is smooth in $B_{R_1}(p_i)$ for time $1$. Rescaling back, we see that $\bM_{\tau}$ is smooth in $B_{e^{-t/2}R_1}(e^{-t/2}p_i)$ for $\tau\in[-e^{-t},0]$. Because $t$ can be chosen arbitrarily large, $\bM_{\tau}\cap (B_{\delta}(0)\backslash\{0\})$ is smooth for $\tau\in[-1,0]$ for some $\delta>0$ (for example, we can choose it to be $e^{-T/2}K_0\sqrt{T}$ for some fixed sufficiently large $T$).
		
		Moreover, the region $\frac{K_0}{2}\sqrt t\leq |y|\leq K_0\sqrt t,\ t\in [T,\infty)$ for the rescaled mean curvature flow scale, corresponds to the region $\frac{K_0}{2} (-\log (-\tau))^{-1/2}\leq |y|\leq K_0 (-\log (-\tau))^{-1/2}, \ \tau \in [-e^{-T}, 0)$, which shrinks to the origin $0$ as $\tau\to 0$. Thus, the singularity is isolated in the backward $\delta$-spacetime neighborhood. 
	\end{proof}
	
	
	\subsection{Mean convex neighborhood and Type-I}\label{SS:Mean convex neighborhood and Type-I}
	Lemma \ref{lem:plcyl} shows that when $T$ is sufficiently large and $\eps_0$ is sufficiently small, the mean curvature flow is still a graph over the cylinder. Then the rescaled mean curvature flow is roughly the graph of the function $\sim \eps_0|V|^2$. This gives a description of the neighborhood of the non-degenerate singularity. 
	
	
	
	\begin{proof}[Proof of Theorem \ref{prop:Mconnbhd}]
		We only need to prove that for the corresponding rescaled mean curvature flow $\{M_t\}$, when $t$ is sufficiently large, $M_t\cap B_{\delta_2e^{t/2}}$ is mean convex.
		
		First, we show that when $t$ is sufficiently large, $M_t\cap B_{\delta_2\sqrt{t}}$ is mean convex, where $\delta_2$ is chosen slightly smaller than $K$ in Proposition \ref{thm:mainthmfor sqrtt}. In previous section, we have proved that $M_t\cap B_{\delta_2\sqrt{t}}$ can be written as the graph of a function $u(\cdot,t)$ over $\cC_{n,k}$, with $u(\theta,y,t)$ converges to $\varrho\sqrt{1+\frac{1}{2t}{\sum_{i=1}^k h_2(y_i)}}-\varrho$ in $C^1$-norm for $y$ in the ball of radius $\delta_2\sqrt{t}$. In particular, this shows that $u(\theta,y,t)$ converges to $\frac{\varrho|y|^2}{1+\sqrt{1+(2t)^{-1}|y|^2}}$ in $C^1$-norm in the ball of radius $\delta_2\sqrt{t}$. We also recall that Proposition \ref{thm:mainthmfor sqrtt} shows that $\|u\|_{C^2(\delta_2\sqrt{t})}<\eps$ for any given $\eps$, when $t$ is sufficiently large.
		
		The mean curvature of the graph of $u(\theta,y,t)=\frac{\varrho t^{-1}|y|^2}{1+\sqrt{1+t^{-1}|y|^2}}+w(\theta,y,t)$ is explicitly given by \eqref{eq:H}, and it is bounded from below by $\varrho/2-C\eps$, where $\varrho/2$ is the mean curvature of $\cC_{n,k}$, and $\eps$ is the $C^2$-bound of $u$. Thus, if initially, we fixed $\eps$ to be a sufficiently small number, we see the mean curvature is positive.
		
		Finally, from the scale $\delta_2\sqrt{t}$ to $\delta_2 e^{t/2}$, we use the same argument as the proof of Theorem \ref{ThmIsolate}. When $t$ is sufficiently large, we can decompose $M_t\cap B_{\delta_2\sqrt{t}}$ into the union of several pieces, each piece is very close to a cylinder with various radii. Under the evolution of mean curvature flow, each piece remains close to some cylinder. In particular, when $t$ is sufficiently large, each piece is mean convex. This concludes the proof.
		
	\end{proof}
	
	
	\begin{proof}[Proof of Theorem \ref{ThmtypeI}]
		We only need to prove that for the corresponding rescaled mean curvature flow $\{M_t\}$, when $t$ is sufficiently large, $M_t\cap B_{\delta_2e^{t/2}}$ has bounded $|A|$. This is true just as the proof of Theorem \ref{prop:Mconnbhd}.
	\end{proof}
	
	
	\subsection{Vanishing of zero modes}
	The result of this subsection is not used anywhere in this paper, but will be used in \cite{SunWangXue2_RegSing}. We have seen that the $H^1$-normal form of the rescaled mean curvature flow is given by $\sum_{i\in\cI}\frac{\varrho}{4t}(y_i^2-2)+O(1/t^2)$, where $\cI\subset\{1,2,\cdots,k\}$. When $\cI=\emptyset$, we expect that the next eigenmodes of $L_{\cC_{n,k}}$, that have strictly positive eigenvalues, should dominate the evolution. Then we should expect a faster decay of the graph function.
	
	\begin{proposition}[``Moreover'' part of Theorem \ref{ThmNF-sqrtt}]\label{prop:decay order zero normal form}
		Suppose $K>0$, $\br(t)=K\sqrt{t}$ and the rescaled mean curvature flow is a graph of the function $u(\cdot,t)$ over $\cC_{n,k}\cap B_{\br(t)}$, and $\cI=\emptyset$ in the $H^1$-normal form. Then  $\|u(\cdot,t)\|_{L^2(\cC_{n,k}\cap B_{\br(t)})}\leq Ce^{-K_0^2 t}$ for some $K_0>0$. 
	\end{proposition}
	
	\begin{proof}
		In  Proposition \ref{LmCoefficient}, we choose the graphical radius $\mathbf{r}(t)=t^\kappa,\ \kappa<1/2$, we have the dichotomy: either $\|u\|_{H^1}$ decays faster than $e^{-\frac{\mathbf{r}(t)^2}{4}}$ or it decays like $1/t$ with $\mathcal I\neq \emptyset.$ 
		
		We next show that the first case is indeed exponentially small in time. In the first case, we have $\mathcal I=\emptyset.$ Using the $C^1$ normal form over the ball $B_{Kt^{1/2}}$, we get that $\|u\|_{C^1(B_{Kt^{1/2}})}=o(1)$. This, combined with the first case in the dichotomy, we get that $\|u\|_{H^1(B_{Kt^{1/2}})}\leq e^{-\frac{\mathbf{r}(t)^2}{4}}$. Next, given $\eps>0$, by choosing $t$ large, we get $\|u(\cdot,t)\|_{C^1(B_{K_0\sqrt{t}})}\leq \eps/2$ by the $C^1$ normal form. We next use the pseudolocality Theorem \ref{thm:K0sqrt t to K sqrt t} to bound $\|\nabla ^2u(\cdot,t)\|_{C^0(K\sqrt{t})}\leq \eps/2$, thus, we get the $C^2$ estimate $\|u(\cdot,t)\|_{C^2(B_{K\sqrt{t}})}\leq \eps$.
		
		We choose $K_0<K$ small such that there is no eigenvalue in the interval $(-K_0^2/4,0)$. Applying this $C^2$ bound to equation \eqref{Eq3Eqs} in Lemma \ref{LmContinuation}, we get that $\eps_1$ -- the bound of $\|u\|_{C^2(B_{t^\kappa})}$ in \eqref{Eq3Eqs} -- is a bounded small number and $\omega(t)=e^{-K_0^2t/4}$ over the ball $B_{K_0\sqrt{t}}$. Then the argument in the proof of Lemma \ref{LmContinuation} gives that either $\|u\|_{H^1(B(K_0\sqrt{t}))}\leq e^{-K_0^2t/4}$, or modes in $E^0$ dominates. Thus, when $\mathcal I=\emptyset,$ the only possibility is that $\|u\|_{H^1(B_{K_0\sqrt{t}})}$ decays exponentially fast. Finally, using the bound $\|u\|_{C^1(B_{K\sqrt{t}})}=o(1)$ for the region $B_{K\sqrt{t}}\setminus B_{K_0\sqrt{t}}$, we get $\|u\|_{H^1(B_{K\sqrt{t}})}\leq C e^{-K_0^2t/4}$ for any $K>0$.
		
	\end{proof}
	
	\begin{remark}
		If $K_0$ is sufficiently large, $-K_0^2/4$ will be smaller than some eigenvalue of $L$. Then the exponential decay is given by $\|u\|_{H^1(B(K_0\sqrt{t}))}\leq Ce^{(\lambda+\eps) t}$ where $\lambda$ is the largest negative eigenvalue of $L$ and $\eps>0$ is a small number. This is almost sharp, see \cite{SunWangXue2_RegSing}.
	\end{remark}
	
	
	
	\appendix
	
	\section{Estimates of the nonlinear term}\label{appendix:Nonlinear estimate}	
	In this appendix, we derive the equation of motion for the graphical function $u$ of a manifold evolving under rescaled mean curvature flow approaching a cylinder $\cC_{n,k}=\bS^{n-k}(\varrho)\times \R^{k}$ where $\varrho=\sqrt{2(n-k)}$ is the radius of the sphere. We will use coordinates $z=(\theta, y)\in \cC_{n,k}$ where $\theta\in \R^{n-k+1}$ denotes point on $\bS^{n-k}(\varrho)$ and $y$ denotes point on $\R^{k}$. Note $|\theta|=\varrho$. The following computations are locally around a point $z_0=(\theta_0,y_0)$, and we choose local orthonormal frame $\{\theta_\alpha\}$ and $\{y_i\}$. Greek letters correspond to the spherical part and $i,j,k$'s correspond to the $\R^{k}$ part. We will use $\pr_\alpha$ and $\pr_i$ to simplify $\pr_{\theta_\alpha}$ and $\pr_{y_i}$ respectively. 
	\begin{proposition}\label{PropQ}Let $M_t$ be a rescaled mean curvature flow converging to a cylinder $\cC_{n,k}$ in the $C^\infty_{loc}$ sense as $t\to\infty$. Writing $M_t$ as a normal graph of a function $u$ over $\cC_{n,k}$ within the graphical radius, then we have  
		\begin{enumerate}
			\item 	we have
			$\partial_tu=L u+\cQ(J^2u),\ J^2u:=(u,\nabla u, \nabla^2u),$
			where we have 	\begin{equation}\label{EqQ<}
				|\cQ(J^2u)|\leq 
				C(|\nabla u|^4+|\nabla u|^2|\Hess_u|+|\nabla u|^2+u^2+|u||\Hess_u|)
			\end{equation}
			if $\|u\|_{C^2}\leq \eps_0$ for some $\eps_0$ small. 
			\item 	The leading terms in $\cQ$ is given explicitly as 
			\begin{equation}\label{EqQ}
				\cQ(J^2u)=-(2\varrho)^{-1}\left(u^2+4u\Delta_\theta u+2|\nabla_\theta u|^2\right) +\cC(J^2u),	
			\end{equation}
			where $\cC(J^2u)$ consists of terms cubic and higher power in $J^2u$, satisfying
			\begin{equation}
				|\cC(J^2u)|\leq C(|u|^3+|\nabla u|^3+|\nabla^2u|(u^2+|\nabla u|^2))\leq C\eps_0(u^2+|\nabla u|^2).
			\end{equation} 
			\item For three functions $v, u_1,u_2:\ \Sigma\cap B_R\to \R$ and the cutoff function $\chi$ that is 1 on $B_{R-1}$ and $0$ outside $B_R$, we have $($denoting $w=u_1-u_2)$\begin{equation*}\label{EqdC}
				\int_{\Sigma\cap B_R} |v 	(\cC(J^2(\chi u_1))-	\cC(J^2(\chi u_2)))|e^{-\frac{|x|^2}{4}}\leq C\max\{\|u_1\|_{C^2},\|u_2\|_{C^2}\}^2(\|w\|\cdot \|v\|+e^{-\frac{R^2}{4}}).
			\end{equation*}
		\end{enumerate}
		
		
		
	\end{proposition}
	
	\begin{proof}
		
		Note that the unit normal vector $\bn=\frac{\theta}{\varrho}$. Later we will use the following fact:
		\[\pr_i \theta=\pr_i\pr_j\theta=0,\ \ \ 
		\pr_\alpha\theta=\theta_\alpha,\ \ \ \pr_{\alpha}\pr_{\beta} \theta = -\frac{\delta_{\alpha\beta}}{\varrho^2}\theta.
		\]
		We consider a graph $\cC_{n,k}$ locally given by $\{\widetilde{F}(z)=(z)+u(z)\bn= (z)+u(z)\frac{\theta}{\varrho}\}$, which induces a frame on $\cC_{n,k}$ given by
		$
		\widetilde{\pr}_\alpha=(1+\frac{u}{\varrho})\theta_{\alpha}+\frac{\pr_\alpha u}{\varrho} \cdot\theta,\ \ \ 
		\widetilde{\pr}_{i}=\pr_i +\frac{\pr_iu}{\varrho} \cdot \theta.
		$
		The induced metric is given by
		\begin{equation*}
			\widetilde{g}_{\alpha\beta}=(1+\frac{u}{\varrho})^2 g_{\alpha\beta} + \pr_\alpha u \pr_\beta u
			,
			\ \ \ 
			\widetilde{g}_{i j}
			=
			g_{ij}+\pr_i u\pr_j u
			,
			\ \ \ 
			\widetilde{g}_{\alpha i}
			=
			\pr_{\alpha} u \pr_i u.
		\end{equation*}
		The inverse matrix is given by
		\begin{equation*}
			\begin{aligned}
				\widetilde{g}^{\alpha\beta}
				&=
				(1+\frac{u}{\varrho})^{-2}g_{\alpha\beta} - (1+\frac{u}{\varrho})^{-4}\pr_\alpha u \pr_\beta u+ m_{\alpha\beta},
				\quad	\\
				\widetilde{g}^{i j}
				&=
				g_{ij} - \pr_i u\pr_j u 
				+ m_{ij},\quad 
				\widetilde{g}^{\alpha i}
				=
				-\pr_\alpha u\pr_i u
				+ m_{\alpha i}.
			\end{aligned}
		\end{equation*}
		Here $m$ are functions of products of $\nabla u$, at least quadratically.	We find a unit normal vector field given by
		\begin{equation*}
			\widetilde{\bn}
			=
			\frac{\frac{\theta}{\varrho}-(1+\frac{u}{\varrho})^{-1}\pr_\alpha u\theta_{\alpha} -\pr_i u \pr_i}
			{\sqrt{
					1+\sum_{\alpha}(1+\frac{u}{\varrho})^{-2}(\pr_\alpha u)^2	
					+\sum_i (\pr_i u)^2.
			}}
			=:
			\frac{\frac{\theta}{\varrho}-(1+\frac{u}{\varrho})^{-1}\pr_\alpha u\theta_{\alpha} -\pr_i u \pr_i}{S}.
		\end{equation*}
		
		Now we calculate the 2nd fundamental form. We have
		\begin{equation*}
			\begin{aligned}
				\pr_{\alpha}\pr_{\beta} F
				&=
				-(1+\frac{u}{\varrho})\frac{\delta_{\alpha\beta}}{\varrho^2}\theta
				+
				\frac{\pr_{\beta} u}{\varrho}\theta_{\alpha}
				+
				\frac{\pr_{\alpha }u}{\varrho}\theta_\beta
				+
				\frac{\pr_{\alpha\beta}u}{\varrho}\theta,
				\\
				\pr_i\pr_{\alpha} F
				&	=
				\frac{\pr_{i\alpha} u}{\varrho}\theta
				+\frac{\pr_i u}{\varrho}\theta_{\alpha},
				\quad		\pr_i\pr_j F
				=
				\frac{\pr_{ij}u}{\varrho}\theta.
			\end{aligned}
		\end{equation*}
		
		Taking inner product with $\widetilde{\bn}$, we get the 2nd fundamental forms:
		\begin{equation*}
			\begin{aligned}
				\tilde{A}_{\alpha\beta}
				&=
				S^{-1}\left(
				-(1+\frac{u}{\varrho})\frac{\delta_{\alpha\beta}}{\varrho}
				-2(1+\frac{u}{\varrho})^{-1}\frac{\pr_\alpha u \pr_{\beta}u}{\varrho}
				+\pr_{\alpha\beta} u
				\right)	,\\
				\tilde{A}_{i\alpha}
				&=
				S^{-1}\left(
				\pr_{i\alpha} u 
				-
				(1+\frac{u}{\varrho})^{-1}\frac{\pr_{\alpha}u \pr_i u}{\varrho}
				\right)	,
				\quad 		\widetilde{A}_{ij}
				=
				S^{-1}\left(\pr_{ij}u
				\right)	.
			\end{aligned}	
		\end{equation*}
		
		In conclusion, we have (note that the convention in mean curvature flow is $H=-\tr A$)
		\begin{equation}\label{eq:H}
			\begin{split}
				-\widetilde{H}
				=&
				S^{-1}
				\left(
				(1+\frac{u}{\varrho})^{-2}\Delta_\theta u
				+\Delta_{\R^k} u
				-\pr_{ij}u\pr_i u\pr_j u
				-\pr_{i\alpha} u\pr_i u\pr_\alpha u
				\right.
				\\
				&
				\left.-(1+\frac{u}{\varrho})^{-4}\pr_\alpha u \pr_\beta u\pr_{\alpha\beta}u-(1+\frac{u}{\varrho})^{-1}\frac{k}{\varrho}
				-(1+\frac{u}{\varrho})^{-3}\frac{|\nabla_\theta u|^2}{\varrho}
				+m
				\right)
			\end{split}
		\end{equation}
		where $m$ is a remainder function of the linear combinations of products of $\nabla u$, with order at least $4$. In fact, we have
		\begin{equation}\label{Eqm1}
			|m|\leq C|\nabla u|^4,\quad  |\nabla m|\leq C|\nabla u|^3.
		\end{equation}
		
		Next we consider the term $\frac{1}{2}\langle \widetilde z,\widetilde \bn\rangle$. On $\cC_{n,k}$ near $z_0=(\theta_0,y_0)$, we have
		\begin{equation*}
			\begin{split}
				\langle \widetilde z,\widetilde\bn\rangle
				=&
				S^{-1}\left\langle (z)+\frac{u}{\varrho}\theta
				, \frac{\theta}{\varrho}-(1+\frac{u}{\varrho})^{-1}\pr_\alpha u\theta_{\alpha} -\pr_i u \pr_i
				\right\rangle
				=
				S^{-1}
				\left(
				\varrho+u
				-y_i\pr_i u
				\right)
			\end{split}
		\end{equation*}
		Therefore, we can obtain the equation of $u$ from the rescaled mean curvature flow equation (modulo diffeomorphism) $(\pr_t \widetilde z)^{\bot}=-(\widetilde H-\frac{\langle \widetilde z,\widetilde\bn\rangle}{2})\widetilde \bn$. Take an inner product of the equation with $\widetilde \bn$, we get 
		\begin{equation*}
			\begin{split}
				S^{-1} \pr_t u
				=&
				S^{-1}
				\left(
				(1+\frac{u}{\varrho})^{-2}\Delta_\theta u
				+\Delta_{\R^k} u
				-\pr_{ij}u\pr_i u\pr_j u
				-\pr_{i\alpha} u\pr_i u\pr_\alpha u
				-(1+\frac{u}{\varrho})^{-4}\pr_\alpha u \pr_\beta u\pr_{\alpha\beta}u
				\right.
				\\
				&
				\left.-(1+\frac{u}{\varrho})^{-1}\frac{k}{\varrho}
				-(1+\frac{u}{\varrho})^{-3}\frac{|\nabla_\theta u|^2}{\varrho}
				+m
				\right)+S^{-1}\frac{1}{2}(\varrho+u-y_i\pr_i u).
			\end{split}
		\end{equation*}
		Thus we obtain the equation of $u$ as follows:
		\begin{equation}\label{EqRMCFFull}
			\begin{split}
				&\pr_t u
				=
				\left(
				(1+\frac{u}{\varrho})^{-2}\Delta_\theta u
				+\Delta_{\R^k} u
				-\pr_{ij}u\pr_i u\pr_j u
				-\pr_{i\alpha} u\pr_i u\pr_\alpha u	-(1+\frac{u}{\varrho})^{-1}\frac{n-k}{\varrho}\right.\\
				&\left.-(1+\frac{u}{\varrho})^{-4}\pr_\alpha u \pr_\beta u\pr_{\alpha\beta}u
				-(1+\frac{u}{\varrho})^{-3}\frac{|\nabla_\theta u|^2}{\varrho}
				+m
				\right)
				+\frac{1}{2}(\varrho+u-y_i\pr_i u).
			\end{split}
		\end{equation}
		When $|u|_{C^0}\leq \eps_0$, we can further rewrite the equation as 
		$\pr_t u
		=
		Lu+\cQ(J^2u),$
		where $\cQ(J^2u)=m_1+m_2+m_3+m_4$, with $m_4$ consisting of all terms like $\pr u\pr u\pr^2 u$ and 
		\begin{equation*}
			m_1=m,
			\ \ \ m_2=-\frac{1}{2\varrho}\frac{u^2}{1+\frac{u}{\varrho}}-(1+\frac{u}{\varrho})^{-3}\frac{|\nabla_\theta u|^2}{\varrho},
			\ \ \ 
			m_3=-\frac{\frac{2u}{\varrho}+\frac{u^2}{\varrho^2}}{(1+\frac{u}{\varrho})^2}\Delta_\theta u.
		\end{equation*}
		From \eqref{EqRMCFFull}, we get \eqref{EqQ} immediately. Indeed, the leading terms come from $m_2$ and $m_3$ respectively. The terms $m_1$ (see \eqref{Eqm1}) and $m_4$ contribute only to $\cC(J^2u)$. Estimate \eqref{EqQ<} follows also immediately from \eqref{EqRMCFFull}. In particular, when $\|u\|_{C^{2}}\leq \eps_0$, we have the estimate
		$		|\cQ|\leq C(|\nabla u|^2+\eps_0u).
		$	
		For two different functions $u_1$ and $u_2$, the fundamental theorem of calculus shows that
		\begin{equation*}
			|\cQ(J^2u_1)-\cQ(J^2u_2)|\leq C\eps_0(|u_1-u_2|+|\nabla u_1-\nabla u_2|+|\Hess_{u_1-u_2}|).
		\end{equation*}
		We next consider item (3). To get the terms $\|u\|\|v\|$ on the RHS, we need to perform a step of integration by parts for terms of the form $v\mathrm{Hess}_u$, which also gives us the boundary term $e^{-R^2/4}$. Taking terms $\partial_iu\partial_j u\partial_{ij}u$ in $m_4$ as an example, we have $$\partial_iu\partial_j u\partial_{ij}u=(\nabla u)^T\nabla^2u\nabla u=\frac{1}{2}\nabla\cdot (|\nabla u|^2\nabla u)-\frac{1}{2}\Delta u \cdot |\nabla u|^2. $$ 
		Consider the first term on the RHS
		\begin{equation*}
			\begin{aligned}
				&\nabla\cdot (|\nabla (\chi u_1)|^2\nabla (\chi u_1))-\nabla\cdot (|\nabla (\chi u_2)|^2\nabla (\chi u_2))\\
				&=\nabla\cdot ((\nabla (\chi w)\cdot\nabla(\chi u_1+\chi u_2)) \nabla (\chi u_1))+\nabla\cdot (|\nabla (\chi u_2)|^2\nabla (\chi w))). 		
			\end{aligned}
		\end{equation*}
		When multiplied by $v$ and taking integration by parts, we see that it can be bounded by the RHS of \eqref{EqdC}. 
		When the derivative during the integration by parts hits the Gaussian weight $e^{-\frac{|x|^2}{4}}$, we shall get $\int v|x|e^{-\frac{|x|^2}{4}}$ which is bounded by $\|v\|_{H^1}$ using the Ecker's inequality $\int v^2|x|^2e^{-\frac{|x|^2}{4}}\leq C\int (v^2+|\nabla v|^2) e^{-\frac{|x|^2}{4}}. $ 
		All other terms can be treated similarly and are easier, so we get \eqref{EqdC}. 
	\end{proof}
	\bibliographystyle{alpha}
	\bibliography{GMT}
\end{document} 

